\documentclass[
  doc,
  floatsintext,
  longtable,
  a4paper,
  nolmodern,
  notxfonts,
  notimes,
  12pt,
  colorlinks=true,linkcolor=blue,citecolor=blue,urlcolor=blue]{apa7}

\usepackage{amsmath}
\usepackage{amssymb}

\geometry{inner=1in, outer=1in}
\fancyhfoffset[LE,RO]{0cm}


\usepackage[bidi=default]{babel}
\babelprovide[main,import]{spanish}


% get rid of language-specific shorthands (see #6817):
\let\LanguageShortHands\languageshorthands
\def\languageshorthands#1{}

\RequirePackage{longtable}
\RequirePackage{threeparttablex}

\makeatletter
\renewcommand{\paragraph}{\@startsection{paragraph}{4}{\parindent}%
	{0\baselineskip \@plus 0.2ex \@minus 0.2ex}%
	{-.5em}%
	{\normalfont\normalsize\bfseries\typesectitle}}

\renewcommand{\subparagraph}[1]{\@startsection{subparagraph}{5}{0.5em}%
	{0\baselineskip \@plus 0.2ex \@minus 0.2ex}%
	{-\z@\relax}%
	{\normalfont\normalsize\bfseries\itshape\hspace{\parindent}{#1}\textit{\addperi}}{\relax}}
\makeatother




\usepackage{longtable, booktabs, multirow, multicol, colortbl, hhline, caption, array, float, xpatch}
\usepackage{subcaption}


\renewcommand\thesubfigure{\Alph{subfigure}}
\setcounter{topnumber}{2}
\setcounter{bottomnumber}{2}
\setcounter{totalnumber}{4}
\renewcommand{\topfraction}{0.85}
\renewcommand{\bottomfraction}{0.85}
\renewcommand{\textfraction}{0.15}
\renewcommand{\floatpagefraction}{0.7}

\usepackage{tcolorbox}
\tcbuselibrary{listings,theorems, breakable, skins}
\usepackage{fontawesome5}

\definecolor{quarto-callout-color}{HTML}{909090}
\definecolor{quarto-callout-note-color}{HTML}{0758E5}
\definecolor{quarto-callout-important-color}{HTML}{CC1914}
\definecolor{quarto-callout-warning-color}{HTML}{EB9113}
\definecolor{quarto-callout-tip-color}{HTML}{00A047}
\definecolor{quarto-callout-caution-color}{HTML}{FC5300}
\definecolor{quarto-callout-color-frame}{HTML}{ACACAC}
\definecolor{quarto-callout-note-color-frame}{HTML}{4582EC}
\definecolor{quarto-callout-important-color-frame}{HTML}{D9534F}
\definecolor{quarto-callout-warning-color-frame}{HTML}{F0AD4E}
\definecolor{quarto-callout-tip-color-frame}{HTML}{02B875}
\definecolor{quarto-callout-caution-color-frame}{HTML}{FD7E14}

%\newlength\Oldarrayrulewidth
%\newlength\Oldtabcolsep


\usepackage{hyperref}



\usepackage{color}
\usepackage{fancyvrb}
\newcommand{\VerbBar}{|}
\newcommand{\VERB}{\Verb[commandchars=\\\{\}]}
\DefineVerbatimEnvironment{Highlighting}{Verbatim}{commandchars=\\\{\}}
% Add ',fontsize=\small' for more characters per line
\usepackage{framed}
\definecolor{shadecolor}{RGB}{241,243,245}
\newenvironment{Shaded}{\begin{snugshade}}{\end{snugshade}}
\newcommand{\AlertTok}[1]{\textcolor[rgb]{0.68,0.00,0.00}{#1}}
\newcommand{\AnnotationTok}[1]{\textcolor[rgb]{0.37,0.37,0.37}{#1}}
\newcommand{\AttributeTok}[1]{\textcolor[rgb]{0.40,0.45,0.13}{#1}}
\newcommand{\BaseNTok}[1]{\textcolor[rgb]{0.68,0.00,0.00}{#1}}
\newcommand{\BuiltInTok}[1]{\textcolor[rgb]{0.00,0.23,0.31}{#1}}
\newcommand{\CharTok}[1]{\textcolor[rgb]{0.13,0.47,0.30}{#1}}
\newcommand{\CommentTok}[1]{\textcolor[rgb]{0.37,0.37,0.37}{#1}}
\newcommand{\CommentVarTok}[1]{\textcolor[rgb]{0.37,0.37,0.37}{\textit{#1}}}
\newcommand{\ConstantTok}[1]{\textcolor[rgb]{0.56,0.35,0.01}{#1}}
\newcommand{\ControlFlowTok}[1]{\textcolor[rgb]{0.00,0.23,0.31}{\textbf{#1}}}
\newcommand{\DataTypeTok}[1]{\textcolor[rgb]{0.68,0.00,0.00}{#1}}
\newcommand{\DecValTok}[1]{\textcolor[rgb]{0.68,0.00,0.00}{#1}}
\newcommand{\DocumentationTok}[1]{\textcolor[rgb]{0.37,0.37,0.37}{\textit{#1}}}
\newcommand{\ErrorTok}[1]{\textcolor[rgb]{0.68,0.00,0.00}{#1}}
\newcommand{\ExtensionTok}[1]{\textcolor[rgb]{0.00,0.23,0.31}{#1}}
\newcommand{\FloatTok}[1]{\textcolor[rgb]{0.68,0.00,0.00}{#1}}
\newcommand{\FunctionTok}[1]{\textcolor[rgb]{0.28,0.35,0.67}{#1}}
\newcommand{\ImportTok}[1]{\textcolor[rgb]{0.00,0.46,0.62}{#1}}
\newcommand{\InformationTok}[1]{\textcolor[rgb]{0.37,0.37,0.37}{#1}}
\newcommand{\KeywordTok}[1]{\textcolor[rgb]{0.00,0.23,0.31}{\textbf{#1}}}
\newcommand{\NormalTok}[1]{\textcolor[rgb]{0.00,0.23,0.31}{#1}}
\newcommand{\OperatorTok}[1]{\textcolor[rgb]{0.37,0.37,0.37}{#1}}
\newcommand{\OtherTok}[1]{\textcolor[rgb]{0.00,0.23,0.31}{#1}}
\newcommand{\PreprocessorTok}[1]{\textcolor[rgb]{0.68,0.00,0.00}{#1}}
\newcommand{\RegionMarkerTok}[1]{\textcolor[rgb]{0.00,0.23,0.31}{#1}}
\newcommand{\SpecialCharTok}[1]{\textcolor[rgb]{0.37,0.37,0.37}{#1}}
\newcommand{\SpecialStringTok}[1]{\textcolor[rgb]{0.13,0.47,0.30}{#1}}
\newcommand{\StringTok}[1]{\textcolor[rgb]{0.13,0.47,0.30}{#1}}
\newcommand{\VariableTok}[1]{\textcolor[rgb]{0.07,0.07,0.07}{#1}}
\newcommand{\VerbatimStringTok}[1]{\textcolor[rgb]{0.13,0.47,0.30}{#1}}
\newcommand{\WarningTok}[1]{\textcolor[rgb]{0.37,0.37,0.37}{\textit{#1}}}

\providecommand{\tightlist}{%
  \setlength{\itemsep}{0pt}\setlength{\parskip}{0pt}}
\usepackage{longtable,booktabs,array}
\usepackage{calc} % for calculating minipage widths
% Correct order of tables after \paragraph or \subparagraph
\usepackage{etoolbox}
\makeatletter
\patchcmd\longtable{\par}{\if@noskipsec\mbox{}\fi\par}{}{}
\makeatother
% Allow footnotes in longtable head/foot
\IfFileExists{footnotehyper.sty}{\usepackage{footnotehyper}}{\usepackage{footnote}}
\makesavenoteenv{longtable}

\usepackage{graphicx}
\makeatletter
\newsavebox\pandoc@box
\newcommand*\pandocbounded[1]{% scales image to fit in text height/width
  \sbox\pandoc@box{#1}%
  \Gscale@div\@tempa{\textheight}{\dimexpr\ht\pandoc@box+\dp\pandoc@box\relax}%
  \Gscale@div\@tempb{\linewidth}{\wd\pandoc@box}%
  \ifdim\@tempb\p@<\@tempa\p@\let\@tempa\@tempb\fi% select the smaller of both
  \ifdim\@tempa\p@<\p@\scalebox{\@tempa}{\usebox\pandoc@box}%
  \else\usebox{\pandoc@box}%
  \fi%
}
% Set default figure placement to htbp
\def\fps@figure{htbp}
\makeatother







\usepackage{newtx}

\defaultfontfeatures{Scale=MatchLowercase}
\defaultfontfeatures[\rmfamily]{Ligatures=TeX,Scale=1}





\title{Validación cruzada y pipelines Python}


\shorttitle{VALIDACIÓN Y MODELO}


\usepackage{etoolbox}



\ccoppy{\textcopyright~2022}



\author{Edison Achalma}



\affiliation{
{Escuela Profesional de Economía, Universidad Nacional de San Cristóbal
de Huamanga}}




\leftheader{Achalma}

\date{2022-10-31}


\abstract{Este abstract será actualizado una vez que se complete el
contenido final del artículo. }

\keywords{keyword1, keyword2}

\authornote{\par{\addORCIDlink{Edison Achalma}{0000-0001-6996-3364}} 
\par{ }
\par{   El autor no tiene conflictos de interés que revelar.    Los
roles de autor se clasificaron utilizando la taxonomía de roles de
colaborador (CRediT; https://credit.niso.org/) de la siguiente
manera:  Edison Achalma:   conceptualización, metodología, análisis
formal, investigación, recursos, curación de
datos, redacción, visualización, supervisión, administración del
proyecto}
\par{La correspondencia relativa a este artículo debe dirigirse a Edison
Achalma, Escuela Profesional de Economía, Universidad Nacional de San
Cristóbal de Huamanga, Portal Independencia N°
57, Ayacucho, AYA 5001, Perú, Email: \href{mailto:elmer.achalma.09@unsch.edu.pe}{elmer.achalma.09@unsch.edu.pe}}
}

\makeatletter
\let\endoldlt\endlongtable
\def\endlongtable{
\hline
\endoldlt
}
\makeatother

\urlstyle{same}



\makeatletter
\@ifpackageloaded{caption}{}{\usepackage{caption}}
\AtBeginDocument{%
\ifdefined\contentsname
  \renewcommand*\contentsname{Tabla de contenidos}
\else
  \newcommand\contentsname{Tabla de contenidos}
\fi
\ifdefined\listfigurename
  \renewcommand*\listfigurename{Lista de Figuras}
\else
  \newcommand\listfigurename{Lista de Figuras}
\fi
\ifdefined\listtablename
  \renewcommand*\listtablename{Lista de Tablas}
\else
  \newcommand\listtablename{Lista de Tablas}
\fi
\ifdefined\figurename
  \renewcommand*\figurename{Figura}
\else
  \newcommand\figurename{Figura}
\fi
\ifdefined\tablename
  \renewcommand*\tablename{Tabla}
\else
  \newcommand\tablename{Tabla}
\fi
}
\@ifpackageloaded{float}{}{\usepackage{float}}
\floatstyle{ruled}
\@ifundefined{c@chapter}{\newfloat{codelisting}{h}{lop}}{\newfloat{codelisting}{h}{lop}[chapter]}
\floatname{codelisting}{Listado}
\newcommand*\listoflistings{\listof{codelisting}{Lista de Listados}}
\makeatother
\makeatletter
\makeatother
\makeatletter
\@ifpackageloaded{caption}{}{\usepackage{caption}}
\@ifpackageloaded{subcaption}{}{\usepackage{subcaption}}
\makeatother
\makeatletter
\@ifpackageloaded{fontawesome5}{}{\usepackage{fontawesome5}}
\makeatother

% From https://tex.stackexchange.com/a/645996/211326
%%% apa7 doesn't want to add appendix section titles in the toc
%%% let's make it do it
\makeatletter
\xpatchcmd{\appendix}
  {\par}
  {\addcontentsline{toc}{section}{\@currentlabelname}\par}
  {}{}
\makeatother

%% Disable longtable counter
%% https://tex.stackexchange.com/a/248395/211326

\usepackage{etoolbox}

\makeatletter
\patchcmd{\LT@caption}
  {\bgroup}
  {\bgroup\global\LTpatch@captiontrue}
  {}{}
\patchcmd{\longtable}
  {\par}
  {\par\global\LTpatch@captionfalse}
  {}{}
\apptocmd{\endlongtable}
  {\ifLTpatch@caption\else\addtocounter{table}{-1}\fi}
  {}{}
\newif\ifLTpatch@caption
\makeatother

\begin{document}

\maketitle


\hypertarget{toc}{}
\tableofcontents
\newpage
\section[Introduction]{Validación cruzada y pipelines Python}

\setcounter{secnumdepth}{3}

\setlength\LTleft{0pt}




En esta undécima y última guía exploraremos técnicas esenciales para
evaluar y optimizar modelos de machine learning: validación cruzada,
búsqueda de hiperparámetros, y análisis de contribución de datos y
variables. Estas técnicas son fundamentales para desarrollar modelos
robustos que generalicen bien a datos no vistos.

\section{El problema del overfitting}\label{el-problema-del-overfitting}

\subsection{Fundamento teórico}\label{fundamento-teuxf3rico}

El \textbf{overfitting} (sobreajuste) ocurre cuando un modelo captura no
solo los patrones verdaderos en los datos, sino también el ruido
aleatorio. Este es uno de los problemas fundamentales en machine
learning y econometría.

\subsubsection{Descomposición del error de
predicción}\label{descomposiciuxf3n-del-error-de-predicciuxf3n}

Para cualquier modelo de predicción, el error esperado fuera de muestra
se descompone en tres componentes:

\[
\mathbb{E}[(Y - \hat{f}(\mathbf{X}))^2] = \text{Bias}^2[\hat{f}(\mathbf{X})] + \text{Var}[\hat{f}(\mathbf{X})] + \sigma^2
\]

donde:

\textbf{1. Sesgo (Bias)}: \[
\text{Bias}[\hat{f}(\mathbf{X})] = \mathbb{E}[\hat{f}(\mathbf{X})] - f(\mathbf{X})
\]

Mide qué tan lejos está la predicción promedio del valor verdadero.
Modelos simples tienen alto sesgo (underfitting).

\textbf{2. Varianza (Variance)}: \[
\text{Var}[\hat{f}(\mathbf{X})] = \mathbb{E}[(\hat{f}(\mathbf{X}) - \mathbb{E}[\hat{f}(\mathbf{X})])^2]
\]

Mide cuánto varían las predicciones si se entrena con diferentes
muestras. Modelos complejos tienen alta varianza (overfitting).

\textbf{3. Error irreducible}: \[
\sigma^2 = \text{Var}(\varepsilon)
\]

Ruido inherente en los datos que ningún modelo puede reducir.

\subsubsection{Trade-off sesgo-varianza}\label{trade-off-sesgo-varianza}

Aumentar la complejidad del modelo:

\begin{itemize}
\tightlist
\item
  \textbf{Reduce sesgo}: El modelo se ajusta mejor a los datos
\item
  \textbf{Aumenta varianza}: El modelo se vuelve más sensible a cambios
  en los datos
\end{itemize}

\textbf{Objetivo}: Encontrar el nivel de complejidad que minimiza el
error total.

\subsubsection{Señales de overfitting}\label{seuxf1ales-de-overfitting}

\begin{enumerate}
\def\labelenumi{\arabic{enumi}.}
\item
  \textbf{Gap entre train y test}: \[
  \text{Performance}_{\text{train}} >> \text{Performance}_{\text{test}}
  \]
\item
  \textbf{Alta complejidad del modelo}: Muchos parámetros relativos a
  observaciones
\item
  \textbf{Rendimiento perfecto en training}: Accuracy = 100\%, R² = 1
\item
  \textbf{Inestabilidad}: Pequeños cambios en datos → grandes cambios en
  modelo
\end{enumerate}

\begin{Shaded}
\begin{Highlighting}[]
\ImportTok{import}\NormalTok{ pandas }\ImportTok{as}\NormalTok{ pd}
\ImportTok{import}\NormalTok{ numpy }\ImportTok{as}\NormalTok{ np}
\ImportTok{import}\NormalTok{ matplotlib.pyplot }\ImportTok{as}\NormalTok{ plt}
\ImportTok{import}\NormalTok{ seaborn }\ImportTok{as}\NormalTok{ sns}

\CommentTok{\# Scikit{-}learn}
\ImportTok{from}\NormalTok{ sklearn.datasets }\ImportTok{import}\NormalTok{ load\_iris}
\ImportTok{from}\NormalTok{ sklearn.model\_selection }\ImportTok{import}\NormalTok{ (}
\NormalTok{    train\_test\_split, cross\_val\_score, cross\_validate,}
\NormalTok{    GridSearchCV, RandomizedSearchCV, learning\_curve,}
\NormalTok{    validation\_curve, KFold, StratifiedKFold}
\NormalTok{)}
\ImportTok{from}\NormalTok{ sklearn.linear\_model }\ImportTok{import}\NormalTok{ LogisticRegression}
\ImportTok{from}\NormalTok{ sklearn.neighbors }\ImportTok{import}\NormalTok{ KNeighborsClassifier}
\ImportTok{from}\NormalTok{ sklearn.tree }\ImportTok{import}\NormalTok{ DecisionTreeClassifier}
\ImportTok{from}\NormalTok{ sklearn.ensemble }\ImportTok{import}\NormalTok{ RandomForestClassifier}
\ImportTok{from}\NormalTok{ sklearn.metrics }\ImportTok{import}\NormalTok{ (}
\NormalTok{    accuracy\_score, precision\_score, recall\_score,}
\NormalTok{    f1\_score, roc\_auc\_score, roc\_curve,}
\NormalTok{    confusion\_matrix, classification\_report}
\NormalTok{)}

\CommentTok{\# Configuración}
\NormalTok{plt.style.use(}\StringTok{\textquotesingle{}seaborn{-}v0\_8{-}whitegrid\textquotesingle{}}\NormalTok{)}
\NormalTok{plt.rcParams[}\StringTok{\textquotesingle{}figure.figsize\textquotesingle{}}\NormalTok{] }\OperatorTok{=}\NormalTok{ (}\DecValTok{12}\NormalTok{, }\DecValTok{6}\NormalTok{)}
\NormalTok{np.random.seed(}\DecValTok{42}\NormalTok{)}

\BuiltInTok{print}\NormalTok{(}\StringTok{"Librerías importadas exitosamente"}\NormalTok{)}
\end{Highlighting}
\end{Shaded}

\begin{verbatim}
Librerías importadas exitosamente
\end{verbatim}

\subsection{Ilustración del
overfitting}\label{ilustraciuxf3n-del-overfitting}

\begin{Shaded}
\begin{Highlighting}[]
\CommentTok{\# Cargar datos}
\NormalTok{iris }\OperatorTok{=}\NormalTok{ load\_iris()}
\NormalTok{X }\OperatorTok{=}\NormalTok{ iris.data[:, :}\DecValTok{2}\NormalTok{]  }\CommentTok{\# Solo 2 características para visualización}
\NormalTok{y }\OperatorTok{=}\NormalTok{ (iris.target }\OperatorTok{==} \DecValTok{2}\NormalTok{).astype(}\BuiltInTok{int}\NormalTok{)  }\CommentTok{\# Clasificación binaria}

\BuiltInTok{print}\NormalTok{(}\StringTok{"ILUSTRACIÓN DE OVERFITTING"}\NormalTok{)}
\BuiltInTok{print}\NormalTok{(}\StringTok{"="} \OperatorTok{*} \DecValTok{70}\NormalTok{)}

\CommentTok{\# Dividir datos}
\NormalTok{X\_train, X\_test, y\_train, y\_test }\OperatorTok{=}\NormalTok{ train\_test\_split(}
\NormalTok{    X, y, test\_size}\OperatorTok{=}\FloatTok{0.3}\NormalTok{, random\_state}\OperatorTok{=}\DecValTok{42}
\NormalTok{)}

\CommentTok{\# Entrenar modelos con diferente complejidad}
\NormalTok{modelos }\OperatorTok{=}\NormalTok{ \{}
    \StringTok{\textquotesingle{}Simple (k=50)\textquotesingle{}}\NormalTok{: KNeighborsClassifier(n\_neighbors}\OperatorTok{=}\DecValTok{50}\NormalTok{),}
    \StringTok{\textquotesingle{}Balanceado (k=5)\textquotesingle{}}\NormalTok{: KNeighborsClassifier(n\_neighbors}\OperatorTok{=}\DecValTok{5}\NormalTok{),}
    \StringTok{\textquotesingle{}Complejo (k=1)\textquotesingle{}}\NormalTok{: KNeighborsClassifier(n\_neighbors}\OperatorTok{=}\DecValTok{1}\NormalTok{)}
\NormalTok{\}}

\NormalTok{resultados }\OperatorTok{=}\NormalTok{ []}

\ControlFlowTok{for}\NormalTok{ nombre, modelo }\KeywordTok{in}\NormalTok{ modelos.items():}
\NormalTok{    modelo.fit(X\_train, y\_train)}
    
\NormalTok{    acc\_train }\OperatorTok{=}\NormalTok{ modelo.score(X\_train, y\_train)}
\NormalTok{    acc\_test }\OperatorTok{=}\NormalTok{ modelo.score(X\_test, y\_test)}
\NormalTok{    gap }\OperatorTok{=}\NormalTok{ acc\_train }\OperatorTok{{-}}\NormalTok{ acc\_test}
    
\NormalTok{    resultados.append(\{}
        \StringTok{\textquotesingle{}Modelo\textquotesingle{}}\NormalTok{: nombre,}
        \StringTok{\textquotesingle{}Accuracy Train\textquotesingle{}}\NormalTok{: acc\_train,}
        \StringTok{\textquotesingle{}Accuracy Test\textquotesingle{}}\NormalTok{: acc\_test,}
        \StringTok{\textquotesingle{}Gap (Train{-}Test)\textquotesingle{}}\NormalTok{: gap}
\NormalTok{    \})}
    
    \BuiltInTok{print}\NormalTok{(}\SpecialStringTok{f"}\CharTok{\textbackslash{}n}\SpecialCharTok{\{}\NormalTok{nombre}\SpecialCharTok{\}}\SpecialStringTok{:"}\NormalTok{)}
    \BuiltInTok{print}\NormalTok{(}\SpecialStringTok{f"  Train: }\SpecialCharTok{\{}\NormalTok{acc\_train}\SpecialCharTok{:.4f\}}\SpecialStringTok{"}\NormalTok{)}
    \BuiltInTok{print}\NormalTok{(}\SpecialStringTok{f"  Test:  }\SpecialCharTok{\{}\NormalTok{acc\_test}\SpecialCharTok{:.4f\}}\SpecialStringTok{"}\NormalTok{)}
    \BuiltInTok{print}\NormalTok{(}\SpecialStringTok{f"  Gap:   }\SpecialCharTok{\{}\NormalTok{gap}\SpecialCharTok{:.4f\}}\SpecialStringTok{"}\NormalTok{)}
    
    \ControlFlowTok{if}\NormalTok{ gap }\OperatorTok{\textgreater{}} \FloatTok{0.15}\NormalTok{:}
        \BuiltInTok{print}\NormalTok{(}\SpecialStringTok{f"  ⚠️  Posible overfitting detectado"}\NormalTok{)}
    \ControlFlowTok{elif}\NormalTok{ acc\_test }\OperatorTok{\textless{}} \FloatTok{0.70}\NormalTok{:}
        \BuiltInTok{print}\NormalTok{(}\SpecialStringTok{f"  ⚠️  Posible underfitting detectado"}\NormalTok{)}
    \ControlFlowTok{else}\NormalTok{:}
        \BuiltInTok{print}\NormalTok{(}\SpecialStringTok{f"  ✓ Balance adecuado"}\NormalTok{)}

\NormalTok{df\_resultados }\OperatorTok{=}\NormalTok{ pd.DataFrame(resultados)}
\end{Highlighting}
\end{Shaded}

\begin{verbatim}
ILUSTRACIÓN DE OVERFITTING
======================================================================

Simple (k=50):
  Train: 0.8190
  Test:  0.8222
  Gap:   -0.0032
  ✓ Balance adecuado

Balanceado (k=5):
  Train: 0.8476
  Test:  0.7778
  Gap:   0.0698
  ✓ Balance adecuado

Complejo (k=1):
  Train: 0.9619
  Test:  0.7111
  Gap:   0.2508
  ⚠️  Posible overfitting detectado
\end{verbatim}

\subsection{Visualización del
trade-off}\label{visualizaciuxf3n-del-trade-off}

\begin{Shaded}
\begin{Highlighting}[]
\CommentTok{\# Evaluar KNN con diferentes k}
\NormalTok{k\_values }\OperatorTok{=} \BuiltInTok{range}\NormalTok{(}\DecValTok{1}\NormalTok{, }\DecValTok{51}\NormalTok{)}
\NormalTok{train\_scores }\OperatorTok{=}\NormalTok{ []}
\NormalTok{test\_scores }\OperatorTok{=}\NormalTok{ []}

\ControlFlowTok{for}\NormalTok{ k }\KeywordTok{in}\NormalTok{ k\_values:}
\NormalTok{    knn }\OperatorTok{=}\NormalTok{ KNeighborsClassifier(n\_neighbors}\OperatorTok{=}\NormalTok{k)}
\NormalTok{    knn.fit(X\_train, y\_train)}
    
\NormalTok{    train\_scores.append(knn.score(X\_train, y\_train))}
\NormalTok{    test\_scores.append(knn.score(X\_test, y\_test))}

\CommentTok{\# Graficar}
\NormalTok{plt.figure(figsize}\OperatorTok{=}\NormalTok{(}\DecValTok{12}\NormalTok{, }\DecValTok{6}\NormalTok{))}

\NormalTok{plt.plot(k\_values, train\_scores, }\StringTok{\textquotesingle{}o{-}\textquotesingle{}}\NormalTok{, label}\OperatorTok{=}\StringTok{\textquotesingle{}Training\textquotesingle{}}\NormalTok{, linewidth}\OperatorTok{=}\DecValTok{2}\NormalTok{)}
\NormalTok{plt.plot(k\_values, test\_scores, }\StringTok{\textquotesingle{}s{-}\textquotesingle{}}\NormalTok{, label}\OperatorTok{=}\StringTok{\textquotesingle{}Test\textquotesingle{}}\NormalTok{, linewidth}\OperatorTok{=}\DecValTok{2}\NormalTok{)}

\CommentTok{\# Marcar mejor k}
\NormalTok{best\_k }\OperatorTok{=}\NormalTok{ k\_values[np.argmax(test\_scores)]}
\NormalTok{best\_score }\OperatorTok{=} \BuiltInTok{max}\NormalTok{(test\_scores)}
\NormalTok{plt.axvline(x}\OperatorTok{=}\NormalTok{best\_k, color}\OperatorTok{=}\StringTok{\textquotesingle{}red\textquotesingle{}}\NormalTok{, linestyle}\OperatorTok{=}\StringTok{\textquotesingle{}{-}{-}\textquotesingle{}}\NormalTok{, alpha}\OperatorTok{=}\FloatTok{0.7}\NormalTok{,}
\NormalTok{            label}\OperatorTok{=}\SpecialStringTok{f\textquotesingle{}Mejor k=}\SpecialCharTok{\{}\NormalTok{best\_k}\SpecialCharTok{\}}\SpecialStringTok{\textquotesingle{}}\NormalTok{)}

\NormalTok{plt.xlabel(}\StringTok{\textquotesingle{}k (número de vecinos)\textquotesingle{}}\NormalTok{)}
\NormalTok{plt.ylabel(}\StringTok{\textquotesingle{}Accuracy\textquotesingle{}}\NormalTok{)}
\NormalTok{plt.title(}\StringTok{\textquotesingle{}Trade{-}off Sesgo{-}Varianza en KNN\textquotesingle{}}\NormalTok{)}
\NormalTok{plt.legend()}
\NormalTok{plt.grid(}\VariableTok{True}\NormalTok{, alpha}\OperatorTok{=}\FloatTok{0.3}\NormalTok{)}

\CommentTok{\# Anotar regiones}
\NormalTok{plt.text(}\DecValTok{5}\NormalTok{, }\FloatTok{0.65}\NormalTok{, }\StringTok{\textquotesingle{}Alta varianza}\CharTok{\textbackslash{}n}\StringTok{(Overfitting)\textquotesingle{}}\NormalTok{, }
\NormalTok{         ha}\OperatorTok{=}\StringTok{\textquotesingle{}center\textquotesingle{}}\NormalTok{, fontsize}\OperatorTok{=}\DecValTok{10}\NormalTok{, bbox}\OperatorTok{=}\BuiltInTok{dict}\NormalTok{(boxstyle}\OperatorTok{=}\StringTok{\textquotesingle{}round\textquotesingle{}}\NormalTok{, facecolor}\OperatorTok{=}\StringTok{\textquotesingle{}red\textquotesingle{}}\NormalTok{, alpha}\OperatorTok{=}\FloatTok{0.3}\NormalTok{))}
\NormalTok{plt.text(}\DecValTok{40}\NormalTok{, }\FloatTok{0.65}\NormalTok{, }\StringTok{\textquotesingle{}Alto sesgo}\CharTok{\textbackslash{}n}\StringTok{(Underfitting)\textquotesingle{}}\NormalTok{, }
\NormalTok{         ha}\OperatorTok{=}\StringTok{\textquotesingle{}center\textquotesingle{}}\NormalTok{, fontsize}\OperatorTok{=}\DecValTok{10}\NormalTok{, bbox}\OperatorTok{=}\BuiltInTok{dict}\NormalTok{(boxstyle}\OperatorTok{=}\StringTok{\textquotesingle{}round\textquotesingle{}}\NormalTok{, facecolor}\OperatorTok{=}\StringTok{\textquotesingle{}blue\textquotesingle{}}\NormalTok{, alpha}\OperatorTok{=}\FloatTok{0.3}\NormalTok{))}

\NormalTok{plt.tight\_layout()}
\NormalTok{plt.savefig(}\StringTok{\textquotesingle{}bias\_variance\_tradeoff.png\textquotesingle{}}\NormalTok{, dpi}\OperatorTok{=}\DecValTok{300}\NormalTok{, bbox\_inches}\OperatorTok{=}\StringTok{\textquotesingle{}tight\textquotesingle{}}\NormalTok{)}
\BuiltInTok{print}\NormalTok{(}\StringTok{"}\CharTok{\textbackslash{}n}\StringTok{Gráfico guardado como \textquotesingle{}bias\_variance\_tradeoff.png\textquotesingle{}"}\NormalTok{)}
\end{Highlighting}
\end{Shaded}

\begin{verbatim}

Gráfico guardado como 'bias_variance_tradeoff.png'
\end{verbatim}

\pandocbounded{\includegraphics[keepaspectratio]{index_files/figure-pdf/cell-4-output-2.pdf}}

\section{Validación cruzada
(Cross-Validation)}\label{validaciuxf3n-cruzada-cross-validation}

\subsection{Fundamento matemático}\label{fundamento-matemuxe1tico}

La validación cruzada es una técnica para estimar el rendimiento de un
modelo en datos no vistos mediante la división estratégica del conjunto
de datos disponible.

\subsubsection{Problema con validación
simple}\label{problema-con-validaciuxf3n-simple}

Con una sola división train-test:

\begin{itemize}
\tightlist
\item
  El resultado depende de qué observaciones caen en cada conjunto
\item
  Alta varianza en la estimación del error
\item
  Desperdicio de datos (test set no se usa para entrenar)
\end{itemize}

\subsubsection{K-Fold Cross-Validation}\label{k-fold-cross-validation}

\textbf{Algoritmo}:

\begin{enumerate}
\def\labelenumi{\arabic{enumi}.}
\item
  Dividir datos en \(K\) folds (particiones) de tamaño aproximadamente
  igual
\item
  Para cada fold \(k = 1, \ldots, K\):

  \begin{itemize}
  \tightlist
  \item
    Entrenar modelo en \(K-1\) folds
  \item
    Validar en fold \(k\)
  \item
    Calcular métrica de performance \(\text{Score}_k\)
  \end{itemize}
\item
  Promedio de scores: \[
  \text{CV Score} = \frac{1}{K} \sum_{k=1}^{K} \text{Score}_k
  \]
\item
  Desviación estándar: \[
  \text{SD} = \sqrt{\frac{1}{K-1} \sum_{k=1}^{K} (\text{Score}_k - \overline{\text{Score}})^2}
  \]
\end{enumerate}

\textbf{Propiedades}:

\begin{itemize}
\item
  \textbf{Sesgo}: \[
  \mathbb{E}[\text{CV Score}] \approx \text{Error verdadero}
  \] Estimación casi insesgada del error de generalización
\item
  \textbf{Varianza}: Disminuye con \(K\) más grande
\item
  \textbf{Trade-off K}:

  \begin{itemize}
  \tightlist
  \item
    \(K\) pequeño (ej: 5): Más sesgo, menos varianza, más rápido
  \item
    \(K\) grande (ej: 10): Menos sesgo, más varianza, más lento
  \item
    \(K = n\) (LOOCV): Sesgo mínimo, varianza máxima, muy lento
  \end{itemize}
\end{itemize}

\subsubsection{Variantes de
Cross-Validation}\label{variantes-de-cross-validation}

\textbf{1. Stratified K-Fold}:

Mantiene proporción de clases en cada fold: \[
\frac{n_k^{(c)}}{n_k} \approx \frac{n^{(c)}}{n} \quad \forall k, c
\]

donde \(n_k^{(c)}\) es número de observaciones de clase \(c\) en fold
\(k\).

\textbf{Crucial para}: Datasets desbalanceados

\textbf{2. Leave-One-Out CV (LOOCV)}:

Caso especial donde \(K = n\): \[
\text{LOOCV} = \frac{1}{n} \sum_{i=1}^{n} L(y_i, \hat{f}^{(-i)}(\mathbf{X}_i))
\]

donde \(\hat{f}^{(-i)}\) es modelo entrenado sin observación \(i\).

\textbf{Ventajas}: Sesgo mínimo \textbf{Desventajas}: Computacionalmente
costoso, alta varianza

\textbf{3. Repeated K-Fold}:

Repetir K-Fold con diferentes particiones aleatorias y promediar.

\subsection{Implementación básica}\label{implementaciuxf3n-buxe1sica}

\begin{Shaded}
\begin{Highlighting}[]
\BuiltInTok{print}\NormalTok{(}\StringTok{"}\CharTok{\textbackslash{}n\textbackslash{}n}\StringTok{VALIDACIÓN CRUZADA"}\NormalTok{)}
\BuiltInTok{print}\NormalTok{(}\StringTok{"="} \OperatorTok{*} \DecValTok{70}\NormalTok{)}

\CommentTok{\# Modelo a evaluar}
\NormalTok{modelo }\OperatorTok{=}\NormalTok{ LogisticRegression(random\_state}\OperatorTok{=}\DecValTok{42}\NormalTok{)}

\CommentTok{\# K{-}Fold CV estándar}
\NormalTok{kfold }\OperatorTok{=}\NormalTok{ KFold(n\_splits}\OperatorTok{=}\DecValTok{5}\NormalTok{, shuffle}\OperatorTok{=}\VariableTok{True}\NormalTok{, random\_state}\OperatorTok{=}\DecValTok{42}\NormalTok{)}

\NormalTok{scores }\OperatorTok{=}\NormalTok{ cross\_val\_score(}
\NormalTok{    modelo, X, y, }
\NormalTok{    cv}\OperatorTok{=}\NormalTok{kfold, }
\NormalTok{    scoring}\OperatorTok{=}\StringTok{\textquotesingle{}accuracy\textquotesingle{}}
\NormalTok{)}

\BuiltInTok{print}\NormalTok{(}\StringTok{"}\CharTok{\textbackslash{}n}\StringTok{K{-}Fold Cross{-}Validation (K=5):"}\NormalTok{)}
\BuiltInTok{print}\NormalTok{(}\SpecialStringTok{f"Scores por fold: }\SpecialCharTok{\{}\NormalTok{scores}\SpecialCharTok{\}}\SpecialStringTok{"}\NormalTok{)}
\BuiltInTok{print}\NormalTok{(}\SpecialStringTok{f"Media: }\SpecialCharTok{\{}\NormalTok{scores}\SpecialCharTok{.}\NormalTok{mean()}\SpecialCharTok{:.4f\}}\SpecialStringTok{"}\NormalTok{)}
\BuiltInTok{print}\NormalTok{(}\SpecialStringTok{f"Desviación estándar: }\SpecialCharTok{\{}\NormalTok{scores}\SpecialCharTok{.}\NormalTok{std()}\SpecialCharTok{:.4f\}}\SpecialStringTok{"}\NormalTok{)}
\BuiltInTok{print}\NormalTok{(}\SpecialStringTok{f"Intervalo 95\% confianza: [}\SpecialCharTok{\{}\NormalTok{scores}\SpecialCharTok{.}\NormalTok{mean() }\OperatorTok{{-}} \FloatTok{1.96}\OperatorTok{*}\NormalTok{scores}\SpecialCharTok{.}\NormalTok{std()}\SpecialCharTok{:.4f\}}\SpecialStringTok{, "}
      \SpecialStringTok{f"}\SpecialCharTok{\{}\NormalTok{scores}\SpecialCharTok{.}\NormalTok{mean() }\OperatorTok{+} \FloatTok{1.96}\OperatorTok{*}\NormalTok{scores}\SpecialCharTok{.}\NormalTok{std()}\SpecialCharTok{:.4f\}}\SpecialStringTok{]"}\NormalTok{)}
\end{Highlighting}
\end{Shaded}

\begin{verbatim}


VALIDACIÓN CRUZADA
======================================================================

K-Fold Cross-Validation (K=5):
Scores por fold: [0.9        0.83333333 0.76666667 0.7        0.73333333]
Media: 0.7867
Desviación estándar: 0.0718
Intervalo 95% confianza: [0.6459, 0.9274]
\end{verbatim}

\subsection{Stratified
Cross-Validation}\label{stratified-cross-validation}

\begin{Shaded}
\begin{Highlighting}[]
\CommentTok{\# Verificar distribución de clases}
\BuiltInTok{print}\NormalTok{(}\SpecialStringTok{f"}\CharTok{\textbackslash{}n\textbackslash{}n}\SpecialStringTok{Distribución de clases:"}\NormalTok{)}
\BuiltInTok{print}\NormalTok{(}\SpecialStringTok{f"Clase 0: }\SpecialCharTok{\{}\NormalTok{(y}\OperatorTok{==}\DecValTok{0}\NormalTok{)}\SpecialCharTok{.}\BuiltInTok{sum}\NormalTok{()}\SpecialCharTok{\}}\SpecialStringTok{ (}\SpecialCharTok{\{}\NormalTok{(y}\OperatorTok{==}\DecValTok{0}\NormalTok{)}\SpecialCharTok{.}\BuiltInTok{sum}\NormalTok{()}\OperatorTok{/}\BuiltInTok{len}\NormalTok{(y)}\OperatorTok{*}\DecValTok{100}\SpecialCharTok{:.1f\}}\SpecialStringTok{\%)"}\NormalTok{)}
\BuiltInTok{print}\NormalTok{(}\SpecialStringTok{f"Clase 1: }\SpecialCharTok{\{}\NormalTok{(y}\OperatorTok{==}\DecValTok{1}\NormalTok{)}\SpecialCharTok{.}\BuiltInTok{sum}\NormalTok{()}\SpecialCharTok{\}}\SpecialStringTok{ (}\SpecialCharTok{\{}\NormalTok{(y}\OperatorTok{==}\DecValTok{1}\NormalTok{)}\SpecialCharTok{.}\BuiltInTok{sum}\NormalTok{()}\OperatorTok{/}\BuiltInTok{len}\NormalTok{(y)}\OperatorTok{*}\DecValTok{100}\SpecialCharTok{:.1f\}}\SpecialStringTok{\%)"}\NormalTok{)}

\CommentTok{\# Stratified K{-}Fold}
\NormalTok{stratified\_kfold }\OperatorTok{=}\NormalTok{ StratifiedKFold(n\_splits}\OperatorTok{=}\DecValTok{5}\NormalTok{, shuffle}\OperatorTok{=}\VariableTok{True}\NormalTok{, random\_state}\OperatorTok{=}\DecValTok{42}\NormalTok{)}

\NormalTok{scores\_stratified }\OperatorTok{=}\NormalTok{ cross\_val\_score(}
\NormalTok{    modelo, X, y,}
\NormalTok{    cv}\OperatorTok{=}\NormalTok{stratified\_kfold,}
\NormalTok{    scoring}\OperatorTok{=}\StringTok{\textquotesingle{}accuracy\textquotesingle{}}
\NormalTok{)}

\BuiltInTok{print}\NormalTok{(}\SpecialStringTok{f"}\CharTok{\textbackslash{}n}\SpecialStringTok{Stratified K{-}Fold CV:"}\NormalTok{)}
\BuiltInTok{print}\NormalTok{(}\SpecialStringTok{f"Media: }\SpecialCharTok{\{}\NormalTok{scores\_stratified}\SpecialCharTok{.}\NormalTok{mean()}\SpecialCharTok{:.4f\}}\SpecialStringTok{"}\NormalTok{)}
\BuiltInTok{print}\NormalTok{(}\SpecialStringTok{f"Desv. est: }\SpecialCharTok{\{}\NormalTok{scores\_stratified}\SpecialCharTok{.}\NormalTok{std()}\SpecialCharTok{:.4f\}}\SpecialStringTok{"}\NormalTok{)}

\CommentTok{\# Comparación}
\BuiltInTok{print}\NormalTok{(}\SpecialStringTok{f"}\CharTok{\textbackslash{}n}\SpecialStringTok{Diferencia con K{-}Fold estándar: "}
      \SpecialStringTok{f"}\SpecialCharTok{\{}\NormalTok{scores\_stratified}\SpecialCharTok{.}\NormalTok{mean() }\OperatorTok{{-}}\NormalTok{ scores}\SpecialCharTok{.}\NormalTok{mean()}\SpecialCharTok{:.4f\}}\SpecialStringTok{"}\NormalTok{)}
\end{Highlighting}
\end{Shaded}

\begin{verbatim}


Distribución de clases:
Clase 0: 100 (66.7%)
Clase 1: 50 (33.3%)

Stratified K-Fold CV:
Media: 0.7867
Desv. est: 0.0686

Diferencia con K-Fold estándar: 0.0000
\end{verbatim}

\subsection{Cross-validation con múltiples
métricas}\label{cross-validation-con-muxfaltiples-muxe9tricas}

\begin{Shaded}
\begin{Highlighting}[]
\CommentTok{\# Evaluar múltiples métricas simultáneamente}
\NormalTok{scoring }\OperatorTok{=}\NormalTok{ \{}
    \StringTok{\textquotesingle{}accuracy\textquotesingle{}}\NormalTok{: }\StringTok{\textquotesingle{}accuracy\textquotesingle{}}\NormalTok{,}
    \StringTok{\textquotesingle{}precision\textquotesingle{}}\NormalTok{: }\StringTok{\textquotesingle{}precision\textquotesingle{}}\NormalTok{,}
    \StringTok{\textquotesingle{}recall\textquotesingle{}}\NormalTok{: }\StringTok{\textquotesingle{}recall\textquotesingle{}}\NormalTok{,}
    \StringTok{\textquotesingle{}f1\textquotesingle{}}\NormalTok{: }\StringTok{\textquotesingle{}f1\textquotesingle{}}\NormalTok{,}
    \StringTok{\textquotesingle{}roc\_auc\textquotesingle{}}\NormalTok{: }\StringTok{\textquotesingle{}roc\_auc\textquotesingle{}}
\NormalTok{\}}

\NormalTok{cv\_results }\OperatorTok{=}\NormalTok{ cross\_validate(}
\NormalTok{    modelo, X, y,}
\NormalTok{    cv}\OperatorTok{=}\NormalTok{stratified\_kfold,}
\NormalTok{    scoring}\OperatorTok{=}\NormalTok{scoring,}
\NormalTok{    return\_train\_score}\OperatorTok{=}\VariableTok{True}
\NormalTok{)}

\BuiltInTok{print}\NormalTok{(}\StringTok{"}\CharTok{\textbackslash{}n\textbackslash{}n}\StringTok{RESULTADOS CROSS{-}VALIDATION {-} MÚLTIPLES MÉTRICAS"}\NormalTok{)}
\BuiltInTok{print}\NormalTok{(}\StringTok{"="} \OperatorTok{*} \DecValTok{70}\NormalTok{)}

\ControlFlowTok{for}\NormalTok{ metric }\KeywordTok{in}\NormalTok{ [}\StringTok{\textquotesingle{}accuracy\textquotesingle{}}\NormalTok{, }\StringTok{\textquotesingle{}precision\textquotesingle{}}\NormalTok{, }\StringTok{\textquotesingle{}recall\textquotesingle{}}\NormalTok{, }\StringTok{\textquotesingle{}f1\textquotesingle{}}\NormalTok{, }\StringTok{\textquotesingle{}roc\_auc\textquotesingle{}}\NormalTok{]:}
\NormalTok{    train\_scores }\OperatorTok{=}\NormalTok{ cv\_results[}\SpecialStringTok{f\textquotesingle{}train\_}\SpecialCharTok{\{}\NormalTok{metric}\SpecialCharTok{\}}\SpecialStringTok{\textquotesingle{}}\NormalTok{]}
\NormalTok{    test\_scores }\OperatorTok{=}\NormalTok{ cv\_results[}\SpecialStringTok{f\textquotesingle{}test\_}\SpecialCharTok{\{}\NormalTok{metric}\SpecialCharTok{\}}\SpecialStringTok{\textquotesingle{}}\NormalTok{]}
    
    \BuiltInTok{print}\NormalTok{(}\SpecialStringTok{f"}\CharTok{\textbackslash{}n}\SpecialCharTok{\{}\NormalTok{metric}\SpecialCharTok{.}\NormalTok{upper()}\SpecialCharTok{\}}\SpecialStringTok{:"}\NormalTok{)}
    \BuiltInTok{print}\NormalTok{(}\SpecialStringTok{f"  Train: }\SpecialCharTok{\{}\NormalTok{train\_scores}\SpecialCharTok{.}\NormalTok{mean()}\SpecialCharTok{:.4f\}}\SpecialStringTok{ ± }\SpecialCharTok{\{}\NormalTok{train\_scores}\SpecialCharTok{.}\NormalTok{std()}\SpecialCharTok{:.4f\}}\SpecialStringTok{"}\NormalTok{)}
    \BuiltInTok{print}\NormalTok{(}\SpecialStringTok{f"  Test:  }\SpecialCharTok{\{}\NormalTok{test\_scores}\SpecialCharTok{.}\NormalTok{mean()}\SpecialCharTok{:.4f\}}\SpecialStringTok{ ± }\SpecialCharTok{\{}\NormalTok{test\_scores}\SpecialCharTok{.}\NormalTok{std()}\SpecialCharTok{:.4f\}}\SpecialStringTok{"}\NormalTok{)}
    \BuiltInTok{print}\NormalTok{(}\SpecialStringTok{f"  Gap:   }\SpecialCharTok{\{}\NormalTok{train\_scores}\SpecialCharTok{.}\NormalTok{mean() }\OperatorTok{{-}}\NormalTok{ test\_scores}\SpecialCharTok{.}\NormalTok{mean()}\SpecialCharTok{:.4f\}}\SpecialStringTok{"}\NormalTok{)}
\end{Highlighting}
\end{Shaded}

\begin{verbatim}


RESULTADOS CROSS-VALIDATION - MÚLTIPLES MÉTRICAS
======================================================================

ACCURACY:
  Train: 0.8067 ± 0.0178
  Test:  0.7867 ± 0.0686
  Gap:   0.0200

PRECISION:
  Train: 0.7260 ± 0.0285
  Test:  0.6984 ± 0.0893
  Gap:   0.0276

RECALL:
  Train: 0.6750 ± 0.0316
  Test:  0.6400 ± 0.2154
  Gap:   0.0350

F1:
  Train: 0.6994 ± 0.0282
  Test:  0.6479 ± 0.1519
  Gap:   0.0514

ROC_AUC:
  Train: 0.8888 ± 0.0162
  Test:  0.8935 ± 0.0611
  Gap:   -0.0047
\end{verbatim}

\subsection{Visualización de resultados
CV}\label{visualizaciuxf3n-de-resultados-cv}

\begin{Shaded}
\begin{Highlighting}[]
\CommentTok{\# Crear DataFrame para visualización}
\NormalTok{metrics\_data }\OperatorTok{=}\NormalTok{ []}

\ControlFlowTok{for}\NormalTok{ metric }\KeywordTok{in}\NormalTok{ [}\StringTok{\textquotesingle{}accuracy\textquotesingle{}}\NormalTok{, }\StringTok{\textquotesingle{}precision\textquotesingle{}}\NormalTok{, }\StringTok{\textquotesingle{}recall\textquotesingle{}}\NormalTok{, }\StringTok{\textquotesingle{}f1\textquotesingle{}}\NormalTok{, }\StringTok{\textquotesingle{}roc\_auc\textquotesingle{}}\NormalTok{]:}
    \ControlFlowTok{for}\NormalTok{ fold }\KeywordTok{in} \BuiltInTok{range}\NormalTok{(}\DecValTok{5}\NormalTok{):}
\NormalTok{        metrics\_data.append(\{}
            \StringTok{\textquotesingle{}Métrica\textquotesingle{}}\NormalTok{: metric,}
            \StringTok{\textquotesingle{}Fold\textquotesingle{}}\NormalTok{: fold }\OperatorTok{+} \DecValTok{1}\NormalTok{,}
            \StringTok{\textquotesingle{}Train\textquotesingle{}}\NormalTok{: cv\_results[}\SpecialStringTok{f\textquotesingle{}train\_}\SpecialCharTok{\{}\NormalTok{metric}\SpecialCharTok{\}}\SpecialStringTok{\textquotesingle{}}\NormalTok{][fold],}
            \StringTok{\textquotesingle{}Test\textquotesingle{}}\NormalTok{: cv\_results[}\SpecialStringTok{f\textquotesingle{}test\_}\SpecialCharTok{\{}\NormalTok{metric}\SpecialCharTok{\}}\SpecialStringTok{\textquotesingle{}}\NormalTok{][fold]}
\NormalTok{        \})}

\NormalTok{df\_metrics }\OperatorTok{=}\NormalTok{ pd.DataFrame(metrics\_data)}

\CommentTok{\# Graficar}
\NormalTok{fig, axes }\OperatorTok{=}\NormalTok{ plt.subplots(}\DecValTok{2}\NormalTok{, }\DecValTok{3}\NormalTok{, figsize}\OperatorTok{=}\NormalTok{(}\DecValTok{15}\NormalTok{, }\DecValTok{10}\NormalTok{))}
\NormalTok{axes }\OperatorTok{=}\NormalTok{ axes.ravel()}

\ControlFlowTok{for}\NormalTok{ idx, metric }\KeywordTok{in} \BuiltInTok{enumerate}\NormalTok{([}\StringTok{\textquotesingle{}accuracy\textquotesingle{}}\NormalTok{, }\StringTok{\textquotesingle{}precision\textquotesingle{}}\NormalTok{, }\StringTok{\textquotesingle{}recall\textquotesingle{}}\NormalTok{, }\StringTok{\textquotesingle{}f1\textquotesingle{}}\NormalTok{, }\StringTok{\textquotesingle{}roc\_auc\textquotesingle{}}\NormalTok{]):}
\NormalTok{    data }\OperatorTok{=}\NormalTok{ df\_metrics[df\_metrics[}\StringTok{\textquotesingle{}Métrica\textquotesingle{}}\NormalTok{] }\OperatorTok{==}\NormalTok{ metric]}
    
\NormalTok{    x }\OperatorTok{=}\NormalTok{ np.arange(}\DecValTok{1}\NormalTok{, }\DecValTok{6}\NormalTok{)}
\NormalTok{    axes[idx].plot(x, data[}\StringTok{\textquotesingle{}Train\textquotesingle{}}\NormalTok{].values, }\StringTok{\textquotesingle{}o{-}\textquotesingle{}}\NormalTok{, label}\OperatorTok{=}\StringTok{\textquotesingle{}Train\textquotesingle{}}\NormalTok{, linewidth}\OperatorTok{=}\DecValTok{2}\NormalTok{)}
\NormalTok{    axes[idx].plot(x, data[}\StringTok{\textquotesingle{}Test\textquotesingle{}}\NormalTok{].values, }\StringTok{\textquotesingle{}s{-}\textquotesingle{}}\NormalTok{, label}\OperatorTok{=}\StringTok{\textquotesingle{}Test\textquotesingle{}}\NormalTok{, linewidth}\OperatorTok{=}\DecValTok{2}\NormalTok{)}
    
\NormalTok{    axes[idx].set\_xlabel(}\StringTok{\textquotesingle{}Fold\textquotesingle{}}\NormalTok{)}
\NormalTok{    axes[idx].set\_ylabel(}\StringTok{\textquotesingle{}Score\textquotesingle{}}\NormalTok{)}
\NormalTok{    axes[idx].set\_title(}\SpecialStringTok{f\textquotesingle{}}\SpecialCharTok{\{}\NormalTok{metric}\SpecialCharTok{.}\NormalTok{upper()}\SpecialCharTok{\}}\SpecialStringTok{\textquotesingle{}}\NormalTok{)}
\NormalTok{    axes[idx].legend()}
\NormalTok{    axes[idx].grid(}\VariableTok{True}\NormalTok{, alpha}\OperatorTok{=}\FloatTok{0.3}\NormalTok{)}
\NormalTok{    axes[idx].set\_ylim([}\DecValTok{0}\NormalTok{, }\FloatTok{1.05}\NormalTok{])}

\CommentTok{\# Ocultar último subplot}
\NormalTok{axes[}\DecValTok{5}\NormalTok{].axis(}\StringTok{\textquotesingle{}off\textquotesingle{}}\NormalTok{)}

\NormalTok{plt.tight\_layout()}
\NormalTok{plt.savefig(}\StringTok{\textquotesingle{}cv\_multiple\_metrics.png\textquotesingle{}}\NormalTok{, dpi}\OperatorTok{=}\DecValTok{300}\NormalTok{, bbox\_inches}\OperatorTok{=}\StringTok{\textquotesingle{}tight\textquotesingle{}}\NormalTok{)}
\BuiltInTok{print}\NormalTok{(}\StringTok{"}\CharTok{\textbackslash{}n}\StringTok{Gráfico guardado como \textquotesingle{}cv\_multiple\_metrics.png\textquotesingle{}"}\NormalTok{)}
\end{Highlighting}
\end{Shaded}

\begin{verbatim}

Gráfico guardado como 'cv_multiple_metrics.png'
\end{verbatim}

\pandocbounded{\includegraphics[keepaspectratio]{index_files/figure-pdf/cell-8-output-2.pdf}}

\section{Grid Search y optimización de
hiperparámetros}\label{grid-search-y-optimizaciuxf3n-de-hiperparuxe1metros}

\subsection{Fundamento matemático}\label{fundamento-matemuxe1tico-1}

La búsqueda de hiperparámetros óptimos es un problema de optimización
sobre un espacio discreto de configuraciones.

\subsubsection{Formulación del
problema}\label{formulaciuxf3n-del-problema}

Dado:

\begin{itemize}
\tightlist
\item
  Modelo \(f(\mathbf{X}; \boldsymbol{\theta}, \boldsymbol{\lambda})\)

  \begin{itemize}
  \tightlist
  \item
    \(\boldsymbol{\theta}\): Parámetros (aprendidos de datos)
  \item
    \(\boldsymbol{\lambda}\): Hiperparámetros (especificados a priori)
  \end{itemize}
\item
  Espacio de hiperparámetros
  \(\Lambda = \Lambda_1 \times \Lambda_2 \times \cdots \times \Lambda_m\)
\item
  Métrica de performance \(S(\boldsymbol{\lambda})\) estimada por CV
\end{itemize}

\textbf{Objetivo}: Encontrar \[
\boldsymbol{\lambda}^* = \arg\max_{\boldsymbol{\lambda} \in \Lambda} S(\boldsymbol{\lambda})
\]

\subsubsection{Grid Search (Búsqueda
exhaustiva)}\label{grid-search-buxfasqueda-exhaustiva}

\textbf{Algoritmo}:

\begin{enumerate}
\def\labelenumi{\arabic{enumi}.}
\item
  Definir grilla de hiperparámetros: \[
  \Lambda = \{\lambda_1^{(1)}, \ldots, \lambda_1^{(n_1)}\} \times \cdots \times \{\lambda_m^{(1)}, \ldots, \lambda_m^{(n_m)}\}
  \]
\item
  Para cada combinación \(\boldsymbol{\lambda} \in \Lambda\):

  \begin{itemize}
  \tightlist
  \item
    Entrenar modelo con CV
  \item
    Calcular \(S(\boldsymbol{\lambda})\)
  \end{itemize}
\item
  Seleccionar \(\boldsymbol{\lambda}^*\) con mejor \(S\)
\end{enumerate}

\textbf{Complejidad}:
\(O(n_1 \times n_2 \times \cdots \times n_m \times K)\)

donde \(K\) es número de folds.

\textbf{Ventajas}:

\begin{itemize}
\tightlist
\item
  Exhaustivo: Garantiza encontrar mejor combinación en grilla
\item
  Paralelizable: Evaluaciones independientes
\end{itemize}

\textbf{Desventajas}:

\begin{itemize}
\tightlist
\item
  Costoso: Crece exponencialmente con número de hiperparámetros
\item
  Solo explora puntos discretos
\end{itemize}

\subsubsection{Random Search}\label{random-search}

En lugar de grilla, muestrear aleatoriamente del espacio de
hiperparámetros:

\[
\boldsymbol{\lambda}^{(i)} \sim P(\boldsymbol{\lambda}), \quad i = 1, \ldots, N
\]

\textbf{Ventajas sobre Grid Search}:

\begin{itemize}
\tightlist
\item
  Más eficiente: Puede encontrar buenos valores con menos evaluaciones
\item
  Explora mejor espacios continuos
\item
  Menos sensible a hiperparámetros irrelevantes
\end{itemize}

\textbf{Teorema (Bergstra \& Bengio, 2012)}: Random Search converge al
óptimo más rápido que Grid Search cuando pocos hiperparámetros importan.

\subsection{Implementación de Grid
Search}\label{implementaciuxf3n-de-grid-search}

\begin{Shaded}
\begin{Highlighting}[]
\BuiltInTok{print}\NormalTok{(}\StringTok{"}\CharTok{\textbackslash{}n\textbackslash{}n}\StringTok{GRID SEARCH {-} OPTIMIZACIÓN DE HIPERPARÁMETROS"}\NormalTok{)}
\BuiltInTok{print}\NormalTok{(}\StringTok{"="} \OperatorTok{*} \DecValTok{70}\NormalTok{)}

\CommentTok{\# Definir modelo y grilla de hiperparámetros}
\NormalTok{modelo }\OperatorTok{=}\NormalTok{ RandomForestClassifier(random\_state}\OperatorTok{=}\DecValTok{42}\NormalTok{)}

\NormalTok{param\_grid }\OperatorTok{=}\NormalTok{ \{}
    \StringTok{\textquotesingle{}n\_estimators\textquotesingle{}}\NormalTok{: [}\DecValTok{50}\NormalTok{, }\DecValTok{100}\NormalTok{, }\DecValTok{200}\NormalTok{],}
    \StringTok{\textquotesingle{}max\_depth\textquotesingle{}}\NormalTok{: [}\DecValTok{3}\NormalTok{, }\DecValTok{5}\NormalTok{, }\DecValTok{10}\NormalTok{, }\VariableTok{None}\NormalTok{],}
    \StringTok{\textquotesingle{}min\_samples\_split\textquotesingle{}}\NormalTok{: [}\DecValTok{2}\NormalTok{, }\DecValTok{5}\NormalTok{, }\DecValTok{10}\NormalTok{],}
    \StringTok{\textquotesingle{}min\_samples\_leaf\textquotesingle{}}\NormalTok{: [}\DecValTok{1}\NormalTok{, }\DecValTok{2}\NormalTok{, }\DecValTok{4}\NormalTok{]}
\NormalTok{\}}

\BuiltInTok{print}\NormalTok{(}\SpecialStringTok{f"Espacio de búsqueda:"}\NormalTok{)}
\BuiltInTok{print}\NormalTok{(}\SpecialStringTok{f"  n\_estimators: }\SpecialCharTok{\{}\NormalTok{param\_grid[}\StringTok{\textquotesingle{}n\_estimators\textquotesingle{}}\NormalTok{]}\SpecialCharTok{\}}\SpecialStringTok{"}\NormalTok{)}
\BuiltInTok{print}\NormalTok{(}\SpecialStringTok{f"  max\_depth: }\SpecialCharTok{\{}\NormalTok{param\_grid[}\StringTok{\textquotesingle{}max\_depth\textquotesingle{}}\NormalTok{]}\SpecialCharTok{\}}\SpecialStringTok{"}\NormalTok{)}
\BuiltInTok{print}\NormalTok{(}\SpecialStringTok{f"  min\_samples\_split: }\SpecialCharTok{\{}\NormalTok{param\_grid[}\StringTok{\textquotesingle{}min\_samples\_split\textquotesingle{}}\NormalTok{]}\SpecialCharTok{\}}\SpecialStringTok{"}\NormalTok{)}
\BuiltInTok{print}\NormalTok{(}\SpecialStringTok{f"  min\_samples\_leaf: }\SpecialCharTok{\{}\NormalTok{param\_grid[}\StringTok{\textquotesingle{}min\_samples\_leaf\textquotesingle{}}\NormalTok{]}\SpecialCharTok{\}}\SpecialStringTok{"}\NormalTok{)}

\NormalTok{n\_combinations }\OperatorTok{=}\NormalTok{ (}\BuiltInTok{len}\NormalTok{(param\_grid[}\StringTok{\textquotesingle{}n\_estimators\textquotesingle{}}\NormalTok{]) }\OperatorTok{*} 
                  \BuiltInTok{len}\NormalTok{(param\_grid[}\StringTok{\textquotesingle{}max\_depth\textquotesingle{}}\NormalTok{]) }\OperatorTok{*}
                  \BuiltInTok{len}\NormalTok{(param\_grid[}\StringTok{\textquotesingle{}min\_samples\_split\textquotesingle{}}\NormalTok{]) }\OperatorTok{*}
                  \BuiltInTok{len}\NormalTok{(param\_grid[}\StringTok{\textquotesingle{}min\_samples\_leaf\textquotesingle{}}\NormalTok{]))}

\BuiltInTok{print}\NormalTok{(}\SpecialStringTok{f"}\CharTok{\textbackslash{}n}\SpecialStringTok{Total de combinaciones: }\SpecialCharTok{\{}\NormalTok{n\_combinations}\SpecialCharTok{\}}\SpecialStringTok{"}\NormalTok{)}
\BuiltInTok{print}\NormalTok{(}\SpecialStringTok{f"Con CV de 5{-}fold: }\SpecialCharTok{\{}\NormalTok{n\_combinations }\OperatorTok{*} \DecValTok{5}\SpecialCharTok{\}}\SpecialStringTok{ modelos a entrenar"}\NormalTok{)}
\end{Highlighting}
\end{Shaded}

\begin{verbatim}


GRID SEARCH - OPTIMIZACIÓN DE HIPERPARÁMETROS
======================================================================
Espacio de búsqueda:
  n_estimators: [50, 100, 200]
  max_depth: [3, 5, 10, None]
  min_samples_split: [2, 5, 10]
  min_samples_leaf: [1, 2, 4]

Total de combinaciones: 108
Con CV de 5-fold: 540 modelos a entrenar
\end{verbatim}

\begin{Shaded}
\begin{Highlighting}[]
\CommentTok{\# Ejecutar Grid Search}
\NormalTok{grid\_search }\OperatorTok{=}\NormalTok{ GridSearchCV(}
\NormalTok{    estimator}\OperatorTok{=}\NormalTok{modelo,}
\NormalTok{    param\_grid}\OperatorTok{=}\NormalTok{param\_grid,}
\NormalTok{    cv}\OperatorTok{=}\NormalTok{StratifiedKFold(n\_splits}\OperatorTok{=}\DecValTok{5}\NormalTok{, shuffle}\OperatorTok{=}\VariableTok{True}\NormalTok{, random\_state}\OperatorTok{=}\DecValTok{42}\NormalTok{),}
\NormalTok{    scoring}\OperatorTok{=}\StringTok{\textquotesingle{}roc\_auc\textquotesingle{}}\NormalTok{,}
\NormalTok{    n\_jobs}\OperatorTok{={-}}\DecValTok{1}\NormalTok{,  }\CommentTok{\# Usar todos los cores}
\NormalTok{    verbose}\OperatorTok{=}\DecValTok{1}\NormalTok{,}
\NormalTok{    return\_train\_score}\OperatorTok{=}\VariableTok{True}
\NormalTok{)}

\BuiltInTok{print}\NormalTok{(}\StringTok{"}\CharTok{\textbackslash{}n}\StringTok{Ejecutando Grid Search..."}\NormalTok{)}
\NormalTok{grid\_search.fit(X, y)}

\BuiltInTok{print}\NormalTok{(}\StringTok{"}\CharTok{\textbackslash{}n\textbackslash{}n}\StringTok{RESULTADOS GRID SEARCH"}\NormalTok{)}
\BuiltInTok{print}\NormalTok{(}\StringTok{"="} \OperatorTok{*} \DecValTok{70}\NormalTok{)}
\BuiltInTok{print}\NormalTok{(}\SpecialStringTok{f"}\CharTok{\textbackslash{}n}\SpecialStringTok{Mejores hiperparámetros:"}\NormalTok{)}
\ControlFlowTok{for}\NormalTok{ param, value }\KeywordTok{in}\NormalTok{ grid\_search.best\_params\_.items():}
    \BuiltInTok{print}\NormalTok{(}\SpecialStringTok{f"  }\SpecialCharTok{\{}\NormalTok{param}\SpecialCharTok{\}}\SpecialStringTok{: }\SpecialCharTok{\{}\NormalTok{value}\SpecialCharTok{\}}\SpecialStringTok{"}\NormalTok{)}

\BuiltInTok{print}\NormalTok{(}\SpecialStringTok{f"}\CharTok{\textbackslash{}n}\SpecialStringTok{Mejor score (ROC{-}AUC): }\SpecialCharTok{\{}\NormalTok{grid\_search}\SpecialCharTok{.}\NormalTok{best\_score\_}\SpecialCharTok{:.4f\}}\SpecialStringTok{"}\NormalTok{)}

\CommentTok{\# Evaluar mejor modelo}
\NormalTok{best\_model }\OperatorTok{=}\NormalTok{ grid\_search.best\_estimator\_}
\NormalTok{y\_pred }\OperatorTok{=}\NormalTok{ best\_model.predict(X)}
\NormalTok{y\_proba }\OperatorTok{=}\NormalTok{ best\_model.predict\_proba(X)[:, }\DecValTok{1}\NormalTok{]}

\BuiltInTok{print}\NormalTok{(}\SpecialStringTok{f"}\CharTok{\textbackslash{}n}\SpecialStringTok{Performance del mejor modelo:"}\NormalTok{)}
\BuiltInTok{print}\NormalTok{(}\SpecialStringTok{f"  Accuracy: }\SpecialCharTok{\{}\NormalTok{accuracy\_score(y, y\_pred)}\SpecialCharTok{:.4f\}}\SpecialStringTok{"}\NormalTok{)}
\BuiltInTok{print}\NormalTok{(}\SpecialStringTok{f"  Precision: }\SpecialCharTok{\{}\NormalTok{precision\_score(y, y\_pred)}\SpecialCharTok{:.4f\}}\SpecialStringTok{"}\NormalTok{)}
\BuiltInTok{print}\NormalTok{(}\SpecialStringTok{f"  Recall: }\SpecialCharTok{\{}\NormalTok{recall\_score(y, y\_pred)}\SpecialCharTok{:.4f\}}\SpecialStringTok{"}\NormalTok{)}
\BuiltInTok{print}\NormalTok{(}\SpecialStringTok{f"  F1{-}Score: }\SpecialCharTok{\{}\NormalTok{f1\_score(y, y\_pred)}\SpecialCharTok{:.4f\}}\SpecialStringTok{"}\NormalTok{)}
\BuiltInTok{print}\NormalTok{(}\SpecialStringTok{f"  ROC{-}AUC: }\SpecialCharTok{\{}\NormalTok{roc\_auc\_score(y, y\_proba)}\SpecialCharTok{:.4f\}}\SpecialStringTok{"}\NormalTok{)}
\end{Highlighting}
\end{Shaded}

\begin{verbatim}

Ejecutando Grid Search...
Fitting 5 folds for each of 108 candidates, totalling 540 fits


RESULTADOS GRID SEARCH
======================================================================

Mejores hiperparámetros:
  max_depth: 3
  min_samples_leaf: 1
  min_samples_split: 10
  n_estimators: 200

Mejor score (ROC-AUC): 0.8810

Performance del mejor modelo:
  Accuracy: 0.8400
  Precision: 0.7500
  Recall: 0.7800
  F1-Score: 0.7647
  ROC-AUC: 0.9141
\end{verbatim}

\subsection{Análisis de resultados Grid
Search}\label{anuxe1lisis-de-resultados-grid-search}

\begin{Shaded}
\begin{Highlighting}[]
\CommentTok{\# Convertir resultados a DataFrame}
\NormalTok{results }\OperatorTok{=}\NormalTok{ pd.DataFrame(grid\_search.cv\_results\_)}

\CommentTok{\# Top 10 configuraciones}
\NormalTok{top\_10 }\OperatorTok{=}\NormalTok{ results.nsmallest(}\DecValTok{10}\NormalTok{, }\StringTok{\textquotesingle{}rank\_test\_score\textquotesingle{}}\NormalTok{)[}
\NormalTok{    [}\StringTok{\textquotesingle{}params\textquotesingle{}}\NormalTok{, }\StringTok{\textquotesingle{}mean\_test\_score\textquotesingle{}}\NormalTok{, }\StringTok{\textquotesingle{}std\_test\_score\textquotesingle{}}\NormalTok{, }
     \StringTok{\textquotesingle{}mean\_train\_score\textquotesingle{}}\NormalTok{, }\StringTok{\textquotesingle{}mean\_fit\_time\textquotesingle{}}\NormalTok{]}
\NormalTok{]}

\BuiltInTok{print}\NormalTok{(}\StringTok{"}\CharTok{\textbackslash{}n\textbackslash{}n}\StringTok{TOP 10 CONFIGURACIONES"}\NormalTok{)}
\BuiltInTok{print}\NormalTok{(}\StringTok{"="} \OperatorTok{*} \DecValTok{70}\NormalTok{)}
\ControlFlowTok{for}\NormalTok{ idx, row }\KeywordTok{in}\NormalTok{ top\_10.iterrows():}
    \BuiltInTok{print}\NormalTok{(}\SpecialStringTok{f"}\CharTok{\textbackslash{}n}\SpecialStringTok{Rank }\SpecialCharTok{\{}\BuiltInTok{int}\NormalTok{(results.loc[idx, }\StringTok{\textquotesingle{}rank\_test\_score\textquotesingle{}}\NormalTok{])}\SpecialCharTok{\}}\SpecialStringTok{:"}\NormalTok{)}
    \BuiltInTok{print}\NormalTok{(}\SpecialStringTok{f"  Parámetros: }\SpecialCharTok{\{}\NormalTok{row[}\StringTok{\textquotesingle{}params\textquotesingle{}}\NormalTok{]}\SpecialCharTok{\}}\SpecialStringTok{"}\NormalTok{)}
    \BuiltInTok{print}\NormalTok{(}\SpecialStringTok{f"  Test Score: }\SpecialCharTok{\{}\NormalTok{row[}\StringTok{\textquotesingle{}mean\_test\_score\textquotesingle{}}\NormalTok{]}\SpecialCharTok{:.4f\}}\SpecialStringTok{ ± }\SpecialCharTok{\{}\NormalTok{row[}\StringTok{\textquotesingle{}std\_test\_score\textquotesingle{}}\NormalTok{]}\SpecialCharTok{:.4f\}}\SpecialStringTok{"}\NormalTok{)}
    \BuiltInTok{print}\NormalTok{(}\SpecialStringTok{f"  Train Score: }\SpecialCharTok{\{}\NormalTok{row[}\StringTok{\textquotesingle{}mean\_train\_score\textquotesingle{}}\NormalTok{]}\SpecialCharTok{:.4f\}}\SpecialStringTok{"}\NormalTok{)}
    \BuiltInTok{print}\NormalTok{(}\SpecialStringTok{f"  Tiempo: }\SpecialCharTok{\{}\NormalTok{row[}\StringTok{\textquotesingle{}mean\_fit\_time\textquotesingle{}}\NormalTok{]}\SpecialCharTok{:.2f\}}\SpecialStringTok{s"}\NormalTok{)}
\end{Highlighting}
\end{Shaded}

\begin{verbatim}


TOP 10 CONFIGURACIONES
======================================================================

Rank 1:
  Parámetros: {'max_depth': 3, 'min_samples_leaf': 1, 'min_samples_split': 10, 'n_estimators': 200}
  Test Score: 0.8810 ± 0.0557
  Train Score: 0.9245
  Tiempo: 1.14s

Rank 2:
  Parámetros: {'max_depth': 3, 'min_samples_leaf': 1, 'min_samples_split': 5, 'n_estimators': 200}
  Test Score: 0.8785 ± 0.0569
  Train Score: 0.9262
  Tiempo: 1.11s

Rank 3:
  Parámetros: {'max_depth': 3, 'min_samples_leaf': 4, 'min_samples_split': 2, 'n_estimators': 200}
  Test Score: 0.8780 ± 0.0645
  Train Score: 0.9183
  Tiempo: 1.15s

Rank 3:
  Parámetros: {'max_depth': 3, 'min_samples_leaf': 4, 'min_samples_split': 5, 'n_estimators': 200}
  Test Score: 0.8780 ± 0.0645
  Train Score: 0.9183
  Tiempo: 1.17s

Rank 5:
  Parámetros: {'max_depth': 3, 'min_samples_leaf': 4, 'min_samples_split': 10, 'n_estimators': 200}
  Test Score: 0.8770 ± 0.0615
  Train Score: 0.9170
  Tiempo: 1.20s

Rank 6:
  Parámetros: {'max_depth': 3, 'min_samples_leaf': 2, 'min_samples_split': 10, 'n_estimators': 200}
  Test Score: 0.8770 ± 0.0595
  Train Score: 0.9228
  Tiempo: 1.07s

Rank 7:
  Parámetros: {'max_depth': 3, 'min_samples_leaf': 1, 'min_samples_split': 2, 'n_estimators': 200}
  Test Score: 0.8765 ± 0.0567
  Train Score: 0.9275
  Tiempo: 1.09s

Rank 8:
  Parámetros: {'max_depth': 3, 'min_samples_leaf': 2, 'min_samples_split': 2, 'n_estimators': 200}
  Test Score: 0.8755 ± 0.0565
  Train Score: 0.9244
  Tiempo: 1.19s

Rank 9:
  Parámetros: {'max_depth': 3, 'min_samples_leaf': 1, 'min_samples_split': 5, 'n_estimators': 100}
  Test Score: 0.8750 ± 0.0527
  Train Score: 0.9266
  Tiempo: 0.53s

Rank 9:
  Parámetros: {'max_depth': 3, 'min_samples_leaf': 1, 'min_samples_split': 10, 'n_estimators': 50}
  Test Score: 0.8750 ± 0.0576
  Train Score: 0.9241
  Tiempo: 0.28s
\end{verbatim}

\subsection{Visualización de Grid
Search}\label{visualizaciuxf3n-de-grid-search}

\begin{Shaded}
\begin{Highlighting}[]
\CommentTok{\# Analizar efecto de hiperparámetros clave}
\NormalTok{fig, axes }\OperatorTok{=}\NormalTok{ plt.subplots(}\DecValTok{2}\NormalTok{, }\DecValTok{2}\NormalTok{, figsize}\OperatorTok{=}\NormalTok{(}\DecValTok{14}\NormalTok{, }\DecValTok{10}\NormalTok{))}

\CommentTok{\# 1. n\_estimators vs score}
\NormalTok{param }\OperatorTok{=} \StringTok{\textquotesingle{}param\_n\_estimators\textquotesingle{}}
\NormalTok{grouped }\OperatorTok{=}\NormalTok{ results.groupby(param).agg(\{}
    \StringTok{\textquotesingle{}mean\_test\_score\textquotesingle{}}\NormalTok{: }\StringTok{\textquotesingle{}mean\textquotesingle{}}\NormalTok{,}
    \StringTok{\textquotesingle{}std\_test\_score\textquotesingle{}}\NormalTok{: }\StringTok{\textquotesingle{}mean\textquotesingle{}}
\NormalTok{\}).reset\_index()}

\NormalTok{axes[}\DecValTok{0}\NormalTok{, }\DecValTok{0}\NormalTok{].errorbar(grouped[param], grouped[}\StringTok{\textquotesingle{}mean\_test\_score\textquotesingle{}}\NormalTok{],}
\NormalTok{                    yerr}\OperatorTok{=}\NormalTok{grouped[}\StringTok{\textquotesingle{}std\_test\_score\textquotesingle{}}\NormalTok{], marker}\OperatorTok{=}\StringTok{\textquotesingle{}o\textquotesingle{}}\NormalTok{, capsize}\OperatorTok{=}\DecValTok{5}\NormalTok{)}
\NormalTok{axes[}\DecValTok{0}\NormalTok{, }\DecValTok{0}\NormalTok{].set\_xlabel(}\StringTok{\textquotesingle{}n\_estimators\textquotesingle{}}\NormalTok{)}
\NormalTok{axes[}\DecValTok{0}\NormalTok{, }\DecValTok{0}\NormalTok{].set\_ylabel(}\StringTok{\textquotesingle{}ROC{-}AUC\textquotesingle{}}\NormalTok{)}
\NormalTok{axes[}\DecValTok{0}\NormalTok{, }\DecValTok{0}\NormalTok{].set\_title(}\StringTok{\textquotesingle{}Efecto de n\_estimators\textquotesingle{}}\NormalTok{)}
\NormalTok{axes[}\DecValTok{0}\NormalTok{, }\DecValTok{0}\NormalTok{].grid(}\VariableTok{True}\NormalTok{, alpha}\OperatorTok{=}\FloatTok{0.3}\NormalTok{)}

\CommentTok{\# 2. max\_depth vs score}
\NormalTok{param }\OperatorTok{=} \StringTok{\textquotesingle{}param\_max\_depth\textquotesingle{}}
\NormalTok{grouped }\OperatorTok{=}\NormalTok{ results.groupby(param).agg(\{}
    \StringTok{\textquotesingle{}mean\_test\_score\textquotesingle{}}\NormalTok{: }\StringTok{\textquotesingle{}mean\textquotesingle{}}\NormalTok{,}
    \StringTok{\textquotesingle{}std\_test\_score\textquotesingle{}}\NormalTok{: }\StringTok{\textquotesingle{}mean\textquotesingle{}}
\NormalTok{\}).reset\_index()}

\CommentTok{\# Reemplazar None con string para graficar}
\NormalTok{grouped[param] }\OperatorTok{=}\NormalTok{ grouped[param].astype(}\BuiltInTok{str}\NormalTok{)}

\NormalTok{axes[}\DecValTok{0}\NormalTok{, }\DecValTok{1}\NormalTok{].errorbar(}\BuiltInTok{range}\NormalTok{(}\BuiltInTok{len}\NormalTok{(grouped)), grouped[}\StringTok{\textquotesingle{}mean\_test\_score\textquotesingle{}}\NormalTok{],}
\NormalTok{                    yerr}\OperatorTok{=}\NormalTok{grouped[}\StringTok{\textquotesingle{}std\_test\_score\textquotesingle{}}\NormalTok{], marker}\OperatorTok{=}\StringTok{\textquotesingle{}s\textquotesingle{}}\NormalTok{, capsize}\OperatorTok{=}\DecValTok{5}\NormalTok{)}
\NormalTok{axes[}\DecValTok{0}\NormalTok{, }\DecValTok{1}\NormalTok{].set\_xticks(}\BuiltInTok{range}\NormalTok{(}\BuiltInTok{len}\NormalTok{(grouped)))}
\NormalTok{axes[}\DecValTok{0}\NormalTok{, }\DecValTok{1}\NormalTok{].set\_xticklabels(grouped[param])}
\NormalTok{axes[}\DecValTok{0}\NormalTok{, }\DecValTok{1}\NormalTok{].set\_xlabel(}\StringTok{\textquotesingle{}max\_depth\textquotesingle{}}\NormalTok{)}
\NormalTok{axes[}\DecValTok{0}\NormalTok{, }\DecValTok{1}\NormalTok{].set\_ylabel(}\StringTok{\textquotesingle{}ROC{-}AUC\textquotesingle{}}\NormalTok{)}
\NormalTok{axes[}\DecValTok{0}\NormalTok{, }\DecValTok{1}\NormalTok{].set\_title(}\StringTok{\textquotesingle{}Efecto de max\_depth\textquotesingle{}}\NormalTok{)}
\NormalTok{axes[}\DecValTok{0}\NormalTok{, }\DecValTok{1}\NormalTok{].grid(}\VariableTok{True}\NormalTok{, alpha}\OperatorTok{=}\FloatTok{0.3}\NormalTok{)}

\CommentTok{\# 3. Heatmap: n\_estimators vs max\_depth}
\NormalTok{pivot\_data }\OperatorTok{=}\NormalTok{ results.pivot\_table(}
\NormalTok{    values}\OperatorTok{=}\StringTok{\textquotesingle{}mean\_test\_score\textquotesingle{}}\NormalTok{,}
\NormalTok{    index}\OperatorTok{=}\StringTok{\textquotesingle{}param\_max\_depth\textquotesingle{}}\NormalTok{,}
\NormalTok{    columns}\OperatorTok{=}\StringTok{\textquotesingle{}param\_n\_estimators\textquotesingle{}}\NormalTok{,}
\NormalTok{    aggfunc}\OperatorTok{=}\StringTok{\textquotesingle{}mean\textquotesingle{}}
\NormalTok{)}

\NormalTok{sns.heatmap(pivot\_data, annot}\OperatorTok{=}\VariableTok{True}\NormalTok{, fmt}\OperatorTok{=}\StringTok{\textquotesingle{}.3f\textquotesingle{}}\NormalTok{, cmap}\OperatorTok{=}\StringTok{\textquotesingle{}YlOrRd\textquotesingle{}}\NormalTok{, }
\NormalTok{            ax}\OperatorTok{=}\NormalTok{axes[}\DecValTok{1}\NormalTok{, }\DecValTok{0}\NormalTok{], cbar\_kws}\OperatorTok{=}\NormalTok{\{}\StringTok{\textquotesingle{}label\textquotesingle{}}\NormalTok{: }\StringTok{\textquotesingle{}ROC{-}AUC\textquotesingle{}}\NormalTok{\})}
\NormalTok{axes[}\DecValTok{1}\NormalTok{, }\DecValTok{0}\NormalTok{].set\_title(}\StringTok{\textquotesingle{}Heatmap: n\_estimators vs max\_depth\textquotesingle{}}\NormalTok{)}

\CommentTok{\# 4. Train vs Test scores (top 20)}
\NormalTok{top\_20 }\OperatorTok{=}\NormalTok{ results.nsmallest(}\DecValTok{20}\NormalTok{, }\StringTok{\textquotesingle{}rank\_test\_score\textquotesingle{}}\NormalTok{)}

\NormalTok{axes[}\DecValTok{1}\NormalTok{, }\DecValTok{1}\NormalTok{].scatter(top\_20[}\StringTok{\textquotesingle{}mean\_train\_score\textquotesingle{}}\NormalTok{], top\_20[}\StringTok{\textquotesingle{}mean\_test\_score\textquotesingle{}}\NormalTok{],}
\NormalTok{                  alpha}\OperatorTok{=}\FloatTok{0.6}\NormalTok{, s}\OperatorTok{=}\DecValTok{100}\NormalTok{)}
\NormalTok{axes[}\DecValTok{1}\NormalTok{, }\DecValTok{1}\NormalTok{].plot([}\FloatTok{0.8}\NormalTok{, }\FloatTok{1.0}\NormalTok{], [}\FloatTok{0.8}\NormalTok{, }\FloatTok{1.0}\NormalTok{], }\StringTok{\textquotesingle{}r{-}{-}\textquotesingle{}}\NormalTok{, label}\OperatorTok{=}\StringTok{\textquotesingle{}Train=Test\textquotesingle{}}\NormalTok{)}
\NormalTok{axes[}\DecValTok{1}\NormalTok{, }\DecValTok{1}\NormalTok{].set\_xlabel(}\StringTok{\textquotesingle{}Train Score\textquotesingle{}}\NormalTok{)}
\NormalTok{axes[}\DecValTok{1}\NormalTok{, }\DecValTok{1}\NormalTok{].set\_ylabel(}\StringTok{\textquotesingle{}Test Score\textquotesingle{}}\NormalTok{)}
\NormalTok{axes[}\DecValTok{1}\NormalTok{, }\DecValTok{1}\NormalTok{].set\_title(}\StringTok{\textquotesingle{}Train vs Test Scores (Top 20)\textquotesingle{}}\NormalTok{)}
\NormalTok{axes[}\DecValTok{1}\NormalTok{, }\DecValTok{1}\NormalTok{].legend()}
\NormalTok{axes[}\DecValTok{1}\NormalTok{, }\DecValTok{1}\NormalTok{].grid(}\VariableTok{True}\NormalTok{, alpha}\OperatorTok{=}\FloatTok{0.3}\NormalTok{)}

\NormalTok{plt.tight\_layout()}
\NormalTok{plt.savefig(}\StringTok{\textquotesingle{}grid\_search\_analysis.png\textquotesingle{}}\NormalTok{, dpi}\OperatorTok{=}\DecValTok{300}\NormalTok{, bbox\_inches}\OperatorTok{=}\StringTok{\textquotesingle{}tight\textquotesingle{}}\NormalTok{)}
\BuiltInTok{print}\NormalTok{(}\StringTok{"}\CharTok{\textbackslash{}n}\StringTok{Gráfico guardado como \textquotesingle{}grid\_search\_analysis.png\textquotesingle{}"}\NormalTok{)}
\end{Highlighting}
\end{Shaded}

\begin{verbatim}

Gráfico guardado como 'grid_search_analysis.png'
\end{verbatim}

\pandocbounded{\includegraphics[keepaspectratio]{index_files/figure-pdf/cell-12-output-2.pdf}}

\subsection{Random Search
implementation}\label{random-search-implementation}

\begin{Shaded}
\begin{Highlighting}[]
\ImportTok{from}\NormalTok{ sklearn.model\_selection }\ImportTok{import}\NormalTok{ RandomizedSearchCV}
\ImportTok{from}\NormalTok{ scipy.stats }\ImportTok{import}\NormalTok{ randint, uniform}

\BuiltInTok{print}\NormalTok{(}\StringTok{"}\CharTok{\textbackslash{}n\textbackslash{}n}\StringTok{RANDOM SEARCH"}\NormalTok{)}
\BuiltInTok{print}\NormalTok{(}\StringTok{"="} \OperatorTok{*} \DecValTok{70}\NormalTok{)}

\CommentTok{\# Definir distribuciones para Random Search}
\NormalTok{param\_distributions }\OperatorTok{=}\NormalTok{ \{}
    \StringTok{\textquotesingle{}n\_estimators\textquotesingle{}}\NormalTok{: randint(}\DecValTok{50}\NormalTok{, }\DecValTok{300}\NormalTok{),}
    \StringTok{\textquotesingle{}max\_depth\textquotesingle{}}\NormalTok{: [}\DecValTok{3}\NormalTok{, }\DecValTok{5}\NormalTok{, }\DecValTok{7}\NormalTok{, }\DecValTok{10}\NormalTok{, }\DecValTok{15}\NormalTok{, }\VariableTok{None}\NormalTok{],}
    \StringTok{\textquotesingle{}min\_samples\_split\textquotesingle{}}\NormalTok{: randint(}\DecValTok{2}\NormalTok{, }\DecValTok{20}\NormalTok{),}
    \StringTok{\textquotesingle{}min\_samples\_leaf\textquotesingle{}}\NormalTok{: randint(}\DecValTok{1}\NormalTok{, }\DecValTok{10}\NormalTok{),}
    \StringTok{\textquotesingle{}max\_features\textquotesingle{}}\NormalTok{: [}\StringTok{\textquotesingle{}sqrt\textquotesingle{}}\NormalTok{, }\StringTok{\textquotesingle{}log2\textquotesingle{}}\NormalTok{, }\VariableTok{None}\NormalTok{]}
\NormalTok{\}}

\CommentTok{\# Ejecutar Random Search}
\NormalTok{random\_search }\OperatorTok{=}\NormalTok{ RandomizedSearchCV(}
\NormalTok{    estimator}\OperatorTok{=}\NormalTok{RandomForestClassifier(random\_state}\OperatorTok{=}\DecValTok{42}\NormalTok{),}
\NormalTok{    param\_distributions}\OperatorTok{=}\NormalTok{param\_distributions,}
\NormalTok{    n\_iter}\OperatorTok{=}\DecValTok{100}\NormalTok{,  }\CommentTok{\# Número de combinaciones a probar}
\NormalTok{    cv}\OperatorTok{=}\NormalTok{StratifiedKFold(n\_splits}\OperatorTok{=}\DecValTok{5}\NormalTok{, shuffle}\OperatorTok{=}\VariableTok{True}\NormalTok{, random\_state}\OperatorTok{=}\DecValTok{42}\NormalTok{),}
\NormalTok{    scoring}\OperatorTok{=}\StringTok{\textquotesingle{}roc\_auc\textquotesingle{}}\NormalTok{,}
\NormalTok{    n\_jobs}\OperatorTok{={-}}\DecValTok{1}\NormalTok{,}
\NormalTok{    verbose}\OperatorTok{=}\DecValTok{1}\NormalTok{,}
\NormalTok{    random\_state}\OperatorTok{=}\DecValTok{42}\NormalTok{,}
\NormalTok{    return\_train\_score}\OperatorTok{=}\VariableTok{True}
\NormalTok{)}

\BuiltInTok{print}\NormalTok{(}\StringTok{"Ejecutando Random Search (100 iteraciones)..."}\NormalTok{)}
\NormalTok{random\_search.fit(X, y)}

\BuiltInTok{print}\NormalTok{(}\SpecialStringTok{f"}\CharTok{\textbackslash{}n}\SpecialStringTok{Mejores hiperparámetros:"}\NormalTok{)}
\ControlFlowTok{for}\NormalTok{ param, value }\KeywordTok{in}\NormalTok{ random\_search.best\_params\_.items():}
    \BuiltInTok{print}\NormalTok{(}\SpecialStringTok{f"  }\SpecialCharTok{\{}\NormalTok{param}\SpecialCharTok{\}}\SpecialStringTok{: }\SpecialCharTok{\{}\NormalTok{value}\SpecialCharTok{\}}\SpecialStringTok{"}\NormalTok{)}

\BuiltInTok{print}\NormalTok{(}\SpecialStringTok{f"}\CharTok{\textbackslash{}n}\SpecialStringTok{Mejor score: }\SpecialCharTok{\{}\NormalTok{random\_search}\SpecialCharTok{.}\NormalTok{best\_score\_}\SpecialCharTok{:.4f\}}\SpecialStringTok{"}\NormalTok{)}

\CommentTok{\# Comparar con Grid Search}
\BuiltInTok{print}\NormalTok{(}\SpecialStringTok{f"}\CharTok{\textbackslash{}n}\SpecialStringTok{Comparación:"}\NormalTok{)}
\BuiltInTok{print}\NormalTok{(}\SpecialStringTok{f"  Grid Search score:   }\SpecialCharTok{\{}\NormalTok{grid\_search}\SpecialCharTok{.}\NormalTok{best\_score\_}\SpecialCharTok{:.4f\}}\SpecialStringTok{"}\NormalTok{)}
\BuiltInTok{print}\NormalTok{(}\SpecialStringTok{f"  Random Search score: }\SpecialCharTok{\{}\NormalTok{random\_search}\SpecialCharTok{.}\NormalTok{best\_score\_}\SpecialCharTok{:.4f\}}\SpecialStringTok{"}\NormalTok{)}
\BuiltInTok{print}\NormalTok{(}\SpecialStringTok{f"  Diferencia:          }\SpecialCharTok{\{}\NormalTok{random\_search}\SpecialCharTok{.}\NormalTok{best\_score\_ }\OperatorTok{{-}}\NormalTok{ grid\_search}\SpecialCharTok{.}\NormalTok{best\_score\_}\SpecialCharTok{:.4f\}}\SpecialStringTok{"}\NormalTok{)}
\end{Highlighting}
\end{Shaded}

\begin{verbatim}


RANDOM SEARCH
======================================================================
Ejecutando Random Search (100 iteraciones)...
Fitting 5 folds for each of 100 candidates, totalling 500 fits

Mejores hiperparámetros:
  max_depth: 3
  max_features: sqrt
  min_samples_leaf: 6
  min_samples_split: 17
  n_estimators: 142

Mejor score: 0.8830

Comparación:
  Grid Search score:   0.8810
  Random Search score: 0.8830
  Diferencia:          0.0020
\end{verbatim}

\section{Curvas de aprendizaje}\label{curvas-de-aprendizaje}

\subsection{Fundamento matemático}\label{fundamento-matemuxe1tico-2}

Las curvas de aprendizaje muestran cómo el performance del modelo cambia
con el tamaño del conjunto de entrenamiento.

\subsubsection{Definición}\label{definiciuxf3n}

Para tamaños de muestra \(m \in \{m_1, m_2, \ldots, m_k\}\) donde
\(m_k = n\):

\textbf{Error de entrenamiento}: \[
\text{Error}_{\text{train}}(m) = \frac{1}{m} \sum_{i=1}^{m} L(y_i, \hat{f}_m(\mathbf{X}_i))
\]

\textbf{Error de validación}: \[
\text{Error}_{\text{val}}(m) = \frac{1}{n-m} \sum_{i=m+1}^{n} L(y_i, \hat{f}_m(\mathbf{X}_i))
\]

\subsubsection{Comportamiento típico}\label{comportamiento-tuxedpico}

\textbf{High Bias (Underfitting)}:

\begin{itemize}
\tightlist
\item
  Error de train alto, se mantiene alto
\item
  Error de validación alto, converge cerca de train
\item
  Gap pequeño
\item
  \textbf{Solución}: Modelo más complejo, más features
\end{itemize}

\textbf{High Variance (Overfitting)}:

\begin{itemize}
\tightlist
\item
  Error de train bajo
\item
  Error de validación alto, no converge
\item
  Gap grande
\item
  \textbf{Solución}: Más datos, regularización, modelo más simple
\end{itemize}

\textbf{Good Fit}:

\begin{itemize}
\tightlist
\item
  Errores convergen a valor razonable
\item
  Gap pequeño
\end{itemize}

\subsection{Implementación}\label{implementaciuxf3n}

\begin{Shaded}
\begin{Highlighting}[]
\BuiltInTok{print}\NormalTok{(}\StringTok{"}\CharTok{\textbackslash{}n\textbackslash{}n}\StringTok{CURVAS DE APRENDIZAJE"}\NormalTok{)}
\BuiltInTok{print}\NormalTok{(}\StringTok{"="} \OperatorTok{*} \DecValTok{70}\NormalTok{)}

\CommentTok{\# Función para crear curvas de aprendizaje}
\KeywordTok{def}\NormalTok{ plot\_learning\_curves(modelo, X, y, title, filename):}
    \CommentTok{"""}
\CommentTok{    Crea y visualiza curvas de aprendizaje}
\CommentTok{    """}
    \CommentTok{\# Calcular curvas}
\NormalTok{    train\_sizes, train\_scores, val\_scores }\OperatorTok{=}\NormalTok{ learning\_curve(}
\NormalTok{        modelo, X, y,}
\NormalTok{        cv}\OperatorTok{=}\DecValTok{5}\NormalTok{,}
\NormalTok{        n\_jobs}\OperatorTok{={-}}\DecValTok{1}\NormalTok{,}
\NormalTok{        train\_sizes}\OperatorTok{=}\NormalTok{np.linspace(}\FloatTok{0.1}\NormalTok{, }\FloatTok{1.0}\NormalTok{, }\DecValTok{10}\NormalTok{),}
\NormalTok{        scoring}\OperatorTok{=}\StringTok{\textquotesingle{}accuracy\textquotesingle{}}\NormalTok{,}
\NormalTok{        random\_state}\OperatorTok{=}\DecValTok{42}
\NormalTok{    )}
    
    \CommentTok{\# Calcular medias y desviaciones estándar}
\NormalTok{    train\_mean }\OperatorTok{=}\NormalTok{ train\_scores.mean(axis}\OperatorTok{=}\DecValTok{1}\NormalTok{)}
\NormalTok{    train\_std }\OperatorTok{=}\NormalTok{ train\_scores.std(axis}\OperatorTok{=}\DecValTok{1}\NormalTok{)}
\NormalTok{    val\_mean }\OperatorTok{=}\NormalTok{ val\_scores.mean(axis}\OperatorTok{=}\DecValTok{1}\NormalTok{)}
\NormalTok{    val\_std }\OperatorTok{=}\NormalTok{ val\_scores.std(axis}\OperatorTok{=}\DecValTok{1}\NormalTok{)}
    
    \CommentTok{\# Graficar}
\NormalTok{    plt.figure(figsize}\OperatorTok{=}\NormalTok{(}\DecValTok{10}\NormalTok{, }\DecValTok{6}\NormalTok{))}
    
\NormalTok{    plt.plot(train\_sizes, train\_mean, }\StringTok{\textquotesingle{}o{-}\textquotesingle{}}\NormalTok{, color}\OperatorTok{=}\StringTok{\textquotesingle{}blue\textquotesingle{}}\NormalTok{, }
\NormalTok{             label}\OperatorTok{=}\StringTok{\textquotesingle{}Training score\textquotesingle{}}\NormalTok{, linewidth}\OperatorTok{=}\DecValTok{2}\NormalTok{)}
\NormalTok{    plt.fill\_between(train\_sizes, train\_mean }\OperatorTok{{-}}\NormalTok{ train\_std,}
\NormalTok{                     train\_mean }\OperatorTok{+}\NormalTok{ train\_std, alpha}\OperatorTok{=}\FloatTok{0.1}\NormalTok{, color}\OperatorTok{=}\StringTok{\textquotesingle{}blue\textquotesingle{}}\NormalTok{)}
    
\NormalTok{    plt.plot(train\_sizes, val\_mean, }\StringTok{\textquotesingle{}s{-}\textquotesingle{}}\NormalTok{, color}\OperatorTok{=}\StringTok{\textquotesingle{}orange\textquotesingle{}}\NormalTok{,}
\NormalTok{             label}\OperatorTok{=}\StringTok{\textquotesingle{}Cross{-}validation score\textquotesingle{}}\NormalTok{, linewidth}\OperatorTok{=}\DecValTok{2}\NormalTok{)}
\NormalTok{    plt.fill\_between(train\_sizes, val\_mean }\OperatorTok{{-}}\NormalTok{ val\_std,}
\NormalTok{                     val\_mean }\OperatorTok{+}\NormalTok{ val\_std, alpha}\OperatorTok{=}\FloatTok{0.1}\NormalTok{, color}\OperatorTok{=}\StringTok{\textquotesingle{}orange\textquotesingle{}}\NormalTok{)}
    
\NormalTok{    plt.xlabel(}\StringTok{\textquotesingle{}Training Size\textquotesingle{}}\NormalTok{)}
\NormalTok{    plt.ylabel(}\StringTok{\textquotesingle{}Accuracy\textquotesingle{}}\NormalTok{)}
\NormalTok{    plt.title(title)}
\NormalTok{    plt.legend(loc}\OperatorTok{=}\StringTok{\textquotesingle{}best\textquotesingle{}}\NormalTok{)}
\NormalTok{    plt.grid(}\VariableTok{True}\NormalTok{, alpha}\OperatorTok{=}\FloatTok{0.3}\NormalTok{)}
\NormalTok{    plt.ylim([}\FloatTok{0.5}\NormalTok{, }\FloatTok{1.05}\NormalTok{])}
    
\NormalTok{    plt.tight\_layout()}
\NormalTok{    plt.savefig(filename, dpi}\OperatorTok{=}\DecValTok{300}\NormalTok{, bbox\_inches}\OperatorTok{=}\StringTok{\textquotesingle{}tight\textquotesingle{}}\NormalTok{)}
    \BuiltInTok{print}\NormalTok{(}\SpecialStringTok{f"}\CharTok{\textbackslash{}n}\SpecialStringTok{Gráfico guardado como \textquotesingle{}}\SpecialCharTok{\{}\NormalTok{filename}\SpecialCharTok{\}}\SpecialStringTok{\textquotesingle{}"}\NormalTok{)}
    
    \CommentTok{\# Diagnóstico}
\NormalTok{    final\_gap }\OperatorTok{=}\NormalTok{ train\_mean[}\OperatorTok{{-}}\DecValTok{1}\NormalTok{] }\OperatorTok{{-}}\NormalTok{ val\_mean[}\OperatorTok{{-}}\DecValTok{1}\NormalTok{]}
    \BuiltInTok{print}\NormalTok{(}\SpecialStringTok{f"}\CharTok{\textbackslash{}n}\SpecialStringTok{Diagnóstico para }\SpecialCharTok{\{}\NormalTok{title}\SpecialCharTok{\}}\SpecialStringTok{:"}\NormalTok{)}
    \BuiltInTok{print}\NormalTok{(}\SpecialStringTok{f"  Final training score: }\SpecialCharTok{\{}\NormalTok{train\_mean[}\OperatorTok{{-}}\DecValTok{1}\NormalTok{]}\SpecialCharTok{:.4f\}}\SpecialStringTok{"}\NormalTok{)}
    \BuiltInTok{print}\NormalTok{(}\SpecialStringTok{f"  Final validation score: }\SpecialCharTok{\{}\NormalTok{val\_mean[}\OperatorTok{{-}}\DecValTok{1}\NormalTok{]}\SpecialCharTok{:.4f\}}\SpecialStringTok{"}\NormalTok{)}
    \BuiltInTok{print}\NormalTok{(}\SpecialStringTok{f"  Gap: }\SpecialCharTok{\{}\NormalTok{final\_gap}\SpecialCharTok{:.4f\}}\SpecialStringTok{"}\NormalTok{)}
    
    \ControlFlowTok{if}\NormalTok{ final\_gap }\OperatorTok{\textgreater{}} \FloatTok{0.15}\NormalTok{:}
        \BuiltInTok{print}\NormalTok{(}\StringTok{"  ⚠️  High variance (overfitting)"}\NormalTok{)}
        \BuiltInTok{print}\NormalTok{(}\StringTok{"  Recomendaciones: Más datos, regularización, modelo más simple"}\NormalTok{)}
    \ControlFlowTok{elif}\NormalTok{ val\_mean[}\OperatorTok{{-}}\DecValTok{1}\NormalTok{] }\OperatorTok{\textless{}} \FloatTok{0.70}\NormalTok{:}
        \BuiltInTok{print}\NormalTok{(}\StringTok{"  ⚠️  High bias (underfitting)"}\NormalTok{)}
        \BuiltInTok{print}\NormalTok{(}\StringTok{"  Recomendaciones: Modelo más complejo, más features"}\NormalTok{)}
    \ControlFlowTok{else}\NormalTok{:}
        \BuiltInTok{print}\NormalTok{(}\StringTok{"  ✓ Balance adecuado"}\NormalTok{)}
    
    \ControlFlowTok{return}\NormalTok{ train\_sizes, train\_mean, val\_mean}

\CommentTok{\# Modelos a comparar}
\NormalTok{modelos }\OperatorTok{=}\NormalTok{ \{}
    \StringTok{\textquotesingle{}Logistic Regression\textquotesingle{}}\NormalTok{: LogisticRegression(random\_state}\OperatorTok{=}\DecValTok{42}\NormalTok{),}
    \StringTok{\textquotesingle{}Decision Tree (depth=3)\textquotesingle{}}\NormalTok{: DecisionTreeClassifier(max\_depth}\OperatorTok{=}\DecValTok{3}\NormalTok{, random\_state}\OperatorTok{=}\DecValTok{42}\NormalTok{),}
    \StringTok{\textquotesingle{}Decision Tree (depth=20)\textquotesingle{}}\NormalTok{: DecisionTreeClassifier(max\_depth}\OperatorTok{=}\DecValTok{20}\NormalTok{, random\_state}\OperatorTok{=}\DecValTok{42}\NormalTok{),}
    \StringTok{\textquotesingle{}Random Forest\textquotesingle{}}\NormalTok{: RandomForestClassifier(n\_estimators}\OperatorTok{=}\DecValTok{100}\NormalTok{, random\_state}\OperatorTok{=}\DecValTok{42}\NormalTok{)}
\NormalTok{\}}

\CommentTok{\# Crear curvas para cada modelo}
\ControlFlowTok{for}\NormalTok{ nombre, modelo }\KeywordTok{in}\NormalTok{ modelos.items():}
\NormalTok{    filename }\OperatorTok{=} \SpecialStringTok{f"learning\_curve\_}\SpecialCharTok{\{}\NormalTok{nombre}\SpecialCharTok{.}\NormalTok{replace(}\StringTok{\textquotesingle{} \textquotesingle{}}\NormalTok{, }\StringTok{\textquotesingle{}\_\textquotesingle{}}\NormalTok{)}\SpecialCharTok{.}\NormalTok{lower()}\SpecialCharTok{\}}\SpecialStringTok{.png"}
\NormalTok{    plot\_learning\_curves(modelo, X, y, }\SpecialStringTok{f\textquotesingle{}Learning Curve: }\SpecialCharTok{\{}\NormalTok{nombre}\SpecialCharTok{\}}\SpecialStringTok{\textquotesingle{}}\NormalTok{, filename)}
\end{Highlighting}
\end{Shaded}

\begin{verbatim}


CURVAS DE APRENDIZAJE
======================================================================
\end{verbatim}

\begin{verbatim}

Gráfico guardado como 'learning_curve_logistic_regression.png'

Diagnóstico para Learning Curve: Logistic Regression:
  Final training score: 0.8033
  Final validation score: 0.7867
  Gap: 0.0167
  ✓ Balance adecuado

Gráfico guardado como 'learning_curve_decision_tree_(depth=3).png'

Diagnóstico para Learning Curve: Decision Tree (depth=3):
  Final training score: 0.8283
  Final validation score: 0.7267
  Gap: 0.1017
  ✓ Balance adecuado

Gráfico guardado como 'learning_curve_decision_tree_(depth=20).png'

Diagnóstico para Learning Curve: Decision Tree (depth=20):
  Final training score: 0.9400
  Final validation score: 0.7000
  Gap: 0.2400
  ⚠️  High variance (overfitting)
  Recomendaciones: Más datos, regularización, modelo más simple

Gráfico guardado como 'learning_curve_random_forest.png'

Diagnóstico para Learning Curve: Random Forest:
  Final training score: 0.9367
  Final validation score: 0.7400
  Gap: 0.1967
  ⚠️  High variance (overfitting)
  Recomendaciones: Más datos, regularización, modelo más simple
\end{verbatim}

\pandocbounded{\includegraphics[keepaspectratio]{index_files/figure-pdf/cell-14-output-3.pdf}}

\pandocbounded{\includegraphics[keepaspectratio]{index_files/figure-pdf/cell-14-output-4.pdf}}

\pandocbounded{\includegraphics[keepaspectratio]{index_files/figure-pdf/cell-14-output-5.pdf}}

\pandocbounded{\includegraphics[keepaspectratio]{index_files/figure-pdf/cell-14-output-6.pdf}}

\subsection{Comparación de modelos}\label{comparaciuxf3n-de-modelos}

\begin{Shaded}
\begin{Highlighting}[]
\CommentTok{\# Comparar todas las curvas en un gráfico}
\NormalTok{fig, axes }\OperatorTok{=}\NormalTok{ plt.subplots(}\DecValTok{2}\NormalTok{, }\DecValTok{2}\NormalTok{, figsize}\OperatorTok{=}\NormalTok{(}\DecValTok{14}\NormalTok{, }\DecValTok{10}\NormalTok{))}
\NormalTok{axes }\OperatorTok{=}\NormalTok{ axes.ravel()}

\ControlFlowTok{for}\NormalTok{ idx, (nombre, modelo) }\KeywordTok{in} \BuiltInTok{enumerate}\NormalTok{(modelos.items()):}
    \CommentTok{\# Calcular curvas}
\NormalTok{    train\_sizes, train\_scores, val\_scores }\OperatorTok{=}\NormalTok{ learning\_curve(}
\NormalTok{        modelo, X, y,}
\NormalTok{        cv}\OperatorTok{=}\DecValTok{5}\NormalTok{,}
\NormalTok{        n\_jobs}\OperatorTok{={-}}\DecValTok{1}\NormalTok{,}
\NormalTok{        train\_sizes}\OperatorTok{=}\NormalTok{np.linspace(}\FloatTok{0.1}\NormalTok{, }\FloatTok{1.0}\NormalTok{, }\DecValTok{10}\NormalTok{),}
\NormalTok{        scoring}\OperatorTok{=}\StringTok{\textquotesingle{}accuracy\textquotesingle{}}\NormalTok{,}
\NormalTok{        random\_state}\OperatorTok{=}\DecValTok{42}
\NormalTok{    )}
    
\NormalTok{    train\_mean }\OperatorTok{=}\NormalTok{ train\_scores.mean(axis}\OperatorTok{=}\DecValTok{1}\NormalTok{)}
\NormalTok{    val\_mean }\OperatorTok{=}\NormalTok{ val\_scores.mean(axis}\OperatorTok{=}\DecValTok{1}\NormalTok{)}
    
\NormalTok{    axes[idx].plot(train\_sizes, train\_mean, }\StringTok{\textquotesingle{}o{-}\textquotesingle{}}\NormalTok{, label}\OperatorTok{=}\StringTok{\textquotesingle{}Train\textquotesingle{}}\NormalTok{, linewidth}\OperatorTok{=}\DecValTok{2}\NormalTok{)}
\NormalTok{    axes[idx].plot(train\_sizes, val\_mean, }\StringTok{\textquotesingle{}s{-}\textquotesingle{}}\NormalTok{, label}\OperatorTok{=}\StringTok{\textquotesingle{}Validation\textquotesingle{}}\NormalTok{, linewidth}\OperatorTok{=}\DecValTok{2}\NormalTok{)}
\NormalTok{    axes[idx].set\_xlabel(}\StringTok{\textquotesingle{}Training Size\textquotesingle{}}\NormalTok{)}
\NormalTok{    axes[idx].set\_ylabel(}\StringTok{\textquotesingle{}Accuracy\textquotesingle{}}\NormalTok{)}
\NormalTok{    axes[idx].set\_title(nombre)}
\NormalTok{    axes[idx].legend()}
\NormalTok{    axes[idx].grid(}\VariableTok{True}\NormalTok{, alpha}\OperatorTok{=}\FloatTok{0.3}\NormalTok{)}
\NormalTok{    axes[idx].set\_ylim([}\FloatTok{0.5}\NormalTok{, }\FloatTok{1.05}\NormalTok{])}

\NormalTok{plt.tight\_layout()}
\NormalTok{plt.savefig(}\StringTok{\textquotesingle{}learning\_curves\_comparison.png\textquotesingle{}}\NormalTok{, dpi}\OperatorTok{=}\DecValTok{300}\NormalTok{, bbox\_inches}\OperatorTok{=}\StringTok{\textquotesingle{}tight\textquotesingle{}}\NormalTok{)}
\BuiltInTok{print}\NormalTok{(}\StringTok{"}\CharTok{\textbackslash{}n}\StringTok{Comparación guardada como \textquotesingle{}learning\_curves\_comparison.png\textquotesingle{}"}\NormalTok{)}
\end{Highlighting}
\end{Shaded}

\begin{verbatim}

Comparación guardada como 'learning_curves_comparison.png'
\end{verbatim}

\pandocbounded{\includegraphics[keepaspectratio]{index_files/figure-pdf/cell-15-output-2.pdf}}

\section{Contribución de la data al
modelo}\label{contribuciuxf3n-de-la-data-al-modelo}

\subsection{Análisis de sensibilidad al tamaño de
muestra}\label{anuxe1lisis-de-sensibilidad-al-tamauxf1o-de-muestra}

\begin{Shaded}
\begin{Highlighting}[]
\BuiltInTok{print}\NormalTok{(}\StringTok{"}\CharTok{\textbackslash{}n\textbackslash{}n}\StringTok{CONTRIBUCIÓN DE LA DATA AL MODELO"}\NormalTok{)}
\BuiltInTok{print}\NormalTok{(}\StringTok{"="} \OperatorTok{*} \DecValTok{70}\NormalTok{)}

\CommentTok{\# Analizar cómo performance cambia con tamaño de datos}
\NormalTok{sample\_sizes }\OperatorTok{=}\NormalTok{ [}\FloatTok{0.1}\NormalTok{, }\FloatTok{0.2}\NormalTok{, }\FloatTok{0.3}\NormalTok{, }\FloatTok{0.4}\NormalTok{, }\FloatTok{0.5}\NormalTok{, }\FloatTok{0.6}\NormalTok{, }\FloatTok{0.7}\NormalTok{, }\FloatTok{0.8}\NormalTok{, }\FloatTok{0.9}\NormalTok{, }\FloatTok{1.0}\NormalTok{]}
\NormalTok{modelo }\OperatorTok{=}\NormalTok{ RandomForestClassifier(n\_estimators}\OperatorTok{=}\DecValTok{100}\NormalTok{, random\_state}\OperatorTok{=}\DecValTok{42}\NormalTok{)}

\NormalTok{resultados\_sizes }\OperatorTok{=}\NormalTok{ []}

\ControlFlowTok{for}\NormalTok{ size }\KeywordTok{in}\NormalTok{ sample\_sizes:}
\NormalTok{    n\_samples }\OperatorTok{=} \BuiltInTok{int}\NormalTok{(}\BuiltInTok{len}\NormalTok{(X) }\OperatorTok{*}\NormalTok{ size)}
    
    \CommentTok{\# Muestrear datos}
\NormalTok{    indices }\OperatorTok{=}\NormalTok{ np.random.choice(}\BuiltInTok{len}\NormalTok{(X), n\_samples, replace}\OperatorTok{=}\VariableTok{False}\NormalTok{)}
\NormalTok{    X\_sample }\OperatorTok{=}\NormalTok{ X[indices]}
\NormalTok{    y\_sample }\OperatorTok{=}\NormalTok{ y[indices]}
    
    \CommentTok{\# Cross{-}validation}
\NormalTok{    scores }\OperatorTok{=}\NormalTok{ cross\_val\_score(modelo, X\_sample, y\_sample, cv}\OperatorTok{=}\DecValTok{5}\NormalTok{, scoring}\OperatorTok{=}\StringTok{\textquotesingle{}roc\_auc\textquotesingle{}}\NormalTok{)}
    
\NormalTok{    resultados\_sizes.append(\{}
        \StringTok{\textquotesingle{}Proporción\textquotesingle{}}\NormalTok{: size,}
        \StringTok{\textquotesingle{}N\_samples\textquotesingle{}}\NormalTok{: n\_samples,}
        \StringTok{\textquotesingle{}Mean\_AUC\textquotesingle{}}\NormalTok{: scores.mean(),}
        \StringTok{\textquotesingle{}Std\_AUC\textquotesingle{}}\NormalTok{: scores.std()}
\NormalTok{    \})}
    
    \BuiltInTok{print}\NormalTok{(}\SpecialStringTok{f"}\CharTok{\textbackslash{}n}\SpecialStringTok{Proporción: }\SpecialCharTok{\{}\NormalTok{size}\SpecialCharTok{:.1f\}}\SpecialStringTok{ (n=}\SpecialCharTok{\{}\NormalTok{n\_samples}\SpecialCharTok{\}}\SpecialStringTok{)"}\NormalTok{)}
    \BuiltInTok{print}\NormalTok{(}\SpecialStringTok{f"  AUC: }\SpecialCharTok{\{}\NormalTok{scores}\SpecialCharTok{.}\NormalTok{mean()}\SpecialCharTok{:.4f\}}\SpecialStringTok{ ± }\SpecialCharTok{\{}\NormalTok{scores}\SpecialCharTok{.}\NormalTok{std()}\SpecialCharTok{:.4f\}}\SpecialStringTok{"}\NormalTok{)}

\NormalTok{df\_sizes }\OperatorTok{=}\NormalTok{ pd.DataFrame(resultados\_sizes)}
\end{Highlighting}
\end{Shaded}

\begin{verbatim}


CONTRIBUCIÓN DE LA DATA AL MODELO
======================================================================
\end{verbatim}

\begin{verbatim}

Proporción: 0.1 (n=15)
  AUC: nan ± nan

Proporción: 0.2 (n=30)
  AUC: 0.8500 ± 0.1458

Proporción: 0.3 (n=45)
  AUC: 0.8778 ± 0.0648

Proporción: 0.4 (n=60)
  AUC: 0.7569 ± 0.1366

Proporción: 0.5 (n=75)
  AUC: 0.7460 ± 0.0774

Proporción: 0.6 (n=90)
  AUC: 0.8688 ± 0.0530

Proporción: 0.7 (n=105)
  AUC: 0.8603 ± 0.0990

Proporción: 0.8 (n=120)
  AUC: 0.7810 ± 0.0831

Proporción: 0.9 (n=135)
  AUC: 0.8302 ± 0.0457

Proporción: 1.0 (n=150)
  AUC: 0.8460 ± 0.0724
\end{verbatim}

\subsection{Visualización}\label{visualizaciuxf3n}

\begin{Shaded}
\begin{Highlighting}[]
\CommentTok{\# Graficar relación tamaño{-}performance}
\NormalTok{fig, axes }\OperatorTok{=}\NormalTok{ plt.subplots(}\DecValTok{1}\NormalTok{, }\DecValTok{2}\NormalTok{, figsize}\OperatorTok{=}\NormalTok{(}\DecValTok{14}\NormalTok{, }\DecValTok{5}\NormalTok{))}

\CommentTok{\# Performance vs tamaño}
\NormalTok{axes[}\DecValTok{0}\NormalTok{].errorbar(df\_sizes[}\StringTok{\textquotesingle{}N\_samples\textquotesingle{}}\NormalTok{], df\_sizes[}\StringTok{\textquotesingle{}Mean\_AUC\textquotesingle{}}\NormalTok{],}
\NormalTok{                yerr}\OperatorTok{=}\NormalTok{df\_sizes[}\StringTok{\textquotesingle{}Std\_AUC\textquotesingle{}}\NormalTok{], marker}\OperatorTok{=}\StringTok{\textquotesingle{}o\textquotesingle{}}\NormalTok{, capsize}\OperatorTok{=}\DecValTok{5}\NormalTok{, linewidth}\OperatorTok{=}\DecValTok{2}\NormalTok{)}
\NormalTok{axes[}\DecValTok{0}\NormalTok{].set\_xlabel(}\StringTok{\textquotesingle{}Tamaño de muestra\textquotesingle{}}\NormalTok{)}
\NormalTok{axes[}\DecValTok{0}\NormalTok{].set\_ylabel(}\StringTok{\textquotesingle{}ROC{-}AUC\textquotesingle{}}\NormalTok{)}
\NormalTok{axes[}\DecValTok{0}\NormalTok{].set\_title(}\StringTok{\textquotesingle{}Performance vs Tamaño de Muestra\textquotesingle{}}\NormalTok{)}
\NormalTok{axes[}\DecValTok{0}\NormalTok{].grid(}\VariableTok{True}\NormalTok{, alpha}\OperatorTok{=}\FloatTok{0.3}\NormalTok{)}

\CommentTok{\# Beneficio marginal}
\NormalTok{marginal\_benefit }\OperatorTok{=}\NormalTok{ df\_sizes[}\StringTok{\textquotesingle{}Mean\_AUC\textquotesingle{}}\NormalTok{].diff()}
\NormalTok{axes[}\DecValTok{1}\NormalTok{].plot(df\_sizes[}\StringTok{\textquotesingle{}N\_samples\textquotesingle{}}\NormalTok{][}\DecValTok{1}\NormalTok{:], marginal\_benefit[}\DecValTok{1}\NormalTok{:], }
            \StringTok{\textquotesingle{}o{-}\textquotesingle{}}\NormalTok{, linewidth}\OperatorTok{=}\DecValTok{2}\NormalTok{)}
\NormalTok{axes[}\DecValTok{1}\NormalTok{].axhline(y}\OperatorTok{=}\DecValTok{0}\NormalTok{, color}\OperatorTok{=}\StringTok{\textquotesingle{}r\textquotesingle{}}\NormalTok{, linestyle}\OperatorTok{=}\StringTok{\textquotesingle{}{-}{-}\textquotesingle{}}\NormalTok{, alpha}\OperatorTok{=}\FloatTok{0.5}\NormalTok{)}
\NormalTok{axes[}\DecValTok{1}\NormalTok{].set\_xlabel(}\StringTok{\textquotesingle{}Tamaño de muestra\textquotesingle{}}\NormalTok{)}
\NormalTok{axes[}\DecValTok{1}\NormalTok{].set\_ylabel(}\StringTok{\textquotesingle{}Beneficio marginal (Δ AUC)\textquotesingle{}}\NormalTok{)}
\NormalTok{axes[}\DecValTok{1}\NormalTok{].set\_title(}\StringTok{\textquotesingle{}Beneficio Marginal de Más Datos\textquotesingle{}}\NormalTok{)}
\NormalTok{axes[}\DecValTok{1}\NormalTok{].grid(}\VariableTok{True}\NormalTok{, alpha}\OperatorTok{=}\FloatTok{0.3}\NormalTok{)}

\NormalTok{plt.tight\_layout()}
\NormalTok{plt.savefig(}\StringTok{\textquotesingle{}data\_contribution.png\textquotesingle{}}\NormalTok{, dpi}\OperatorTok{=}\DecValTok{300}\NormalTok{, bbox\_inches}\OperatorTok{=}\StringTok{\textquotesingle{}tight\textquotesingle{}}\NormalTok{)}
\BuiltInTok{print}\NormalTok{(}\StringTok{"}\CharTok{\textbackslash{}n}\StringTok{Gráfico guardado como \textquotesingle{}data\_contribution.png\textquotesingle{}"}\NormalTok{)}

\CommentTok{\# Análisis de rendimientos}
\BuiltInTok{print}\NormalTok{(}\SpecialStringTok{f"}\CharTok{\textbackslash{}n\textbackslash{}n}\SpecialStringTok{ANÁLISIS DE RENDIMIENTOS:"}\NormalTok{)}
\BuiltInTok{print}\NormalTok{(}\SpecialStringTok{f"Con 10\% de datos (n=}\SpecialCharTok{\{}\BuiltInTok{int}\NormalTok{(}\FloatTok{0.1}\OperatorTok{*}\BuiltInTok{len}\NormalTok{(X))}\SpecialCharTok{\}}\SpecialStringTok{): AUC = }\SpecialCharTok{\{}\NormalTok{df\_sizes}\SpecialCharTok{.}\NormalTok{loc[}\DecValTok{0}\NormalTok{, }\StringTok{\textquotesingle{}Mean\_AUC\textquotesingle{}}\NormalTok{]}\SpecialCharTok{:.4f\}}\SpecialStringTok{"}\NormalTok{)}
\BuiltInTok{print}\NormalTok{(}\SpecialStringTok{f"Con 50\% de datos (n=}\SpecialCharTok{\{}\BuiltInTok{int}\NormalTok{(}\FloatTok{0.5}\OperatorTok{*}\BuiltInTok{len}\NormalTok{(X))}\SpecialCharTok{\}}\SpecialStringTok{): AUC = }\SpecialCharTok{\{}\NormalTok{df\_sizes}\SpecialCharTok{.}\NormalTok{loc[}\DecValTok{4}\NormalTok{, }\StringTok{\textquotesingle{}Mean\_AUC\textquotesingle{}}\NormalTok{]}\SpecialCharTok{:.4f\}}\SpecialStringTok{"}\NormalTok{)}
\BuiltInTok{print}\NormalTok{(}\SpecialStringTok{f"Con 100\% de datos (n=}\SpecialCharTok{\{}\BuiltInTok{len}\NormalTok{(X)}\SpecialCharTok{\}}\SpecialStringTok{): AUC = }\SpecialCharTok{\{}\NormalTok{df\_sizes}\SpecialCharTok{.}\NormalTok{loc[}\DecValTok{9}\NormalTok{, }\StringTok{\textquotesingle{}Mean\_AUC\textquotesingle{}}\NormalTok{]}\SpecialCharTok{:.4f\}}\SpecialStringTok{"}\NormalTok{)}

\NormalTok{improvement\_10\_50 }\OperatorTok{=}\NormalTok{ df\_sizes.loc[}\DecValTok{4}\NormalTok{, }\StringTok{\textquotesingle{}Mean\_AUC\textquotesingle{}}\NormalTok{] }\OperatorTok{{-}}\NormalTok{ df\_sizes.loc[}\DecValTok{0}\NormalTok{, }\StringTok{\textquotesingle{}Mean\_AUC\textquotesingle{}}\NormalTok{]}
\NormalTok{improvement\_50\_100 }\OperatorTok{=}\NormalTok{ df\_sizes.loc[}\DecValTok{9}\NormalTok{, }\StringTok{\textquotesingle{}Mean\_AUC\textquotesingle{}}\NormalTok{] }\OperatorTok{{-}}\NormalTok{ df\_sizes.loc[}\DecValTok{4}\NormalTok{, }\StringTok{\textquotesingle{}Mean\_AUC\textquotesingle{}}\NormalTok{]}

\BuiltInTok{print}\NormalTok{(}\SpecialStringTok{f"}\CharTok{\textbackslash{}n}\SpecialStringTok{Mejora 10\%→50\%: }\SpecialCharTok{\{}\NormalTok{improvement\_10\_50}\SpecialCharTok{:.4f\}}\SpecialStringTok{"}\NormalTok{)}
\BuiltInTok{print}\NormalTok{(}\SpecialStringTok{f"Mejora 50\%→100\%: }\SpecialCharTok{\{}\NormalTok{improvement\_50\_100}\SpecialCharTok{:.4f\}}\SpecialStringTok{"}\NormalTok{)}

\ControlFlowTok{if}\NormalTok{ improvement\_50\_100 }\OperatorTok{\textless{}}\NormalTok{ improvement\_10\_50 }\OperatorTok{/} \DecValTok{2}\NormalTok{:}
    \BuiltInTok{print}\NormalTok{(}\StringTok{"}\CharTok{\textbackslash{}n}\StringTok{⚠️  Rendimientos decrecientes: Más datos tiene beneficio limitado"}\NormalTok{)}
\ControlFlowTok{else}\NormalTok{:}
    \BuiltInTok{print}\NormalTok{(}\StringTok{"}\CharTok{\textbackslash{}n}\StringTok{✓ Más datos aún beneficioso"}\NormalTok{)}
\end{Highlighting}
\end{Shaded}

\begin{verbatim}

Gráfico guardado como 'data_contribution.png'


ANÁLISIS DE RENDIMIENTOS:
Con 10% de datos (n=15): AUC = nan
Con 50% de datos (n=75): AUC = 0.7460
Con 100% de datos (n=150): AUC = 0.8460

Mejora 10%→50%: nan
Mejora 50%→100%: 0.1000

✓ Más datos aún beneficioso
\end{verbatim}

\pandocbounded{\includegraphics[keepaspectratio]{index_files/figure-pdf/cell-17-output-2.pdf}}

\section{Contribución de las
variables}\label{contribuciuxf3n-de-las-variables}

\subsection{Importancia de
características}\label{importancia-de-caracteruxedsticas}

\begin{Shaded}
\begin{Highlighting}[]
\BuiltInTok{print}\NormalTok{(}\StringTok{"}\CharTok{\textbackslash{}n\textbackslash{}n}\StringTok{CONTRIBUCIÓN DE LAS VARIABLES AL MODELO"}\NormalTok{)}
\BuiltInTok{print}\NormalTok{(}\StringTok{"="} \OperatorTok{*} \DecValTok{70}\NormalTok{)}

\CommentTok{\# Usar todas las características de Iris}
\NormalTok{X\_full }\OperatorTok{=}\NormalTok{ iris.data}
\NormalTok{y\_full }\OperatorTok{=}\NormalTok{ (iris.target }\OperatorTok{==} \DecValTok{2}\NormalTok{).astype(}\BuiltInTok{int}\NormalTok{)}

\NormalTok{feature\_names }\OperatorTok{=}\NormalTok{ iris.feature\_names}

\CommentTok{\# Entrenar modelo con todas las características}
\NormalTok{rf\_full }\OperatorTok{=}\NormalTok{ RandomForestClassifier(n\_estimators}\OperatorTok{=}\DecValTok{200}\NormalTok{, random\_state}\OperatorTok{=}\DecValTok{42}\NormalTok{)}
\NormalTok{rf\_full.fit(X\_full, y\_full)}

\CommentTok{\# Obtener importancias}
\NormalTok{importances }\OperatorTok{=}\NormalTok{ rf\_full.feature\_importances\_}
\NormalTok{indices }\OperatorTok{=}\NormalTok{ np.argsort(importances)[::}\OperatorTok{{-}}\DecValTok{1}\NormalTok{]}

\BuiltInTok{print}\NormalTok{(}\StringTok{"}\CharTok{\textbackslash{}n}\StringTok{Importancia de características (Random Forest):"}\NormalTok{)}
\BuiltInTok{print}\NormalTok{(}\StringTok{"="} \OperatorTok{*} \DecValTok{70}\NormalTok{)}
\ControlFlowTok{for}\NormalTok{ i, idx }\KeywordTok{in} \BuiltInTok{enumerate}\NormalTok{(indices):}
    \BuiltInTok{print}\NormalTok{(}\SpecialStringTok{f"}\SpecialCharTok{\{}\NormalTok{i}\OperatorTok{+}\DecValTok{1}\SpecialCharTok{\}}\SpecialStringTok{. }\SpecialCharTok{\{}\NormalTok{feature\_names[idx]}\SpecialCharTok{:\textless{}30\}}\SpecialStringTok{ }\SpecialCharTok{\{}\NormalTok{importances[idx]}\SpecialCharTok{:.4f\}}\SpecialStringTok{"}\NormalTok{)}
\end{Highlighting}
\end{Shaded}

\begin{verbatim}


CONTRIBUCIÓN DE LAS VARIABLES AL MODELO
======================================================================

Importancia de características (Random Forest):
======================================================================
1. petal length (cm)              0.4842
2. petal width (cm)               0.3874
3. sepal length (cm)              0.1014
4. sepal width (cm)               0.0269
\end{verbatim}

\subsection{Análisis de
permutación}\label{anuxe1lisis-de-permutaciuxf3n}

\begin{Shaded}
\begin{Highlighting}[]
\ImportTok{from}\NormalTok{ sklearn.inspection }\ImportTok{import}\NormalTok{ permutation\_importance}

\CommentTok{\# Permutation importance (más robusto)}
\NormalTok{perm\_importance }\OperatorTok{=}\NormalTok{ permutation\_importance(}
\NormalTok{    rf\_full, X\_full, y\_full,}
\NormalTok{    n\_repeats}\OperatorTok{=}\DecValTok{30}\NormalTok{,}
\NormalTok{    random\_state}\OperatorTok{=}\DecValTok{42}\NormalTok{,}
\NormalTok{    n\_jobs}\OperatorTok{={-}}\DecValTok{1}
\NormalTok{)}

\CommentTok{\# Comparar con importancia de Random Forest}
\BuiltInTok{print}\NormalTok{(}\StringTok{"}\CharTok{\textbackslash{}n\textbackslash{}n}\StringTok{COMPARACIÓN DE MÉTODOS DE IMPORTANCIA"}\NormalTok{)}
\BuiltInTok{print}\NormalTok{(}\StringTok{"="} \OperatorTok{*} \DecValTok{70}\NormalTok{)}
\BuiltInTok{print}\NormalTok{(}\SpecialStringTok{f"}\SpecialCharTok{\{}\StringTok{\textquotesingle{}Característica\textquotesingle{}}\SpecialCharTok{:\textless{}30\}}\SpecialStringTok{ }\SpecialCharTok{\{}\StringTok{\textquotesingle{}RF Importance\textquotesingle{}}\SpecialCharTok{:\textless{}15\}}\SpecialStringTok{ }\SpecialCharTok{\{}\StringTok{\textquotesingle{}Perm Importance\textquotesingle{}}\SpecialCharTok{:\textless{}15\}}\SpecialStringTok{"}\NormalTok{)}
\BuiltInTok{print}\NormalTok{(}\StringTok{"{-}"} \OperatorTok{*} \DecValTok{70}\NormalTok{)}

\ControlFlowTok{for}\NormalTok{ i }\KeywordTok{in}\NormalTok{ indices:}
    \BuiltInTok{print}\NormalTok{(}\SpecialStringTok{f"}\SpecialCharTok{\{}\NormalTok{feature\_names[i]}\SpecialCharTok{:\textless{}30\}}\SpecialStringTok{ "}
          \SpecialStringTok{f"}\SpecialCharTok{\{}\NormalTok{importances[i]}\SpecialCharTok{:\textless{}15.4f\}}\SpecialStringTok{ "}
          \SpecialStringTok{f"}\SpecialCharTok{\{}\NormalTok{perm\_importance}\SpecialCharTok{.}\NormalTok{importances\_mean[i]}\SpecialCharTok{:\textless{}15.4f\}}\SpecialStringTok{"}\NormalTok{)}
\end{Highlighting}
\end{Shaded}

\begin{verbatim}


COMPARACIÓN DE MÉTODOS DE IMPORTANCIA
======================================================================
Característica                 RF Importance   Perm Importance
----------------------------------------------------------------------
petal length (cm)              0.4842          0.2142         
petal width (cm)               0.3874          0.1627         
sepal length (cm)              0.1014          0.0084         
sepal width (cm)               0.0269          0.0113         
\end{verbatim}

\subsection{Visualización de
importancia}\label{visualizaciuxf3n-de-importancia}

\begin{Shaded}
\begin{Highlighting}[]
\CommentTok{\# Graficar importancias}
\NormalTok{fig, axes }\OperatorTok{=}\NormalTok{ plt.subplots(}\DecValTok{1}\NormalTok{, }\DecValTok{2}\NormalTok{, figsize}\OperatorTok{=}\NormalTok{(}\DecValTok{14}\NormalTok{, }\DecValTok{5}\NormalTok{))}

\CommentTok{\# Random Forest importance}
\NormalTok{axes[}\DecValTok{0}\NormalTok{].barh(}\BuiltInTok{range}\NormalTok{(}\BuiltInTok{len}\NormalTok{(indices)), importances[indices])}
\NormalTok{axes[}\DecValTok{0}\NormalTok{].set\_yticks(}\BuiltInTok{range}\NormalTok{(}\BuiltInTok{len}\NormalTok{(indices)))}
\NormalTok{axes[}\DecValTok{0}\NormalTok{].set\_yticklabels([feature\_names[i] }\ControlFlowTok{for}\NormalTok{ i }\KeywordTok{in}\NormalTok{ indices])}
\NormalTok{axes[}\DecValTok{0}\NormalTok{].set\_xlabel(}\StringTok{\textquotesingle{}Importancia\textquotesingle{}}\NormalTok{)}
\NormalTok{axes[}\DecValTok{0}\NormalTok{].set\_title(}\StringTok{\textquotesingle{}Feature Importance (Random Forest)\textquotesingle{}}\NormalTok{)}
\NormalTok{axes[}\DecValTok{0}\NormalTok{].invert\_yaxis()}

\CommentTok{\# Permutation importance}
\NormalTok{perm\_sorted\_idx }\OperatorTok{=}\NormalTok{ perm\_importance.importances\_mean.argsort()[::}\OperatorTok{{-}}\DecValTok{1}\NormalTok{]}
\NormalTok{axes[}\DecValTok{1}\NormalTok{].boxplot(}
\NormalTok{    [perm\_importance.importances[i] }\ControlFlowTok{for}\NormalTok{ i }\KeywordTok{in}\NormalTok{ perm\_sorted\_idx],}
\NormalTok{    vert}\OperatorTok{=}\VariableTok{False}\NormalTok{,}
\NormalTok{    labels}\OperatorTok{=}\NormalTok{[feature\_names[i] }\ControlFlowTok{for}\NormalTok{ i }\KeywordTok{in}\NormalTok{ perm\_sorted\_idx]}
\NormalTok{)}
\NormalTok{axes[}\DecValTok{1}\NormalTok{].set\_xlabel(}\StringTok{\textquotesingle{}Decrease in Accuracy\textquotesingle{}}\NormalTok{)}
\NormalTok{axes[}\DecValTok{1}\NormalTok{].set\_title(}\StringTok{\textquotesingle{}Permutation Importance\textquotesingle{}}\NormalTok{)}

\NormalTok{plt.tight\_layout()}
\NormalTok{plt.savefig(}\StringTok{\textquotesingle{}feature\_importance.png\textquotesingle{}}\NormalTok{, dpi}\OperatorTok{=}\DecValTok{300}\NormalTok{, bbox\_inches}\OperatorTok{=}\StringTok{\textquotesingle{}tight\textquotesingle{}}\NormalTok{)}
\BuiltInTok{print}\NormalTok{(}\StringTok{"}\CharTok{\textbackslash{}n}\StringTok{Gráfico guardado como \textquotesingle{}feature\_importance.png\textquotesingle{}"}\NormalTok{)}
\end{Highlighting}
\end{Shaded}

\begin{verbatim}

Gráfico guardado como 'feature_importance.png'
\end{verbatim}

\pandocbounded{\includegraphics[keepaspectratio]{index_files/figure-pdf/cell-20-output-2.pdf}}

\subsection{Selección de
características}\label{selecciuxf3n-de-caracteruxedsticas}

\begin{Shaded}
\begin{Highlighting}[]
\ImportTok{from}\NormalTok{ sklearn.feature\_selection }\ImportTok{import}\NormalTok{ RFECV}

\BuiltInTok{print}\NormalTok{(}\StringTok{"}\CharTok{\textbackslash{}n\textbackslash{}n}\StringTok{SELECCIÓN DE CARACTERÍSTICAS"}\NormalTok{)}
\BuiltInTok{print}\NormalTok{(}\StringTok{"="} \OperatorTok{*} \DecValTok{70}\NormalTok{)}

\CommentTok{\# Recursive Feature Elimination con CV}
\NormalTok{modelo\_seleccion }\OperatorTok{=}\NormalTok{ LogisticRegression(random\_state}\OperatorTok{=}\DecValTok{42}\NormalTok{, max\_iter}\OperatorTok{=}\DecValTok{1000}\NormalTok{)}

\NormalTok{rfecv }\OperatorTok{=}\NormalTok{ RFECV(}
\NormalTok{    estimator}\OperatorTok{=}\NormalTok{modelo\_seleccion,}
\NormalTok{    step}\OperatorTok{=}\DecValTok{1}\NormalTok{,}
\NormalTok{    cv}\OperatorTok{=}\NormalTok{StratifiedKFold(}\DecValTok{5}\NormalTok{),}
\NormalTok{    scoring}\OperatorTok{=}\StringTok{\textquotesingle{}roc\_auc\textquotesingle{}}\NormalTok{,}
\NormalTok{    n\_jobs}\OperatorTok{={-}}\DecValTok{1}
\NormalTok{)}

\NormalTok{rfecv.fit(X\_full, y\_full)}

\BuiltInTok{print}\NormalTok{(}\SpecialStringTok{f"}\CharTok{\textbackslash{}n}\SpecialStringTok{Número óptimo de características: }\SpecialCharTok{\{}\NormalTok{rfecv}\SpecialCharTok{.}\NormalTok{n\_features\_}\SpecialCharTok{\}}\SpecialStringTok{"}\NormalTok{)}
\BuiltInTok{print}\NormalTok{(}\SpecialStringTok{f"}\CharTok{\textbackslash{}n}\SpecialStringTok{Características seleccionadas:"}\NormalTok{)}
\ControlFlowTok{for}\NormalTok{ i, selected }\KeywordTok{in} \BuiltInTok{enumerate}\NormalTok{(rfecv.support\_):}
    \ControlFlowTok{if}\NormalTok{ selected:}
        \BuiltInTok{print}\NormalTok{(}\SpecialStringTok{f"  ✓ }\SpecialCharTok{\{}\NormalTok{feature\_names[i]}\SpecialCharTok{\}}\SpecialStringTok{"}\NormalTok{)}
    \ControlFlowTok{else}\NormalTok{:}
        \BuiltInTok{print}\NormalTok{(}\SpecialStringTok{f"  ✗ }\SpecialCharTok{\{}\NormalTok{feature\_names[i]}\SpecialCharTok{\}}\SpecialStringTok{"}\NormalTok{)}

\CommentTok{\# Graficar}
\NormalTok{plt.figure(figsize}\OperatorTok{=}\NormalTok{(}\DecValTok{10}\NormalTok{, }\DecValTok{6}\NormalTok{))}
\NormalTok{plt.plot(}\BuiltInTok{range}\NormalTok{(}\DecValTok{1}\NormalTok{, }\BuiltInTok{len}\NormalTok{(rfecv.cv\_results\_[}\StringTok{\textquotesingle{}mean\_test\_score\textquotesingle{}}\NormalTok{]) }\OperatorTok{+} \DecValTok{1}\NormalTok{),}
\NormalTok{         rfecv.cv\_results\_[}\StringTok{\textquotesingle{}mean\_test\_score\textquotesingle{}}\NormalTok{], }\StringTok{\textquotesingle{}o{-}\textquotesingle{}}\NormalTok{, linewidth}\OperatorTok{=}\DecValTok{2}\NormalTok{)}
\NormalTok{plt.xlabel(}\StringTok{\textquotesingle{}Número de características\textquotesingle{}}\NormalTok{)}
\NormalTok{plt.ylabel(}\StringTok{\textquotesingle{}ROC{-}AUC (CV)\textquotesingle{}}\NormalTok{)}
\NormalTok{plt.title(}\StringTok{\textquotesingle{}Recursive Feature Elimination con Cross{-}Validation\textquotesingle{}}\NormalTok{)}
\NormalTok{plt.grid(}\VariableTok{True}\NormalTok{, alpha}\OperatorTok{=}\FloatTok{0.3}\NormalTok{)}
\NormalTok{plt.axvline(x}\OperatorTok{=}\NormalTok{rfecv.n\_features\_, color}\OperatorTok{=}\StringTok{\textquotesingle{}r\textquotesingle{}}\NormalTok{, linestyle}\OperatorTok{=}\StringTok{\textquotesingle{}{-}{-}\textquotesingle{}}\NormalTok{,}
\NormalTok{            label}\OperatorTok{=}\SpecialStringTok{f\textquotesingle{}Óptimo: }\SpecialCharTok{\{}\NormalTok{rfecv}\SpecialCharTok{.}\NormalTok{n\_features\_}\SpecialCharTok{\}}\SpecialStringTok{ features\textquotesingle{}}\NormalTok{)}
\NormalTok{plt.legend()}
\NormalTok{plt.tight\_layout()}
\NormalTok{plt.savefig(}\StringTok{\textquotesingle{}rfecv\_results.png\textquotesingle{}}\NormalTok{, dpi}\OperatorTok{=}\DecValTok{300}\NormalTok{, bbox\_inches}\OperatorTok{=}\StringTok{\textquotesingle{}tight\textquotesingle{}}\NormalTok{)}
\BuiltInTok{print}\NormalTok{(}\StringTok{"}\CharTok{\textbackslash{}n}\StringTok{Gráfico guardado como \textquotesingle{}rfecv\_results.png\textquotesingle{}"}\NormalTok{)}
\end{Highlighting}
\end{Shaded}

\begin{verbatim}


SELECCIÓN DE CARACTERÍSTICAS
======================================================================

Número óptimo de características: 4

Características seleccionadas:
  ✓ sepal length (cm)
  ✓ sepal width (cm)
  ✓ petal length (cm)
  ✓ petal width (cm)

Gráfico guardado como 'rfecv_results.png'
\end{verbatim}

\pandocbounded{\includegraphics[keepaspectratio]{index_files/figure-pdf/cell-21-output-2.pdf}}

\subsection{Análisis incremental de
variables}\label{anuxe1lisis-incremental-de-variables}

\begin{Shaded}
\begin{Highlighting}[]
\CommentTok{\# Evaluar contribución incremental de cada variable}
\BuiltInTok{print}\NormalTok{(}\StringTok{"}\CharTok{\textbackslash{}n\textbackslash{}n}\StringTok{ANÁLISIS INCREMENTAL DE VARIABLES"}\NormalTok{)}
\BuiltInTok{print}\NormalTok{(}\StringTok{"="} \OperatorTok{*} \DecValTok{70}\NormalTok{)}

\CommentTok{\# Ordenar por importancia}
\NormalTok{ordered\_features }\OperatorTok{=}\NormalTok{ [indices[i] }\ControlFlowTok{for}\NormalTok{ i }\KeywordTok{in} \BuiltInTok{range}\NormalTok{(}\BuiltInTok{len}\NormalTok{(indices))]}

\NormalTok{resultados\_incremental }\OperatorTok{=}\NormalTok{ []}

\ControlFlowTok{for}\NormalTok{ n\_features }\KeywordTok{in} \BuiltInTok{range}\NormalTok{(}\DecValTok{1}\NormalTok{, }\BuiltInTok{len}\NormalTok{(ordered\_features) }\OperatorTok{+} \DecValTok{1}\NormalTok{):}
\NormalTok{    features\_subset }\OperatorTok{=}\NormalTok{ ordered\_features[:n\_features]}
\NormalTok{    X\_subset }\OperatorTok{=}\NormalTok{ X\_full[:, features\_subset]}
    
    \CommentTok{\# Cross{-}validation}
\NormalTok{    scores }\OperatorTok{=}\NormalTok{ cross\_val\_score(}
\NormalTok{        RandomForestClassifier(n\_estimators}\OperatorTok{=}\DecValTok{100}\NormalTok{, random\_state}\OperatorTok{=}\DecValTok{42}\NormalTok{),}
\NormalTok{        X\_subset, y\_full,}
\NormalTok{        cv}\OperatorTok{=}\DecValTok{5}\NormalTok{,}
\NormalTok{        scoring}\OperatorTok{=}\StringTok{\textquotesingle{}roc\_auc\textquotesingle{}}
\NormalTok{    )}
    
\NormalTok{    resultados\_incremental.append(\{}
        \StringTok{\textquotesingle{}N\_features\textquotesingle{}}\NormalTok{: n\_features,}
        \StringTok{\textquotesingle{}Features\textquotesingle{}}\NormalTok{: [feature\_names[i] }\ControlFlowTok{for}\NormalTok{ i }\KeywordTok{in}\NormalTok{ features\_subset],}
        \StringTok{\textquotesingle{}Mean\_AUC\textquotesingle{}}\NormalTok{: scores.mean(),}
        \StringTok{\textquotesingle{}Std\_AUC\textquotesingle{}}\NormalTok{: scores.std()}
\NormalTok{    \})}
    
    \BuiltInTok{print}\NormalTok{(}\SpecialStringTok{f"}\CharTok{\textbackslash{}n}\SpecialCharTok{\{}\NormalTok{n\_features}\SpecialCharTok{\}}\SpecialStringTok{ feature(s): }\SpecialCharTok{\{}\NormalTok{feature\_names[features\_subset[}\OperatorTok{{-}}\DecValTok{1}\NormalTok{]]}\SpecialCharTok{\}}\SpecialStringTok{"}\NormalTok{)}
    \BuiltInTok{print}\NormalTok{(}\SpecialStringTok{f"  AUC: }\SpecialCharTok{\{}\NormalTok{scores}\SpecialCharTok{.}\NormalTok{mean()}\SpecialCharTok{:.4f\}}\SpecialStringTok{ ± }\SpecialCharTok{\{}\NormalTok{scores}\SpecialCharTok{.}\NormalTok{std()}\SpecialCharTok{:.4f\}}\SpecialStringTok{"}\NormalTok{)}

\NormalTok{df\_incremental }\OperatorTok{=}\NormalTok{ pd.DataFrame(resultados\_incremental)}

\CommentTok{\# Graficar}
\NormalTok{plt.figure(figsize}\OperatorTok{=}\NormalTok{(}\DecValTok{10}\NormalTok{, }\DecValTok{6}\NormalTok{))}
\NormalTok{plt.errorbar(df\_incremental[}\StringTok{\textquotesingle{}N\_features\textquotesingle{}}\NormalTok{], df\_incremental[}\StringTok{\textquotesingle{}Mean\_AUC\textquotesingle{}}\NormalTok{],}
\NormalTok{            yerr}\OperatorTok{=}\NormalTok{df\_incremental[}\StringTok{\textquotesingle{}Std\_AUC\textquotesingle{}}\NormalTok{], marker}\OperatorTok{=}\StringTok{\textquotesingle{}o\textquotesingle{}}\NormalTok{, capsize}\OperatorTok{=}\DecValTok{5}\NormalTok{, linewidth}\OperatorTok{=}\DecValTok{2}\NormalTok{)}
\NormalTok{plt.xlabel(}\StringTok{\textquotesingle{}Número de características (agregadas por importancia)\textquotesingle{}}\NormalTok{)}
\NormalTok{plt.ylabel(}\StringTok{\textquotesingle{}ROC{-}AUC\textquotesingle{}}\NormalTok{)}
\NormalTok{plt.title(}\StringTok{\textquotesingle{}Contribución Incremental de Variables\textquotesingle{}}\NormalTok{)}
\NormalTok{plt.grid(}\VariableTok{True}\NormalTok{, alpha}\OperatorTok{=}\FloatTok{0.3}\NormalTok{)}
\NormalTok{plt.xticks(}\BuiltInTok{range}\NormalTok{(}\DecValTok{1}\NormalTok{, }\BuiltInTok{len}\NormalTok{(feature\_names) }\OperatorTok{+} \DecValTok{1}\NormalTok{))}
\NormalTok{plt.tight\_layout()}
\NormalTok{plt.savefig(}\StringTok{\textquotesingle{}incremental\_features.png\textquotesingle{}}\NormalTok{, dpi}\OperatorTok{=}\DecValTok{300}\NormalTok{, bbox\_inches}\OperatorTok{=}\StringTok{\textquotesingle{}tight\textquotesingle{}}\NormalTok{)}
\BuiltInTok{print}\NormalTok{(}\StringTok{"}\CharTok{\textbackslash{}n}\StringTok{Gráfico guardado como \textquotesingle{}incremental\_features.png\textquotesingle{}"}\NormalTok{)}
\end{Highlighting}
\end{Shaded}

\begin{verbatim}


ANÁLISIS INCREMENTAL DE VARIABLES
======================================================================

1 feature(s): petal length (cm)
  AUC: 0.9740 ± 0.0174

2 feature(s): petal width (cm)
  AUC: 0.9880 ± 0.0240

3 feature(s): sepal length (cm)
  AUC: 0.9885 ± 0.0206

4 feature(s): sepal width (cm)
  AUC: 0.9945 ± 0.0110

Gráfico guardado como 'incremental_features.png'
\end{verbatim}

\pandocbounded{\includegraphics[keepaspectratio]{index_files/figure-pdf/cell-22-output-2.pdf}}

\section{Aplicaciones económicas}\label{aplicaciones-econuxf3micas}

\subsection{Aplicación 1: Predicción de incumplimiento
crediticio}\label{aplicaciuxf3n-1-predicciuxf3n-de-incumplimiento-crediticio}

\begin{Shaded}
\begin{Highlighting}[]
\BuiltInTok{print}\NormalTok{(}\StringTok{"}\CharTok{\textbackslash{}n\textbackslash{}n\textbackslash{}n}\StringTok{APLICACIÓN: OPTIMIZACIÓN DE MODELO DE CRÉDITO"}\NormalTok{)}
\BuiltInTok{print}\NormalTok{(}\StringTok{"="} \OperatorTok{*} \DecValTok{70}\NormalTok{)}

\CommentTok{\# Generar datos sintéticos de crédito}
\NormalTok{np.random.seed(}\DecValTok{42}\NormalTok{)}
\NormalTok{n\_clientes }\OperatorTok{=} \DecValTok{1000}

\NormalTok{datos\_credito }\OperatorTok{=}\NormalTok{ pd.DataFrame(\{}
    \StringTok{\textquotesingle{}ingreso\_mensual\textquotesingle{}}\NormalTok{: np.random.lognormal(}\DecValTok{8}\NormalTok{, }\FloatTok{0.5}\NormalTok{, n\_clientes),}
    \StringTok{\textquotesingle{}edad\textquotesingle{}}\NormalTok{: np.random.normal(}\DecValTok{40}\NormalTok{, }\DecValTok{12}\NormalTok{, n\_clientes),}
    \StringTok{\textquotesingle{}antiguedad\_laboral\textquotesingle{}}\NormalTok{: np.random.exponential(}\DecValTok{5}\NormalTok{, n\_clientes),}
    \StringTok{\textquotesingle{}ratio\_deuda\_ingreso\textquotesingle{}}\NormalTok{: np.random.beta(}\DecValTok{2}\NormalTok{, }\DecValTok{5}\NormalTok{, n\_clientes),}
    \StringTok{\textquotesingle{}num\_creditos\_anteriores\textquotesingle{}}\NormalTok{: np.random.poisson(}\DecValTok{2}\NormalTok{, n\_clientes),}
    \StringTok{\textquotesingle{}score\_crediticio\textquotesingle{}}\NormalTok{: np.random.normal(}\DecValTok{650}\NormalTok{, }\DecValTok{100}\NormalTok{, n\_clientes),}
    \StringTok{\textquotesingle{}tiene\_hipoteca\textquotesingle{}}\NormalTok{: np.random.choice([}\DecValTok{0}\NormalTok{, }\DecValTok{1}\NormalTok{], n\_clientes, p}\OperatorTok{=}\NormalTok{[}\FloatTok{0.7}\NormalTok{, }\FloatTok{0.3}\NormalTok{]),}
    \StringTok{\textquotesingle{}tiene\_ahorros\textquotesingle{}}\NormalTok{: np.random.choice([}\DecValTok{0}\NormalTok{, }\DecValTok{1}\NormalTok{], n\_clientes, p}\OperatorTok{=}\NormalTok{[}\FloatTok{0.4}\NormalTok{, }\FloatTok{0.6}\NormalTok{])}
\NormalTok{\})}

\CommentTok{\# Variable dependiente (probabilidad de incumplimiento)}
\NormalTok{prob\_incumplimiento }\OperatorTok{=}\NormalTok{ (}
    \OperatorTok{{-}}\FloatTok{0.01} \OperatorTok{*}\NormalTok{ datos\_credito[}\StringTok{\textquotesingle{}ingreso\_mensual\textquotesingle{}}\NormalTok{] }\OperatorTok{/} \DecValTok{1000} \OperatorTok{+}
    \OperatorTok{{-}}\FloatTok{0.02} \OperatorTok{*}\NormalTok{ datos\_credito[}\StringTok{\textquotesingle{}edad\textquotesingle{}}\NormalTok{] }\OperatorTok{+}
    \OperatorTok{{-}}\FloatTok{0.03} \OperatorTok{*}\NormalTok{ datos\_credito[}\StringTok{\textquotesingle{}antiguedad\_laboral\textquotesingle{}}\NormalTok{] }\OperatorTok{+}
    \FloatTok{1.5} \OperatorTok{*}\NormalTok{ datos\_credito[}\StringTok{\textquotesingle{}ratio\_deuda\_ingreso\textquotesingle{}}\NormalTok{] }\OperatorTok{+}
    \FloatTok{0.05} \OperatorTok{*}\NormalTok{ datos\_credito[}\StringTok{\textquotesingle{}num\_creditos\_anteriores\textquotesingle{}}\NormalTok{] }\OperatorTok{+}
    \OperatorTok{{-}}\FloatTok{0.005} \OperatorTok{*}\NormalTok{ datos\_credito[}\StringTok{\textquotesingle{}score\_crediticio\textquotesingle{}}\NormalTok{] }\OperatorTok{+}
    \OperatorTok{{-}}\FloatTok{0.3} \OperatorTok{*}\NormalTok{ datos\_credito[}\StringTok{\textquotesingle{}tiene\_hipoteca\textquotesingle{}}\NormalTok{] }\OperatorTok{+}
    \OperatorTok{{-}}\FloatTok{0.4} \OperatorTok{*}\NormalTok{ datos\_credito[}\StringTok{\textquotesingle{}tiene\_ahorros\textquotesingle{}}\NormalTok{] }\OperatorTok{+}
\NormalTok{    np.random.normal(}\DecValTok{0}\NormalTok{, }\FloatTok{0.2}\NormalTok{, n\_clientes)}
\NormalTok{)}

\CommentTok{\# Normalizar a [0,1] y generar incumplimientos}
\ImportTok{from}\NormalTok{ scipy.special }\ImportTok{import}\NormalTok{ expit}
\NormalTok{prob\_incumplimiento }\OperatorTok{=}\NormalTok{ expit(prob\_incumplimiento)}
\NormalTok{datos\_credito[}\StringTok{\textquotesingle{}incumplimiento\textquotesingle{}}\NormalTok{] }\OperatorTok{=}\NormalTok{ (np.random.random(n\_clientes) }\OperatorTok{\textless{}}\NormalTok{ prob\_incumplimiento).astype(}\BuiltInTok{int}\NormalTok{)}

\BuiltInTok{print}\NormalTok{(}\SpecialStringTok{f"Dataset de crédito generado:"}\NormalTok{)}
\BuiltInTok{print}\NormalTok{(}\SpecialStringTok{f"  Total clientes: }\SpecialCharTok{\{}\NormalTok{n\_clientes}\SpecialCharTok{\}}\SpecialStringTok{"}\NormalTok{)}
\BuiltInTok{print}\NormalTok{(}\SpecialStringTok{f"  Tasa de incumplimiento: }\SpecialCharTok{\{}\NormalTok{datos\_credito[}\StringTok{\textquotesingle{}incumplimiento\textquotesingle{}}\NormalTok{]}\SpecialCharTok{.}\NormalTok{mean()}\SpecialCharTok{:.2\%\}}\SpecialStringTok{"}\NormalTok{)}

\CommentTok{\# Preparar datos}
\NormalTok{X\_credito }\OperatorTok{=}\NormalTok{ datos\_credito.drop(}\StringTok{\textquotesingle{}incumplimiento\textquotesingle{}}\NormalTok{, axis}\OperatorTok{=}\DecValTok{1}\NormalTok{)}
\NormalTok{y\_credito }\OperatorTok{=}\NormalTok{ datos\_credito[}\StringTok{\textquotesingle{}incumplimiento\textquotesingle{}}\NormalTok{]}
\end{Highlighting}
\end{Shaded}

\begin{verbatim}



APLICACIÓN: OPTIMIZACIÓN DE MODELO DE CRÉDITO
======================================================================
Dataset de crédito generado:
  Total clientes: 1000
  Tasa de incumplimiento: 2.60%
\end{verbatim}

\subsection{Pipeline completo de
optimización}\label{pipeline-completo-de-optimizaciuxf3n}

\begin{Shaded}
\begin{Highlighting}[]
\CommentTok{\# 1. Comparar múltiples modelos con CV}
\NormalTok{modelos\_credito }\OperatorTok{=}\NormalTok{ \{}
    \StringTok{\textquotesingle{}Logistic Regression\textquotesingle{}}\NormalTok{: LogisticRegression(random\_state}\OperatorTok{=}\DecValTok{42}\NormalTok{, max\_iter}\OperatorTok{=}\DecValTok{1000}\NormalTok{),}
    \StringTok{\textquotesingle{}Random Forest\textquotesingle{}}\NormalTok{: RandomForestClassifier(random\_state}\OperatorTok{=}\DecValTok{42}\NormalTok{),}
    \StringTok{\textquotesingle{}Decision Tree\textquotesingle{}}\NormalTok{: DecisionTreeClassifier(random\_state}\OperatorTok{=}\DecValTok{42}\NormalTok{)}
\NormalTok{\}}

\BuiltInTok{print}\NormalTok{(}\StringTok{"}\CharTok{\textbackslash{}n\textbackslash{}n}\StringTok{COMPARACIÓN DE MODELOS"}\NormalTok{)}
\BuiltInTok{print}\NormalTok{(}\StringTok{"="} \OperatorTok{*} \DecValTok{70}\NormalTok{)}

\NormalTok{resultados\_modelos }\OperatorTok{=}\NormalTok{ []}

\ControlFlowTok{for}\NormalTok{ nombre, modelo }\KeywordTok{in}\NormalTok{ modelos\_credito.items():}
\NormalTok{    cv\_scores }\OperatorTok{=}\NormalTok{ cross\_validate(}
\NormalTok{        modelo, X\_credito, y\_credito,}
\NormalTok{        cv}\OperatorTok{=}\DecValTok{5}\NormalTok{,}
\NormalTok{        scoring}\OperatorTok{=}\NormalTok{[}\StringTok{\textquotesingle{}roc\_auc\textquotesingle{}}\NormalTok{, }\StringTok{\textquotesingle{}precision\textquotesingle{}}\NormalTok{, }\StringTok{\textquotesingle{}recall\textquotesingle{}}\NormalTok{],}
\NormalTok{        return\_train\_score}\OperatorTok{=}\VariableTok{True}
\NormalTok{    )}
    
\NormalTok{    resultados\_modelos.append(\{}
        \StringTok{\textquotesingle{}Modelo\textquotesingle{}}\NormalTok{: nombre,}
        \StringTok{\textquotesingle{}AUC\_test\textquotesingle{}}\NormalTok{: cv\_scores[}\StringTok{\textquotesingle{}test\_roc\_auc\textquotesingle{}}\NormalTok{].mean(),}
        \StringTok{\textquotesingle{}Precision\textquotesingle{}}\NormalTok{: cv\_scores[}\StringTok{\textquotesingle{}test\_precision\textquotesingle{}}\NormalTok{].mean(),}
        \StringTok{\textquotesingle{}Recall\textquotesingle{}}\NormalTok{: cv\_scores[}\StringTok{\textquotesingle{}test\_recall\textquotesingle{}}\NormalTok{].mean()}
\NormalTok{    \})}
    
    \BuiltInTok{print}\NormalTok{(}\SpecialStringTok{f"}\CharTok{\textbackslash{}n}\SpecialCharTok{\{}\NormalTok{nombre}\SpecialCharTok{\}}\SpecialStringTok{:"}\NormalTok{)}
    \BuiltInTok{print}\NormalTok{(}\SpecialStringTok{f"  AUC: }\SpecialCharTok{\{}\NormalTok{cv\_scores[}\StringTok{\textquotesingle{}test\_roc\_auc\textquotesingle{}}\NormalTok{]}\SpecialCharTok{.}\NormalTok{mean()}\SpecialCharTok{:.4f\}}\SpecialStringTok{"}\NormalTok{)}
    \BuiltInTok{print}\NormalTok{(}\SpecialStringTok{f"  Precision: }\SpecialCharTok{\{}\NormalTok{cv\_scores[}\StringTok{\textquotesingle{}test\_precision\textquotesingle{}}\NormalTok{]}\SpecialCharTok{.}\NormalTok{mean()}\SpecialCharTok{:.4f\}}\SpecialStringTok{"}\NormalTok{)}
    \BuiltInTok{print}\NormalTok{(}\SpecialStringTok{f"  Recall: }\SpecialCharTok{\{}\NormalTok{cv\_scores[}\StringTok{\textquotesingle{}test\_recall\textquotesingle{}}\NormalTok{]}\SpecialCharTok{.}\NormalTok{mean()}\SpecialCharTok{:.4f\}}\SpecialStringTok{"}\NormalTok{)}

\NormalTok{df\_modelos }\OperatorTok{=}\NormalTok{ pd.DataFrame(resultados\_modelos)}

\CommentTok{\# 2. Optimizar mejor modelo (Random Forest)}
\BuiltInTok{print}\NormalTok{(}\StringTok{"}\CharTok{\textbackslash{}n\textbackslash{}n}\StringTok{OPTIMIZACIÓN DE HIPERPARÁMETROS"}\NormalTok{)}
\BuiltInTok{print}\NormalTok{(}\StringTok{"="} \OperatorTok{*} \DecValTok{70}\NormalTok{)}

\NormalTok{param\_grid\_credito }\OperatorTok{=}\NormalTok{ \{}
    \StringTok{\textquotesingle{}n\_estimators\textquotesingle{}}\NormalTok{: [}\DecValTok{100}\NormalTok{, }\DecValTok{200}\NormalTok{, }\DecValTok{300}\NormalTok{],}
    \StringTok{\textquotesingle{}max\_depth\textquotesingle{}}\NormalTok{: [}\DecValTok{5}\NormalTok{, }\DecValTok{10}\NormalTok{, }\DecValTok{15}\NormalTok{, }\DecValTok{20}\NormalTok{],}
    \StringTok{\textquotesingle{}min\_samples\_split\textquotesingle{}}\NormalTok{: [}\DecValTok{2}\NormalTok{, }\DecValTok{5}\NormalTok{, }\DecValTok{10}\NormalTok{],}
    \StringTok{\textquotesingle{}class\_weight\textquotesingle{}}\NormalTok{: [}\StringTok{\textquotesingle{}balanced\textquotesingle{}}\NormalTok{, }\VariableTok{None}\NormalTok{]}
\NormalTok{\}}

\NormalTok{grid\_credito }\OperatorTok{=}\NormalTok{ GridSearchCV(}
\NormalTok{    RandomForestClassifier(random\_state}\OperatorTok{=}\DecValTok{42}\NormalTok{),}
\NormalTok{    param\_grid\_credito,}
\NormalTok{    cv}\OperatorTok{=}\DecValTok{5}\NormalTok{,}
\NormalTok{    scoring}\OperatorTok{=}\StringTok{\textquotesingle{}roc\_auc\textquotesingle{}}\NormalTok{,}
\NormalTok{    n\_jobs}\OperatorTok{={-}}\DecValTok{1}\NormalTok{,}
\NormalTok{    verbose}\OperatorTok{=}\DecValTok{0}
\NormalTok{)}

\BuiltInTok{print}\NormalTok{(}\StringTok{"Ejecutando Grid Search..."}\NormalTok{)}
\NormalTok{grid\_credito.fit(X\_credito, y\_credito)}

\BuiltInTok{print}\NormalTok{(}\SpecialStringTok{f"}\CharTok{\textbackslash{}n}\SpecialStringTok{Mejores parámetros:"}\NormalTok{)}
\ControlFlowTok{for}\NormalTok{ param, value }\KeywordTok{in}\NormalTok{ grid\_credito.best\_params\_.items():}
    \BuiltInTok{print}\NormalTok{(}\SpecialStringTok{f"  }\SpecialCharTok{\{}\NormalTok{param}\SpecialCharTok{\}}\SpecialStringTok{: }\SpecialCharTok{\{}\NormalTok{value}\SpecialCharTok{\}}\SpecialStringTok{"}\NormalTok{)}
\BuiltInTok{print}\NormalTok{(}\SpecialStringTok{f"}\CharTok{\textbackslash{}n}\SpecialStringTok{Mejor AUC: }\SpecialCharTok{\{}\NormalTok{grid\_credito}\SpecialCharTok{.}\NormalTok{best\_score\_}\SpecialCharTok{:.4f\}}\SpecialStringTok{"}\NormalTok{)}

\CommentTok{\# 3. Análisis de importancia de variables}
\NormalTok{modelo\_final }\OperatorTok{=}\NormalTok{ grid\_credito.best\_estimator\_}
\NormalTok{importances\_credito }\OperatorTok{=}\NormalTok{ modelo\_final.feature\_importances\_}

\BuiltInTok{print}\NormalTok{(}\SpecialStringTok{f"}\CharTok{\textbackslash{}n\textbackslash{}n}\SpecialStringTok{IMPORTANCIA DE VARIABLES"}\NormalTok{)}
\BuiltInTok{print}\NormalTok{(}\StringTok{"="} \OperatorTok{*} \DecValTok{70}\NormalTok{)}
\NormalTok{importances\_df }\OperatorTok{=}\NormalTok{ pd.DataFrame(\{}
    \StringTok{\textquotesingle{}Variable\textquotesingle{}}\NormalTok{: X\_credito.columns,}
    \StringTok{\textquotesingle{}Importancia\textquotesingle{}}\NormalTok{: importances\_credito}
\NormalTok{\}).sort\_values(}\StringTok{\textquotesingle{}Importancia\textquotesingle{}}\NormalTok{, ascending}\OperatorTok{=}\VariableTok{False}\NormalTok{)}

\ControlFlowTok{for}\NormalTok{ \_, row }\KeywordTok{in}\NormalTok{ importances\_df.iterrows():}
\NormalTok{    bar }\OperatorTok{=} \StringTok{\textquotesingle{}█\textquotesingle{}} \OperatorTok{*} \BuiltInTok{int}\NormalTok{(row[}\StringTok{\textquotesingle{}Importancia\textquotesingle{}}\NormalTok{] }\OperatorTok{*} \DecValTok{50}\NormalTok{)}
    \BuiltInTok{print}\NormalTok{(}\SpecialStringTok{f"}\SpecialCharTok{\{}\NormalTok{row[}\StringTok{\textquotesingle{}Variable\textquotesingle{}}\NormalTok{]}\SpecialCharTok{:\textless{}25\}}\SpecialStringTok{ }\SpecialCharTok{\{}\NormalTok{row[}\StringTok{\textquotesingle{}Importancia\textquotesingle{}}\NormalTok{]}\SpecialCharTok{:.4f\}}\SpecialStringTok{ }\SpecialCharTok{\{}\NormalTok{bar}\SpecialCharTok{\}}\SpecialStringTok{"}\NormalTok{)}

\CommentTok{\# 4. Curvas de aprendizaje}
\BuiltInTok{print}\NormalTok{(}\SpecialStringTok{f"}\CharTok{\textbackslash{}n\textbackslash{}n}\SpecialStringTok{CURVAS DE APRENDIZAJE"}\NormalTok{)}
\BuiltInTok{print}\NormalTok{(}\StringTok{"="} \OperatorTok{*} \DecValTok{70}\NormalTok{)}

\NormalTok{train\_sizes, train\_scores, val\_scores }\OperatorTok{=}\NormalTok{ learning\_curve(}
\NormalTok{    modelo\_final, X\_credito, y\_credito,}
\NormalTok{    cv}\OperatorTok{=}\DecValTok{5}\NormalTok{,}
\NormalTok{    n\_jobs}\OperatorTok{={-}}\DecValTok{1}\NormalTok{,}
\NormalTok{    train\_sizes}\OperatorTok{=}\NormalTok{np.linspace(}\FloatTok{0.1}\NormalTok{, }\FloatTok{1.0}\NormalTok{, }\DecValTok{10}\NormalTok{),}
\NormalTok{    scoring}\OperatorTok{=}\StringTok{\textquotesingle{}roc\_auc\textquotesingle{}}
\NormalTok{)}

\NormalTok{plt.figure(figsize}\OperatorTok{=}\NormalTok{(}\DecValTok{10}\NormalTok{, }\DecValTok{6}\NormalTok{))}
\NormalTok{plt.plot(train\_sizes, train\_scores.mean(axis}\OperatorTok{=}\DecValTok{1}\NormalTok{), }\StringTok{\textquotesingle{}o{-}\textquotesingle{}}\NormalTok{, label}\OperatorTok{=}\StringTok{\textquotesingle{}Train\textquotesingle{}}\NormalTok{, linewidth}\OperatorTok{=}\DecValTok{2}\NormalTok{)}
\NormalTok{plt.plot(train\_sizes, val\_scores.mean(axis}\OperatorTok{=}\DecValTok{1}\NormalTok{), }\StringTok{\textquotesingle{}s{-}\textquotesingle{}}\NormalTok{, label}\OperatorTok{=}\StringTok{\textquotesingle{}Validation\textquotesingle{}}\NormalTok{, linewidth}\OperatorTok{=}\DecValTok{2}\NormalTok{)}
\NormalTok{plt.xlabel(}\StringTok{\textquotesingle{}Tamaño de muestra\textquotesingle{}}\NormalTok{)}
\NormalTok{plt.ylabel(}\StringTok{\textquotesingle{}ROC{-}AUC\textquotesingle{}}\NormalTok{)}
\NormalTok{plt.title(}\StringTok{\textquotesingle{}Curvas de Aprendizaje {-} Modelo de Crédito Optimizado\textquotesingle{}}\NormalTok{)}
\NormalTok{plt.legend()}
\NormalTok{plt.grid(}\VariableTok{True}\NormalTok{, alpha}\OperatorTok{=}\FloatTok{0.3}\NormalTok{)}
\NormalTok{plt.tight\_layout()}
\NormalTok{plt.savefig(}\StringTok{\textquotesingle{}learning\_curve\_credito.png\textquotesingle{}}\NormalTok{, dpi}\OperatorTok{=}\DecValTok{300}\NormalTok{, bbox\_inches}\OperatorTok{=}\StringTok{\textquotesingle{}tight\textquotesingle{}}\NormalTok{)}
\BuiltInTok{print}\NormalTok{(}\StringTok{"Curva guardada como \textquotesingle{}learning\_curve\_credito.png\textquotesingle{}"}\NormalTok{)}

\CommentTok{\# Diagnóstico final}
\NormalTok{gap\_final }\OperatorTok{=}\NormalTok{ train\_scores.mean(axis}\OperatorTok{=}\DecValTok{1}\NormalTok{)[}\OperatorTok{{-}}\DecValTok{1}\NormalTok{] }\OperatorTok{{-}}\NormalTok{ val\_scores.mean(axis}\OperatorTok{=}\DecValTok{1}\NormalTok{)[}\OperatorTok{{-}}\DecValTok{1}\NormalTok{]}
\BuiltInTok{print}\NormalTok{(}\SpecialStringTok{f"}\CharTok{\textbackslash{}n}\SpecialStringTok{Diagnóstico:"}\NormalTok{)}
\BuiltInTok{print}\NormalTok{(}\SpecialStringTok{f"  Training AUC: }\SpecialCharTok{\{}\NormalTok{train\_scores}\SpecialCharTok{.}\NormalTok{mean(axis}\OperatorTok{=}\DecValTok{1}\NormalTok{)[}\OperatorTok{{-}}\DecValTok{1}\NormalTok{]}\SpecialCharTok{:.4f\}}\SpecialStringTok{"}\NormalTok{)}
\BuiltInTok{print}\NormalTok{(}\SpecialStringTok{f"  Validation AUC: }\SpecialCharTok{\{}\NormalTok{val\_scores}\SpecialCharTok{.}\NormalTok{mean(axis}\OperatorTok{=}\DecValTok{1}\NormalTok{)[}\OperatorTok{{-}}\DecValTok{1}\NormalTok{]}\SpecialCharTok{:.4f\}}\SpecialStringTok{"}\NormalTok{)}
\BuiltInTok{print}\NormalTok{(}\SpecialStringTok{f"  Gap: }\SpecialCharTok{\{}\NormalTok{gap\_final}\SpecialCharTok{:.4f\}}\SpecialStringTok{"}\NormalTok{)}

\ControlFlowTok{if}\NormalTok{ gap\_final }\OperatorTok{\textless{}} \FloatTok{0.05}\NormalTok{:}
    \BuiltInTok{print}\NormalTok{(}\StringTok{"  ✓ Modelo bien balanceado"}\NormalTok{)}
\ControlFlowTok{elif}\NormalTok{ gap\_final }\OperatorTok{\textless{}} \FloatTok{0.15}\NormalTok{:}
    \BuiltInTok{print}\NormalTok{(}\StringTok{"  ⚠️  Ligero overfitting"}\NormalTok{)}
\ControlFlowTok{else}\NormalTok{:}
    \BuiltInTok{print}\NormalTok{(}\StringTok{"  ⚠️  Overfitting significativo"}\NormalTok{)}
\end{Highlighting}
\end{Shaded}

\begin{verbatim}


COMPARACIÓN DE MODELOS
======================================================================
\end{verbatim}

\begin{verbatim}

Logistic Regression:
  AUC: 0.6643
  Precision: 0.0000
  Recall: 0.0000
\end{verbatim}

\begin{verbatim}

Random Forest:
  AUC: 0.5305
  Precision: 0.0000
  Recall: 0.0000

Decision Tree:
  AUC: 0.5261
  Precision: 0.0833
  Recall: 0.0800


OPTIMIZACIÓN DE HIPERPARÁMETROS
======================================================================
Ejecutando Grid Search...

Mejores parámetros:
  class_weight: balanced
  max_depth: 10
  min_samples_split: 2
  n_estimators: 300

Mejor AUC: 0.6431


IMPORTANCIA DE VARIABLES
======================================================================
antiguedad_laboral        0.2290 ███████████
score_crediticio          0.1828 █████████
ratio_deuda_ingreso       0.1591 ███████
edad                      0.1563 ███████
ingreso_mensual           0.1367 ██████
num_creditos_anteriores   0.0709 ███
tiene_ahorros             0.0384 █
tiene_hipoteca            0.0267 █


CURVAS DE APRENDIZAJE
======================================================================
Curva guardada como 'learning_curve_credito.png'

Diagnóstico:
  Training AUC: 1.0000
  Validation AUC: 0.6431
  Gap: 0.3569
  ⚠️  Overfitting significativo
\end{verbatim}

\pandocbounded{\includegraphics[keepaspectratio]{index_files/figure-pdf/cell-24-output-4.pdf}}

\section{Mejores prácticas}\label{mejores-pruxe1cticas}

\subsection{Checklist de validación de
modelos}\label{checklist-de-validaciuxf3n-de-modelos}

\begin{Shaded}
\begin{Highlighting}[]
\BuiltInTok{print}\NormalTok{(}\StringTok{"}\CharTok{\textbackslash{}n\textbackslash{}n\textbackslash{}n}\StringTok{CHECKLIST DE VALIDACIÓN DE MODELOS"}\NormalTok{)}
\BuiltInTok{print}\NormalTok{(}\StringTok{"="} \OperatorTok{*} \DecValTok{70}\NormalTok{)}

\NormalTok{checklist }\OperatorTok{=} \StringTok{"""}
\StringTok{□ DIVISIÓN DE DATOS}
\StringTok{  ✓ Separar train/test antes de cualquier análisis}
\StringTok{  ✓ Usar stratified split para clases desbalanceadas}
\StringTok{  ✓ Mantener test set intocado hasta evaluación final}

\StringTok{□ VALIDACIÓN CRUZADA}
\StringTok{  ✓ Usar K{-}Fold (típicamente K=5 o K=10)}
\StringTok{  ✓ Stratified K{-}Fold para clasificación}
\StringTok{  ✓ Repetir con diferentes seeds}
\StringTok{  ✓ Reportar media ± desviación estándar}

\StringTok{□ OPTIMIZACIÓN DE HIPERPARÁMETROS}
\StringTok{  ✓ Definir grilla razonable de hiperparámetros}
\StringTok{  ✓ Usar validación cruzada anidada}
\StringTok{  ✓ Considerar Random Search para espacios grandes}
\StringTok{  ✓ Guardar mejores parámetros}

\StringTok{□ EVALUACIÓN DE OVERFITTING}
\StringTok{  ✓ Comparar performance train vs test}
\StringTok{  ✓ Analizar curvas de aprendizaje}
\StringTok{  ✓ Verificar gap \textless{} 0.15}
\StringTok{  ✓ Usar regularización si es necesario}

\StringTok{□ ANÁLISIS DE DATOS}
\StringTok{  ✓ Evaluar si más datos ayudarían}
\StringTok{  ✓ Verificar rendimientos decrecientes}
\StringTok{  ✓ Considerar costo de recolección}

\StringTok{□ SELECCIÓN DE VARIABLES}
\StringTok{  ✓ Analizar importancia de características}
\StringTok{  ✓ Considerar RFECV para reducir dimensionalidad}
\StringTok{  ✓ Evaluar contribución incremental}
\StringTok{  ✓ Remover variables redundantes}

\StringTok{□ MÉTRICAS APROPIADAS}
\StringTok{  ✓ Elegir métrica según objetivo de negocio}
\StringTok{  ✓ Considerar costos de errores (FP vs FN)}
\StringTok{  ✓ Reportar múltiples métricas}
\StringTok{  ✓ Usar curvas ROC/PR para panorama completo}

\StringTok{□ ROBUSTEZ}
\StringTok{  ✓ Probar con diferentes particiones}
\StringTok{  ✓ Validar en período temporal diferente}
\StringTok{  ✓ Verificar estabilidad con diferentes seeds}
\StringTok{  ✓ Detectar data leakage}

\StringTok{□ DOCUMENTACIÓN}
\StringTok{  ✓ Guardar hiperparámetros óptimos}
\StringTok{  ✓ Documentar proceso de selección}
\StringTok{  ✓ Registrar performance en test}
\StringTok{  ✓ Explicar decisiones de diseño}

\StringTok{□ COMPARACIÓN}
\StringTok{  ✓ Establecer baseline simple}
\StringTok{  ✓ Comparar múltiples algoritmos}
\StringTok{  ✓ Usar mismos datos para comparación}
\StringTok{  ✓ Considerar trade{-}offs (interpretabilidad vs accuracy)}
\StringTok{"""}

\BuiltInTok{print}\NormalTok{(checklist)}
\end{Highlighting}
\end{Shaded}

\begin{verbatim}



CHECKLIST DE VALIDACIÓN DE MODELOS
======================================================================

□ DIVISIÓN DE DATOS
  ✓ Separar train/test antes de cualquier análisis
  ✓ Usar stratified split para clases desbalanceadas
  ✓ Mantener test set intocado hasta evaluación final

□ VALIDACIÓN CRUZADA
  ✓ Usar K-Fold (típicamente K=5 o K=10)
  ✓ Stratified K-Fold para clasificación
  ✓ Repetir con diferentes seeds
  ✓ Reportar media ± desviación estándar

□ OPTIMIZACIÓN DE HIPERPARÁMETROS
  ✓ Definir grilla razonable de hiperparámetros
  ✓ Usar validación cruzada anidada
  ✓ Considerar Random Search para espacios grandes
  ✓ Guardar mejores parámetros

□ EVALUACIÓN DE OVERFITTING
  ✓ Comparar performance train vs test
  ✓ Analizar curvas de aprendizaje
  ✓ Verificar gap < 0.15
  ✓ Usar regularización si es necesario

□ ANÁLISIS DE DATOS
  ✓ Evaluar si más datos ayudarían
  ✓ Verificar rendimientos decrecientes
  ✓ Considerar costo de recolección

□ SELECCIÓN DE VARIABLES
  ✓ Analizar importancia de características
  ✓ Considerar RFECV para reducir dimensionalidad
  ✓ Evaluar contribución incremental
  ✓ Remover variables redundantes

□ MÉTRICAS APROPIADAS
  ✓ Elegir métrica según objetivo de negocio
  ✓ Considerar costos de errores (FP vs FN)
  ✓ Reportar múltiples métricas
  ✓ Usar curvas ROC/PR para panorama completo

□ ROBUSTEZ
  ✓ Probar con diferentes particiones
  ✓ Validar en período temporal diferente
  ✓ Verificar estabilidad con diferentes seeds
  ✓ Detectar data leakage

□ DOCUMENTACIÓN
  ✓ Guardar hiperparámetros óptimos
  ✓ Documentar proceso de selección
  ✓ Registrar performance en test
  ✓ Explicar decisiones de diseño

□ COMPARACIÓN
  ✓ Establecer baseline simple
  ✓ Comparar múltiples algoritmos
  ✓ Usar mismos datos para comparación
  ✓ Considerar trade-offs (interpretabilidad vs accuracy)
\end{verbatim}

\subsection{Función de pipeline
completo}\label{funciuxf3n-de-pipeline-completo}

\begin{Shaded}
\begin{Highlighting}[]
\KeywordTok{def}\NormalTok{ pipeline\_validacion\_completo(X, y, modelos, param\_grids, feature\_names}\OperatorTok{=}\VariableTok{None}\NormalTok{):}
    \CommentTok{"""}
\CommentTok{    Pipeline completo de validación y optimización de modelos}
\CommentTok{    }
\CommentTok{    Parameters:}
\CommentTok{    {-}{-}{-}{-}{-}{-}{-}{-}{-}{-}{-}}
\CommentTok{    X : array{-}like}
\CommentTok{        Características}
\CommentTok{    y : array{-}like}
\CommentTok{        Variable objetivo}
\CommentTok{    modelos : dict}
\CommentTok{        Diccionario con modelos a evaluar}
\CommentTok{    param\_grids : dict}
\CommentTok{        Diccionario con grillas de hiperparámetros}
\CommentTok{    feature\_names : list, optional}
\CommentTok{        Nombres de características}
\CommentTok{    }
\CommentTok{    Returns:}
\CommentTok{    {-}{-}{-}{-}{-}{-}{-}{-}}
\CommentTok{    dict}
\CommentTok{        Resultados completos del pipeline}
\CommentTok{    """}
    
    \BuiltInTok{print}\NormalTok{(}\StringTok{"="}\OperatorTok{*}\DecValTok{70}\NormalTok{)}
    \BuiltInTok{print}\NormalTok{(}\StringTok{"PIPELINE DE VALIDACIÓN COMPLETO"}\NormalTok{)}
    \BuiltInTok{print}\NormalTok{(}\StringTok{"="}\OperatorTok{*}\DecValTok{70}\NormalTok{)}
    
\NormalTok{    resultados }\OperatorTok{=}\NormalTok{ \{}
        \StringTok{\textquotesingle{}comparacion\_modelos\textquotesingle{}}\NormalTok{: [],}
        \StringTok{\textquotesingle{}mejores\_modelos\textquotesingle{}}\NormalTok{: \{\},}
        \StringTok{\textquotesingle{}importancias\textquotesingle{}}\NormalTok{: \{\},}
        \StringTok{\textquotesingle{}diagnostico\textquotesingle{}}\NormalTok{: \{\}}
\NormalTok{    \}}
    
    \CommentTok{\# 1. Comparación inicial de modelos}
    \BuiltInTok{print}\NormalTok{(}\StringTok{"}\CharTok{\textbackslash{}n}\StringTok{1. COMPARACIÓN INICIAL DE MODELOS"}\NormalTok{)}
    \BuiltInTok{print}\NormalTok{(}\StringTok{"{-}"}\OperatorTok{*}\DecValTok{70}\NormalTok{)}
    
    \ControlFlowTok{for}\NormalTok{ nombre, modelo }\KeywordTok{in}\NormalTok{ modelos.items():}
        \BuiltInTok{print}\NormalTok{(}\SpecialStringTok{f"}\CharTok{\textbackslash{}n}\SpecialStringTok{Evaluando }\SpecialCharTok{\{}\NormalTok{nombre}\SpecialCharTok{\}}\SpecialStringTok{..."}\NormalTok{)}
        
\NormalTok{        cv\_results }\OperatorTok{=}\NormalTok{ cross\_validate(}
\NormalTok{            modelo, X, y,}
\NormalTok{            cv}\OperatorTok{=}\NormalTok{StratifiedKFold(}\DecValTok{5}\NormalTok{, shuffle}\OperatorTok{=}\VariableTok{True}\NormalTok{, random\_state}\OperatorTok{=}\DecValTok{42}\NormalTok{),}
\NormalTok{            scoring}\OperatorTok{=}\NormalTok{[}\StringTok{\textquotesingle{}roc\_auc\textquotesingle{}}\NormalTok{, }\StringTok{\textquotesingle{}precision\textquotesingle{}}\NormalTok{, }\StringTok{\textquotesingle{}recall\textquotesingle{}}\NormalTok{, }\StringTok{\textquotesingle{}f1\textquotesingle{}}\NormalTok{],}
\NormalTok{            return\_train\_score}\OperatorTok{=}\VariableTok{True}\NormalTok{,}
\NormalTok{            n\_jobs}\OperatorTok{={-}}\DecValTok{1}
\NormalTok{        )}
        
\NormalTok{        resultados[}\StringTok{\textquotesingle{}comparacion\_modelos\textquotesingle{}}\NormalTok{].append(\{}
            \StringTok{\textquotesingle{}Modelo\textquotesingle{}}\NormalTok{: nombre,}
            \StringTok{\textquotesingle{}AUC\_train\textquotesingle{}}\NormalTok{: cv\_results[}\StringTok{\textquotesingle{}train\_roc\_auc\textquotesingle{}}\NormalTok{].mean(),}
            \StringTok{\textquotesingle{}AUC\_test\textquotesingle{}}\NormalTok{: cv\_results[}\StringTok{\textquotesingle{}test\_roc\_auc\textquotesingle{}}\NormalTok{].mean(),}
            \StringTok{\textquotesingle{}AUC\_std\textquotesingle{}}\NormalTok{: cv\_results[}\StringTok{\textquotesingle{}test\_roc\_auc\textquotesingle{}}\NormalTok{].std(),}
            \StringTok{\textquotesingle{}Precision\textquotesingle{}}\NormalTok{: cv\_results[}\StringTok{\textquotesingle{}test\_precision\textquotesingle{}}\NormalTok{].mean(),}
            \StringTok{\textquotesingle{}Recall\textquotesingle{}}\NormalTok{: cv\_results[}\StringTok{\textquotesingle{}test\_recall\textquotesingle{}}\NormalTok{].mean(),}
            \StringTok{\textquotesingle{}F1\textquotesingle{}}\NormalTok{: cv\_results[}\StringTok{\textquotesingle{}test\_f1\textquotesingle{}}\NormalTok{].mean(),}
            \StringTok{\textquotesingle{}Gap\textquotesingle{}}\NormalTok{: cv\_results[}\StringTok{\textquotesingle{}train\_roc\_auc\textquotesingle{}}\NormalTok{].mean() }\OperatorTok{{-}}\NormalTok{ cv\_results[}\StringTok{\textquotesingle{}test\_roc\_auc\textquotesingle{}}\NormalTok{].mean()}
\NormalTok{        \})}
        
        \BuiltInTok{print}\NormalTok{(}\SpecialStringTok{f"  AUC: }\SpecialCharTok{\{}\NormalTok{cv\_results[}\StringTok{\textquotesingle{}test\_roc\_auc\textquotesingle{}}\NormalTok{]}\SpecialCharTok{.}\NormalTok{mean()}\SpecialCharTok{:.4f\}}\SpecialStringTok{ ± "}
              \SpecialStringTok{f"}\SpecialCharTok{\{}\NormalTok{cv\_results[}\StringTok{\textquotesingle{}test\_roc\_auc\textquotesingle{}}\NormalTok{]}\SpecialCharTok{.}\NormalTok{std()}\SpecialCharTok{:.4f\}}\SpecialStringTok{"}\NormalTok{)}
        \BuiltInTok{print}\NormalTok{(}\SpecialStringTok{f"  Gap: }\SpecialCharTok{\{}\NormalTok{cv\_results[}\StringTok{\textquotesingle{}train\_roc\_auc\textquotesingle{}}\NormalTok{]}\SpecialCharTok{.}\NormalTok{mean() }\OperatorTok{{-}}\NormalTok{ cv\_results[}\StringTok{\textquotesingle{}test\_roc\_auc\textquotesingle{}}\NormalTok{]}\SpecialCharTok{.}\NormalTok{mean()}\SpecialCharTok{:.4f\}}\SpecialStringTok{"}\NormalTok{)}
    
\NormalTok{    df\_comparacion }\OperatorTok{=}\NormalTok{ pd.DataFrame(resultados[}\StringTok{\textquotesingle{}comparacion\_modelos\textquotesingle{}}\NormalTok{])}
    
    \CommentTok{\# 2. Optimización de hiperparámetros}
    \BuiltInTok{print}\NormalTok{(}\StringTok{"}\CharTok{\textbackslash{}n\textbackslash{}n}\StringTok{2. OPTIMIZACIÓN DE HIPERPARÁMETROS"}\NormalTok{)}
    \BuiltInTok{print}\NormalTok{(}\StringTok{"{-}"}\OperatorTok{*}\DecValTok{70}\NormalTok{)}
    
    \ControlFlowTok{for}\NormalTok{ nombre }\KeywordTok{in}\NormalTok{ modelos.keys():}
        \ControlFlowTok{if}\NormalTok{ nombre }\KeywordTok{in}\NormalTok{ param\_grids:}
            \BuiltInTok{print}\NormalTok{(}\SpecialStringTok{f"}\CharTok{\textbackslash{}n}\SpecialStringTok{Optimizando }\SpecialCharTok{\{}\NormalTok{nombre}\SpecialCharTok{\}}\SpecialStringTok{..."}\NormalTok{)}
            
\NormalTok{            grid }\OperatorTok{=}\NormalTok{ GridSearchCV(}
\NormalTok{                modelos[nombre],}
\NormalTok{                param\_grids[nombre],}
\NormalTok{                cv}\OperatorTok{=}\DecValTok{5}\NormalTok{,}
\NormalTok{                scoring}\OperatorTok{=}\StringTok{\textquotesingle{}roc\_auc\textquotesingle{}}\NormalTok{,}
\NormalTok{                n\_jobs}\OperatorTok{={-}}\DecValTok{1}\NormalTok{,}
\NormalTok{                verbose}\OperatorTok{=}\DecValTok{0}
\NormalTok{            )}
            
\NormalTok{            grid.fit(X, y)}
            
\NormalTok{            resultados[}\StringTok{\textquotesingle{}mejores\_modelos\textquotesingle{}}\NormalTok{][nombre] }\OperatorTok{=}\NormalTok{ \{}
                \StringTok{\textquotesingle{}modelo\textquotesingle{}}\NormalTok{: grid.best\_estimator\_,}
                \StringTok{\textquotesingle{}params\textquotesingle{}}\NormalTok{: grid.best\_params\_,}
                \StringTok{\textquotesingle{}score\textquotesingle{}}\NormalTok{: grid.best\_score\_}
\NormalTok{            \}}
            
            \BuiltInTok{print}\NormalTok{(}\SpecialStringTok{f"  Mejor AUC: }\SpecialCharTok{\{}\NormalTok{grid}\SpecialCharTok{.}\NormalTok{best\_score\_}\SpecialCharTok{:.4f\}}\SpecialStringTok{"}\NormalTok{)}
            \BuiltInTok{print}\NormalTok{(}\SpecialStringTok{f"  Mejores parámetros: }\SpecialCharTok{\{}\NormalTok{grid}\SpecialCharTok{.}\NormalTok{best\_params\_}\SpecialCharTok{\}}\SpecialStringTok{"}\NormalTok{)}
    
    \CommentTok{\# 3. Análisis de importancia (para modelos con feature\_importances\_)}
    \BuiltInTok{print}\NormalTok{(}\StringTok{"}\CharTok{\textbackslash{}n\textbackslash{}n}\StringTok{3. ANÁLISIS DE IMPORTANCIA DE VARIABLES"}\NormalTok{)}
    \BuiltInTok{print}\NormalTok{(}\StringTok{"{-}"}\OperatorTok{*}\DecValTok{70}\NormalTok{)}
    
    \ControlFlowTok{for}\NormalTok{ nombre, info }\KeywordTok{in}\NormalTok{ resultados[}\StringTok{\textquotesingle{}mejores\_modelos\textquotesingle{}}\NormalTok{].items():}
\NormalTok{        modelo }\OperatorTok{=}\NormalTok{ info[}\StringTok{\textquotesingle{}modelo\textquotesingle{}}\NormalTok{]}
        
        \ControlFlowTok{if} \BuiltInTok{hasattr}\NormalTok{(modelo, }\StringTok{\textquotesingle{}feature\_importances\_\textquotesingle{}}\NormalTok{):}
\NormalTok{            importances }\OperatorTok{=}\NormalTok{ modelo.feature\_importances\_}
\NormalTok{            resultados[}\StringTok{\textquotesingle{}importancias\textquotesingle{}}\NormalTok{][nombre] }\OperatorTok{=}\NormalTok{ importances}
            
            \BuiltInTok{print}\NormalTok{(}\SpecialStringTok{f"}\CharTok{\textbackslash{}n}\SpecialCharTok{\{}\NormalTok{nombre}\SpecialCharTok{\}}\SpecialStringTok{:"}\NormalTok{)}
            \ControlFlowTok{if}\NormalTok{ feature\_names }\KeywordTok{is} \KeywordTok{not} \VariableTok{None}\NormalTok{:}
                \ControlFlowTok{for}\NormalTok{ i }\KeywordTok{in}\NormalTok{ np.argsort(importances)[::}\OperatorTok{{-}}\DecValTok{1}\NormalTok{]:}
                    \BuiltInTok{print}\NormalTok{(}\SpecialStringTok{f"  }\SpecialCharTok{\{}\NormalTok{feature\_names[i]}\SpecialCharTok{:\textless{}30\}}\SpecialStringTok{ }\SpecialCharTok{\{}\NormalTok{importances[i]}\SpecialCharTok{:.4f\}}\SpecialStringTok{"}\NormalTok{)}
            \ControlFlowTok{else}\NormalTok{:}
                \ControlFlowTok{for}\NormalTok{ i, imp }\KeywordTok{in} \BuiltInTok{enumerate}\NormalTok{(importances):}
                    \BuiltInTok{print}\NormalTok{(}\SpecialStringTok{f"  Feature }\SpecialCharTok{\{}\NormalTok{i}\SpecialCharTok{\}}\SpecialStringTok{: }\SpecialCharTok{\{}\NormalTok{imp}\SpecialCharTok{:.4f\}}\SpecialStringTok{"}\NormalTok{)}
    
    \CommentTok{\# 4. Diagnóstico final}
    \BuiltInTok{print}\NormalTok{(}\StringTok{"}\CharTok{\textbackslash{}n\textbackslash{}n}\StringTok{4. DIAGNÓSTICO Y RECOMENDACIONES"}\NormalTok{)}
    \BuiltInTok{print}\NormalTok{(}\StringTok{"{-}"}\OperatorTok{*}\DecValTok{70}\NormalTok{)}
    
\NormalTok{    mejor\_modelo\_nombre }\OperatorTok{=}\NormalTok{ df\_comparacion.loc[df\_comparacion[}\StringTok{\textquotesingle{}AUC\_test\textquotesingle{}}\NormalTok{].idxmax(), }\StringTok{\textquotesingle{}Modelo\textquotesingle{}}\NormalTok{]}
\NormalTok{    mejor\_auc }\OperatorTok{=}\NormalTok{ df\_comparacion[}\StringTok{\textquotesingle{}AUC\_test\textquotesingle{}}\NormalTok{].}\BuiltInTok{max}\NormalTok{()}
    
\NormalTok{    resultados[}\StringTok{\textquotesingle{}diagnostico\textquotesingle{}}\NormalTok{][}\StringTok{\textquotesingle{}mejor\_modelo\textquotesingle{}}\NormalTok{] }\OperatorTok{=}\NormalTok{ mejor\_modelo\_nombre}
\NormalTok{    resultados[}\StringTok{\textquotesingle{}diagnostico\textquotesingle{}}\NormalTok{][}\StringTok{\textquotesingle{}mejor\_auc\textquotesingle{}}\NormalTok{] }\OperatorTok{=}\NormalTok{ mejor\_auc}
    
    \BuiltInTok{print}\NormalTok{(}\SpecialStringTok{f"}\CharTok{\textbackslash{}n}\SpecialStringTok{Mejor modelo: }\SpecialCharTok{\{}\NormalTok{mejor\_modelo\_nombre}\SpecialCharTok{\}}\SpecialStringTok{"}\NormalTok{)}
    \BuiltInTok{print}\NormalTok{(}\SpecialStringTok{f"AUC: }\SpecialCharTok{\{}\NormalTok{mejor\_auc}\SpecialCharTok{:.4f\}}\SpecialStringTok{"}\NormalTok{)}
    
    \CommentTok{\# Verificar overfitting}
\NormalTok{    gap }\OperatorTok{=}\NormalTok{ df\_comparacion.loc[df\_comparacion[}\StringTok{\textquotesingle{}Modelo\textquotesingle{}}\NormalTok{] }\OperatorTok{==}\NormalTok{ mejor\_modelo\_nombre, }\StringTok{\textquotesingle{}Gap\textquotesingle{}}\NormalTok{].values[}\DecValTok{0}\NormalTok{]}
    
    \ControlFlowTok{if}\NormalTok{ gap }\OperatorTok{\textgreater{}} \FloatTok{0.15}\NormalTok{:}
        \BuiltInTok{print}\NormalTok{(}\StringTok{"}\CharTok{\textbackslash{}n}\StringTok{⚠️  OVERFITTING DETECTADO"}\NormalTok{)}
        \BuiltInTok{print}\NormalTok{(}\StringTok{"Recomendaciones:"}\NormalTok{)}
        \BuiltInTok{print}\NormalTok{(}\StringTok{"  {-} Recolectar más datos"}\NormalTok{)}
        \BuiltInTok{print}\NormalTok{(}\StringTok{"  {-} Aumentar regularización"}\NormalTok{)}
        \BuiltInTok{print}\NormalTok{(}\StringTok{"  {-} Simplificar modelo"}\NormalTok{)}
        \BuiltInTok{print}\NormalTok{(}\StringTok{"  {-} Remover variables poco importantes"}\NormalTok{)}
    \ControlFlowTok{elif}\NormalTok{ gap }\OperatorTok{\textgreater{}} \FloatTok{0.10}\NormalTok{:}
        \BuiltInTok{print}\NormalTok{(}\StringTok{"}\CharTok{\textbackslash{}n}\StringTok{⚠️  Ligero overfitting"}\NormalTok{)}
        \BuiltInTok{print}\NormalTok{(}\StringTok{"Recomendaciones:"}\NormalTok{)}
        \BuiltInTok{print}\NormalTok{(}\StringTok{"  {-} Considerar regularización adicional"}\NormalTok{)}
        \BuiltInTok{print}\NormalTok{(}\StringTok{"  {-} Validar en conjunto temporal diferente"}\NormalTok{)}
    \ControlFlowTok{else}\NormalTok{:}
        \BuiltInTok{print}\NormalTok{(}\StringTok{"}\CharTok{\textbackslash{}n}\StringTok{✓ Modelo bien balanceado"}\NormalTok{)}
    
    \ControlFlowTok{return}\NormalTok{ resultados}

\CommentTok{\# Ejemplo de uso}
\BuiltInTok{print}\NormalTok{(}\StringTok{"}\CharTok{\textbackslash{}n\textbackslash{}n}\StringTok{EJEMPLO DE PIPELINE COMPLETO"}\NormalTok{)}
\BuiltInTok{print}\NormalTok{(}\StringTok{"="} \OperatorTok{*} \DecValTok{70}\NormalTok{)}

\NormalTok{modelos\_ejemplo }\OperatorTok{=}\NormalTok{ \{}
    \StringTok{\textquotesingle{}Logistic Regression\textquotesingle{}}\NormalTok{: LogisticRegression(random\_state}\OperatorTok{=}\DecValTok{42}\NormalTok{, max\_iter}\OperatorTok{=}\DecValTok{1000}\NormalTok{),}
    \StringTok{\textquotesingle{}Random Forest\textquotesingle{}}\NormalTok{: RandomForestClassifier(random\_state}\OperatorTok{=}\DecValTok{42}\NormalTok{),}
    \StringTok{\textquotesingle{}Decision Tree\textquotesingle{}}\NormalTok{: DecisionTreeClassifier(random\_state}\OperatorTok{=}\DecValTok{42}\NormalTok{)}
\NormalTok{\}}

\NormalTok{param\_grids\_ejemplo }\OperatorTok{=}\NormalTok{ \{}
    \StringTok{\textquotesingle{}Random Forest\textquotesingle{}}\NormalTok{: \{}
        \StringTok{\textquotesingle{}n\_estimators\textquotesingle{}}\NormalTok{: [}\DecValTok{100}\NormalTok{, }\DecValTok{200}\NormalTok{],}
        \StringTok{\textquotesingle{}max\_depth\textquotesingle{}}\NormalTok{: [}\DecValTok{5}\NormalTok{, }\DecValTok{10}\NormalTok{, }\VariableTok{None}\NormalTok{],}
        \StringTok{\textquotesingle{}min\_samples\_split\textquotesingle{}}\NormalTok{: [}\DecValTok{2}\NormalTok{, }\DecValTok{5}\NormalTok{]}
\NormalTok{    \},}
    \StringTok{\textquotesingle{}Decision Tree\textquotesingle{}}\NormalTok{: \{}
        \StringTok{\textquotesingle{}max\_depth\textquotesingle{}}\NormalTok{: [}\DecValTok{3}\NormalTok{, }\DecValTok{5}\NormalTok{, }\DecValTok{7}\NormalTok{, }\DecValTok{10}\NormalTok{],}
        \StringTok{\textquotesingle{}min\_samples\_split\textquotesingle{}}\NormalTok{: [}\DecValTok{2}\NormalTok{, }\DecValTok{5}\NormalTok{, }\DecValTok{10}\NormalTok{]}
\NormalTok{    \}}
\NormalTok{\}}

\NormalTok{resultados\_pipeline }\OperatorTok{=}\NormalTok{ pipeline\_validacion\_completo(}
\NormalTok{    X\_credito, y\_credito,}
\NormalTok{    modelos\_ejemplo,}
\NormalTok{    param\_grids\_ejemplo,}
\NormalTok{    feature\_names}\OperatorTok{=}\NormalTok{X\_credito.columns.tolist()}
\NormalTok{)}
\end{Highlighting}
\end{Shaded}

\begin{verbatim}


EJEMPLO DE PIPELINE COMPLETO
======================================================================
======================================================================
PIPELINE DE VALIDACIÓN COMPLETO
======================================================================

1. COMPARACIÓN INICIAL DE MODELOS
----------------------------------------------------------------------

Evaluando Logistic Regression...
\end{verbatim}

\begin{verbatim}
  AUC: 0.6901 ± 0.1161
  Gap: 0.0524

Evaluando Random Forest...
\end{verbatim}

\begin{verbatim}
  AUC: 0.5376 ± 0.0664
  Gap: 0.4624

Evaluando Decision Tree...
  AUC: 0.5025 ± 0.0438
  Gap: 0.4975


2. OPTIMIZACIÓN DE HIPERPARÁMETROS
----------------------------------------------------------------------

Optimizando Random Forest...
  Mejor AUC: 0.6043
  Mejores parámetros: {'max_depth': 5, 'min_samples_split': 2, 'n_estimators': 100}

Optimizando Decision Tree...
  Mejor AUC: 0.5301
  Mejores parámetros: {'max_depth': 3, 'min_samples_split': 2}


3. ANÁLISIS DE IMPORTANCIA DE VARIABLES
----------------------------------------------------------------------

Random Forest:
  antiguedad_laboral             0.2271
  ingreso_mensual                0.2201
  score_crediticio               0.1674
  edad                           0.1661
  ratio_deuda_ingreso            0.1243
  num_creditos_anteriores        0.0530
  tiene_ahorros                  0.0284
  tiene_hipoteca                 0.0136

Decision Tree:
  antiguedad_laboral             0.3298
  ingreso_mensual                0.3279
  edad                           0.2752
  score_crediticio               0.0671
  num_creditos_anteriores        0.0000
  tiene_hipoteca                 0.0000
  tiene_ahorros                  0.0000
  ratio_deuda_ingreso            0.0000


4. DIAGNÓSTICO Y RECOMENDACIONES
----------------------------------------------------------------------

Mejor modelo: Logistic Regression
AUC: 0.6901

✓ Modelo bien balanceado
\end{verbatim}

\section{Conclusión}\label{conclusiuxf3n}

La validación cruzada y optimización de modelos son fundamentales para
desarrollar sistemas de machine learning robustos y confiables. La clave
está en:

\begin{enumerate}
\def\labelenumi{\arabic{enumi}.}
\tightlist
\item
  \textbf{No confiar en una sola métrica o división}
\item
  \textbf{Usar validación cruzada sistemáticamente}
\item
  \textbf{Optimizar hiperparámetros con datos de validación}
\item
  \textbf{Diagnosticar y corregir overfitting}
\item
  \textbf{Analizar contribución de datos y variables}
\item
  \textbf{Documentar decisiones y resultados}
\end{enumerate}

Con estas técnicas, los economistas pueden desarrollar modelos
predictivos que generalicen bien a nuevos datos y proporcionen insights
valiosos para la toma de decisiones.

\subsection{Recursos adicionales}\label{recursos-adicionales}

Para profundizar en validación y optimización:

\begin{itemize}
\tightlist
\item
  Hastie, Tibshirani \& Friedman - ``Elements of Statistical Learning''
\item
  James et al.~- ``Introduction to Statistical Learning''
\item
  Scikit-learn Model Selection:
  scikit-learn.org/stable/model\_selection.html
\item
  Cross-validation strategies:
  scikit-learn.org/stable/modules/cross\_validation.html
\item
  Hyperparameter tuning guide:
  scikit-learn.org/stable/modules/grid\_search.html
\end{itemize}

\section{Publicaciones Similares}\label{publicaciones-similares}

Si te interesó este artículo, te recomendamos que explores otros blogs y
recursos relacionados que pueden ampliar tus conocimientos. Aquí te dejo
algunas sugerencias:

\begin{enumerate}
\def\labelenumi{\arabic{enumi}.}
\tightlist
\item
  \href{https://numerus-scriptum.netlify.app/python/2020-06-19-instalacion-de-anaconda/index.pdf}{\faIcon{file-pdf}}
  \href{https://numerus-scriptum.netlify.app/python/2020-06-19-instalacion-de-anaconda}{Instalacion
  De Anaconda}
\item
  \href{https://numerus-scriptum.netlify.app/python/2020-06-20-configurar-entorno-virtual-python-anaconda/index.pdf}{\faIcon{file-pdf}}
  \href{https://numerus-scriptum.netlify.app/python/2020-06-20-configurar-entorno-virtual-python-anaconda}{Configurar
  Entorno Virtual Python Anaconda}
\item
  \href{https://numerus-scriptum.netlify.app/python/2021-04-17-01-introducion-a-la-programacion-con-python/index.pdf}{\faIcon{file-pdf}}
  \href{https://numerus-scriptum.netlify.app/python/2021-04-17-01-introducion-a-la-programacion-con-python}{01
  Introducion A La Programacion Con Python}
\item
  \href{https://numerus-scriptum.netlify.app/python/2021-05-31-02-variables-expresiones-y-statements-con-python/index.pdf}{\faIcon{file-pdf}}
  \href{https://numerus-scriptum.netlify.app/python/2021-05-31-02-variables-expresiones-y-statements-con-python}{02
  Variables Expresiones Y Statements Con Python}
\item
  \href{https://numerus-scriptum.netlify.app/python/2021-06-07-03-objetos-de-python/index.pdf}{\faIcon{file-pdf}}
  \href{https://numerus-scriptum.netlify.app/python/2021-06-07-03-objetos-de-python}{03
  Objetos De Python}
\item
  \href{https://numerus-scriptum.netlify.app/python/2021-06-14-04-ejecucion-condicional-con-python/index.pdf}{\faIcon{file-pdf}}
  \href{https://numerus-scriptum.netlify.app/python/2021-06-14-04-ejecucion-condicional-con-python}{04
  Ejecucion Condicional Con Python}
\item
  \href{https://numerus-scriptum.netlify.app/python/2021-06-21-05-iteraciones-con-python/index.pdf}{\faIcon{file-pdf}}
  \href{https://numerus-scriptum.netlify.app/python/2021-06-21-05-iteraciones-con-python}{05
  Iteraciones Con Python}
\item
  \href{https://numerus-scriptum.netlify.app/python/2021-08-16-06-funciones-con-python/index.pdf}{\faIcon{file-pdf}}
  \href{https://numerus-scriptum.netlify.app/python/2021-08-16-06-funciones-con-python}{06
  Funciones Con Python}
\item
  \href{https://numerus-scriptum.netlify.app/python/2021-08-23-07-dataframes-con-python/index.pdf}{\faIcon{file-pdf}}
  \href{https://numerus-scriptum.netlify.app/python/2021-08-23-07-dataframes-con-python}{07
  Dataframes Con Python}
\item
  \href{https://numerus-scriptum.netlify.app/python/2021-11-29-08-prediccion-y-metrica-de-performance-con-python/index.pdf}{\faIcon{file-pdf}}
  \href{https://numerus-scriptum.netlify.app/python/2021-11-29-08-prediccion-y-metrica-de-performance-con-python}{08
  Prediccion Y Metrica De Performance Con Python}
\item
  \href{https://numerus-scriptum.netlify.app/python/2021-12-06-09-metodos-de-machine-learning-para-clasificacion-con-python/index.pdf}{\faIcon{file-pdf}}
  \href{https://numerus-scriptum.netlify.app/python/2021-12-06-09-metodos-de-machine-learning-para-clasificacion-con-python}{09
  Metodos De Machine Learning Para Clasificacion Con Python}
\item
  \href{https://numerus-scriptum.netlify.app/python/2021-12-13-10-metodos-de-machine-learning-para-regresion-con-python/index.pdf}{\faIcon{file-pdf}}
  \href{https://numerus-scriptum.netlify.app/python/2021-12-13-10-metodos-de-machine-learning-para-regresion-con-python}{10
  Metodos De Machine Learning Para Regresion Con Python}
\item
  \href{https://numerus-scriptum.netlify.app/python/2022-10-31-11-validacion-cruzada-y-composicion-del-modelo-con-python/index.pdf}{\faIcon{file-pdf}}
  \href{https://numerus-scriptum.netlify.app/python/2022-10-31-11-validacion-cruzada-y-composicion-del-modelo-con-python}{11
  Validacion Cruzada Y Composicion Del Modelo Con Python}
\item
  \href{https://numerus-scriptum.netlify.app/python/2025-05-10-visualizacion-de-datos-con-python/index.pdf}{\faIcon{file-pdf}}
  \href{https://numerus-scriptum.netlify.app/python/2025-05-10-visualizacion-de-datos-con-python}{Visualizacion
  De Datos Con Python}
\end{enumerate}

Esperamos que encuentres estas publicaciones igualmente interesantes y
útiles. ¡Disfruta de la lectura!






\end{document}
